\documentclass[12pt]{report}
 
\usepackage{algorithm}
\usepackage{algpseudocode}
\usepackage{amsmath}
\usepackage{amssymb}
\usepackage{amsthm}
\usepackage{physics}
\usepackage[pdfencoding=auto, psdextra]{hyperref}
\usepackage{bookmark}
\usepackage{enumitem}
\usepackage{fullpage}
\usepackage{mathtools}
\usepackage{graphicx}
\usepackage{color}
\usepackage{tikz}
\hypersetup{
    colorlinks=true,
    linktoc=all,
    linkcolor=blue,
}

\usepackage{mathtools}

\usetikzlibrary{cd}

% Packages for algorithm
\usepackage{algorithm}
\usepackage{algorithmicx, algpseudocode}
\usepackage{float}

% Definitions for algorithm

\algblockdefx{MRepeat}{EndRepeat}{\textbf{repeat}}{}
\algnotext{EndRepeat}

% Packages for Korean
\usepackage{kotex} 

\graphicspath{{./images/}}

\newcommand{\piecewise}[1]{
   \foreach \f/\a/\b/\open/\closed in {#1}{%
      \draw[domain=\a:\b, smooth, variable=\x] plot ({\x}, \f);
      \foreach \x[evaluate={\y=\f;}] in \open{%
         \draw[fill=white] (\x,\y) circle (.6mm);
      }
      \foreach \x[evaluate={\y=\f;}] in \closed{%
         \fill (\x,\y) circle (.6mm);
      }
   }
}

%% Theorem, Corollary, Lemma
\theoremstyle{plain}
\newtheorem{theorem}{Theorem}[section]
\newtheorem{corollary}[theorem]{Corollary}
\newtheorem{lemma}[theorem]{Lemma}
\newtheorem{proposition}[theorem]{Proposition}

\theoremstyle{remark}
\newtheorem*{remark}{Remark}

\theoremstyle{definition}
\newtheorem{definition}[theorem]{Definition}
\newtheorem{example}[theorem]{Example}

%% User-defined commands
\newcommand{\N}{\mathbb{N}}
\newcommand{\Z}{\mathbb{Z}}
\newcommand{\Q}{\mathbb{Q}}
\newcommand{\R}{\mathbb{R}}
\newcommand{\C}{\mathbb{C}}

%\newcommand{\DD}{\mathbb{D}}
%\newcommand{\HH}{\mathbb{H}}

\newcommand{\Expec}[2]{\underset{#2}{\mathbb{E}}\left[ #1 \right]}
\newcommand{\ExpecTwo}[3]{\underset{#3}{\underset{#2}{\mathbb{E}}}\left[ #1 \right]}
\newcommand{\ExpecThree}[4]{\underset{#4}{\underset{#3}{\underset{#2}{\mathbb{E}}}}\left[ #1 \right]}
\newcommand{\Prob}[2]{\underset{#2}{\mathbb{P}}\left[ #1 \right]}
\newcommand{\Var}[2]{\underset{#2}{\mathrm{Var}}\left[ #1 \right]}

\newcommand{\diam}[1]{\mathrm{diam}\left( #1 \right)}
\newcommand{\hcap}[1]{\mathrm{hcap}\left( #1 \right)}
\newcommand{\Imm}[1]{\mathrm{Im}\left[ #1 \right]}

\newcommand{\Caraconv}{\stackrel{\mathrm{Cara}}{\longrightarrow}}

\newcommand{\RightSqP}{]}

 
\begin{document}
 
\title{Algebra 1}
\author{SNU Graduate course lecture notes \\ Lecture by Prof. Dohyeong Kim, Notes by Juhyeok Choi}
\date{2025 Spring}

 
\maketitle
\tableofcontents

\chapter{Groups}

\section{Monoids}
\begin{definition}
    A \emph{monoid} is a set $S$ equipped with a map
    \begin{align*}
        S\times S &\rightarrow S \\
        (x,y) &\mapsto xy
    \end{align*}
    satisfying the following conditions:
    \begin{enumerate}
        \item There exists an \emph{identity element} $e\in S$ such that $ex=x=xe$ for all $x\in S$.
        \item For all $x,y,z\in S$, $(xy)z = x(yz)$.
    \end{enumerate}
    The identity element is often written as $1$.
\end{definition}

\begin{definition}
    Let $S$ be a monoid. If
    \begin{align*}
        xy = yx
    \end{align*}
    for all $x,y\in S$, then we say $S$ is \emph{commutative} or \emph{abelian}. In that case, we often adopt the additive notation
    \begin{align*}
        S\times S &\rightarrow S \\
        (x,y) &\mapsto x+y
    \end{align*}
    instead of $(x,y)\mapsto xy$ and write the identity element as $0$.
\end{definition}

If
\begin{align*}
    m:S\times S\rightarrow S
\end{align*}
defines a monoid structure on $S$, we obtain a sequence of maps
\begin{align*}
    m_k:S^k\rightarrow S,\quad k=0,1,2,3,\dots
\end{align*}
where
\begin{align*}
    m_k(x_1,\dots,x_k) = m(m_{k-1}(x_1,\dots,x_{k-1}),x_k).
\end{align*}
Here $m_0(S^0)$ is the identity. Then $m_k$ satisfying the conditions below gives an alternative definition of a monoid.
\begin{enumerate}
    \item $m_1(x) = x$ for all $x\in S$.
    \item $m_k(x_1,\dots,x_k) = m(m_{k-1}(x_1,\dots,x_{k-1}),x_k) = m(x_1,m_{k-1}(x_2,\dots,x_{k}))$.
\end{enumerate}

\begin{definition}
    For a set $X$, let
    \begin{align*}
        \mathrm{str}(X)
    \end{align*}
    be the set of strings with letters in $X$. It is a monoid with respect to concatenation. The identity element is the empty string. There exists an identity map $i:X\rightarrow \mathrm{str}(X)$, mapping $x\in X$ to $\text{`$x$'}$. We can construct a \emph{free monoid} as below:
    \begin{figure}[ht]
    \centering
    \begin{tikzcd}
        X \arrow[dr, "i"'] \arrow[rr, "f"] & & M \\
        & A \arrow[ur, dashed, "\phi"'] &
    \end{tikzcd}
    \end{figure}
    \\
    where $\phi$ is a unique monoid homomorphism (meaning $\phi(xy)=\phi(x)\phi(y)$ for all $x$, $y$) between $\mathrm{str}(X)$ and $M$. Here, we say that $f$ \emph{uniquely factors through} $i$.
\end{definition}
\begin{remark}
    This is an example of a universal property.
\end{remark}


\section{Groups}
\begin{definition}
    A monoid $(S,\cdot)$ is called a \emph{group} if every element $s\in S$ has an \emph{inverse} $s^{-1}$ satisfying
    \begin{align*}
        ss^{-1} = s^{-1}s=1.
    \end{align*}
\end{definition}
\begin{remark}
    Let $(G,\cdot)$ be a group. If $xy=1$, then
    \begin{align*}
        y = (x^{-1}x)y=x^{-1}(xy)=x^{-1},
    \end{align*}
    so $y=x^{-1}$ and $yx=1$.
\end{remark}

\begin{example}
    Let $G$ be a group and $X$ be a nonempty set. Denote by
    \begin{align*}
        C(X,G):=\{f:X\rightarrow G\}
    \end{align*}
    the set of all functions mapping $X$ to $G$. For $f_1,f_2\in C(X,G)$, define $f_1f_2$ by
    \begin{align*}
        (f_1f_2)(x) = f_1(x)f_2(x).
    \end{align*}
    Then $C(X,G)$ is a group. To check, observe that
    \begin{enumerate}
        \item There exists an identity element $e(x) \equiv e$ where $e$ is an identity element of $G$.
        \item For all $f,g,h\in C(X,G)$,
        \begin{align*}
            ((fg)h)(x) = (fg)(x)h(x) = f(x)g(x)h(x) = f(x)((gh)(x))=(f(gh))(x),\quad x\in X
        \end{align*}
        so $(fg)h=f(gh)$.
        \item For all $f\in C(X,G)$, there exists an inverse $g$ where $g(x) = (f(x))^{-1}$.
    \end{enumerate}
\end{example}

\begin{example}[permutations]
    For a set $X$, the set
    \begin{align*}
        \mathrm{Perm}(X)
    \end{align*}
    is a collection of all bijections mapping $X$ to itself. $\mathrm{Perm}(X)$ with composition forms a group. This is called a \emph{permutation group}. For $[n]=\{1,2,\dots,n\}$, we often write
    \begin{align*}
        S_n:=\mathrm{Perm}([n])
    \end{align*}
    and call $S_n$ as the \emph{symmetric group} on $n$ letters.
\end{example}
\begin{remark}
    $S_n$ has $n!$ elements, as the number of permutations of $n$ letters is $n\times (n-1)\times\cdots\times 2\times 1$. (For a permutation $\phi$, observe the number of possibilities for $\phi(i)$ after $\phi(1),\dots,\phi(n-1)$ are chosen.)
\end{remark}
\begin{remark}
    For $n\le 2$, $S_1 = \{i\}$ and $S_2 = \{i, (12)\}$, so $S_n$ is commutative. If $n\ge 3$, observe that $(12),(123)\in S_n$ and
    \begin{align*}
        (12)(123) = (23),\quad (123)(12)=(13).
    \end{align*}
    Hence $S_n$ is not commutative. $S_n$ for $n\ge 3$ is one of the simplest examples of non-commutative groups.
\end{remark}

\begin{example}[general linear group]
    Let $V$ be a vector space over a field $k$. Let
    \begin{align*}
        GL(V)
    \end{align*}
    be the collection of $k$-linear automorphisms on $V$. Then $GL(V)$ with composition forms a group. If $V=k^n$, we often write
    \begin{align*}
        GL(n,k):=GL(k^n).
    \end{align*}
\end{example}

\begin{example}[general linear group over a field $\Z/p\Z$]
    Let $p$ be a prime. We observe $GL(n, \Z/p\Z)$. If $n=1$, $GL(n,\Z/p\Z)$ is equivalent to $\Z/p\Z$, so the group is abelian. For $n\ge 2$, let
    \begin{align*}
        f(x_1,x_2,\dots,x_n) &= (x_1+x_2,x_2,x_3,x_4,\dots,x_n), \\
        g(x_1,x_2,\dots,x_n) &= (x_1,x_1+x_2,x_3,x_4,\dots,x_n).
    \end{align*}
    Then $f,g\in GL(V)$ and
    \begin{align*}
        (f\circ g)(x_1,x_2,0,\cdots,0) &= (2x_1+x_2,x_1+x_2,0,\dots,0), \\
        (g\circ f)(x_1,x_2,0,\cdots,0) &= (x_1+x_2,x_1+2x_2,0,\dots,0).
    \end{align*}
    Therefore $GL(n,k)$ is non-commutative. (This result comes directly from the fact that the product of matrices is not commutative.)
\end{example}

\begin{definition}
    Let $(G,\cdot)$ be a group. A subset $H\subset G$ is called a \emph{subgroup} if
    \begin{enumerate}
        \item $H$ is closed under $\cdot$.
        \item $(H,\cdot)$ is a group.
    \end{enumerate}
    We write $H\le G$.
\end{definition}

\begin{example}
    Let $G$ be a group and let $H\subset G$ satisfies
    \begin{align*}
        x,y\in H\implies xy^{-1}\in H.
    \end{align*}
    First, $1=xx^{-1}\in H$, so the identity element is in $H$. Also, for all $x\in H$, $x^{-1} = 1x^{-1}\in H$. So if $x,y\in H$, $y^{-1}\in H$ and
    \begin{align*}
        xy = x(y^{-1})^{-1}\in H.
    \end{align*}
    Hence $xy$ is closed under $\cdot$. To check that $(H,\cdot)$ is a group, it suffices to show that the identity element of $G$ is in $H$ and the inverses of every element of $H$ are also in $H$. This is already shown, so $H$ is itself a group. Therefore $H\le G$.
\end{example}

\begin{example}[multiplicative group of a field]
    If $k$ is a field, then
    \begin{align*}
        k^\times := k-\{0\}
    \end{align*}
    is a group with respect to the multiplication of $k$. 
\end{example}

\begin{example}
    Let $n\in\N$ and $k$ be a field. Let
    \begin{align*}
        \mu_n(k):=\{a\in k^\times:a^n=1\}.
    \end{align*}
    As $k^\times$ is commutative, for all $a,b\in \mu_n$, $(ab)^n = a^n b^n =1$ and $ab\in\mu_n(k)$. So $\mu_n(k)$ is closed under multiplication. Also, as $1^n=1$, $1\in \mu_n(k)$, and for all $a\in \mu_n(k)$, $a^{-n}=(a^n)^{-1}=1$, so $a^{-1}\in\mu_n(k)$. Therefore $\mu_n(k)$ is itself a group, and we conclude that $\mu_n(k)\le k^\times$.
\end{example}

\begin{definition}
    If $G_1$, $G_2$ are groups, then the \emph{direct product}
    \begin{align*}
        G_1\times G_2
    \end{align*}
    is again a group, where the operation is given by
    \begin{align*}
        (x,y),(z,w)=(xz,yw).
    \end{align*}
\end{definition}

\begin{definition}
    A group $G$ is \emph{cyclic} if there exists $g\in G$ such that
    \begin{align*}
        G=\{g^n:n\in\Z\}.
    \end{align*}
    Such $g$ is called a \emph{generator} of $G$.
\end{definition}

\begin{example}
    $\mu_n(\C)$ for $n\in\N$ is an example of a cyclic group. Observe that
    \begin{align*}
        \mu_n(\C) = \{e^{i\frac{2m}{n\pi}}:m=0,1,\dots,n-1\}.
    \end{align*}
    Therefore $e^{2i/n\pi}$ is a generator of $\mu_n(\C)$.
\end{example}


\section{Homomorphisms}
\begin{definition}
    Let $G$, $G'$ be monoids with identity elements $e\in G$ and $e'\in G'$. If a map
    \begin{align*}
        f:G\rightarrow G'
    \end{align*}
    satisfies
    \begin{align*}
        f(e)=e'
    \end{align*}
    and
    \begin{align*}
        f(xy)=f(x)f(y),\quad x,y\in G,
    \end{align*}
    then we say $f$ is a \emph{homomorphism} between $G$ and $G'$.
\end{definition}
\begin{remark}
    For alternative definition, if a map $f:G\rightarrow G'$ preserves all finite products, then we say $f$ is a homomorphism.
\end{remark}

\begin{definition}
    A homomorphism $\alpha:G\rightarrow G'$ is called an \emph{isomorphism} if there is a homomorphism $\beta:G'\rightarrow G$ such that
    \begin{align*}
        \beta\circ \alpha = 1_G,\quad \alpha\circ\beta = 1_{G'}
    \end{align*}
    where $1_A$ denotes the identity function on a set $A$.
\end{definition}

\begin{example}
    Let $G=\Z/5\Z$ and $G'=\mu_5(\C)$. Let $f:G\rightarrow G'$ such that
    \begin{align*}
        f(a)=\exp \left( \frac{2\pi i}{5}a \right).
    \end{align*}
    Then there exists an inverse map $f^{-1}$. As $f'\circ f=1_G$ and $f\circ f^{-1} = 1_{G'}$, it suffices to show that $f$ and $f^{-1}$ are homomorphisms. Observe that
    \begin{align*}
        f(a+b) = \exp \left(\frac{2\pi i}{5}(a+b)\right) = \exp \left( \frac{2\pi i}{5}a \right) \exp \left( \frac{2\pi i}{5}b \right) = f(a)f(b)
    \end{align*}
    and
    \begin{align*}
        f^{-1} \left( \exp \left( \frac{2\pi i}{5}a \right)\exp \left( \frac{2\pi i}{5}b \right) \right) &= f^{-1} \left( \exp \left( \frac{2\pi i}{5}(a+b) \right) \right) \\
        &= a+b \\
        &= f^{-1}\left( \exp \left( \frac{2\pi i}{5}a \right) \right) + f^{-1} \left( \exp \left( \frac{2\pi i}{5}b \right) \right).
    \end{align*}
    Therefore $f$ is an isomorphism between $G$ and $G'$.
\end{example}

\begin{proposition}
    A bijective group homomorphism is an isomorphism.
\end{proposition}
\begin{proof}
    Let $\phi:G\rightarrow G'$ be a bijective group homomorphism. Then there exists an inverse function $\phi^{-1}$. It suffices to show that $\phi^{-1}$ is also a homomorphism. Let $x,y\in G'$. Then
    \begin{align*}
        \phi^{-1}(x)\phi^{-1}(y) = \phi^{-1}(\phi(\phi^{-1}(x)\phi^{-1}(y))) = \phi^{-1}(\phi(\phi^{-1}(x))\phi(\phi^{-1}(y)))=\phi^{-1}(xy).
    \end{align*}
    Therefore $\phi^{-1}$ is also a homomorphism.
\end{proof}

\begin{definition}
    The \emph{kernel} of a function $\alpha$ on a group $G$, denoted as $\mathrm{ker}(\alpha)$, is defined as
    \begin{align*}
        \mathrm{ker}(\alpha) := \{g\in G:\alpha(g) = 1\}.
    \end{align*}
\end{definition}

\begin{proposition}
    A group homomorphism $\alpha$ is injective if and only if $\mathrm{ker}(\alpha)=\{1\}$.
\end{proposition}
\begin{proof}
    If $\alpha$ is injective, then by definition $\mathrm{ker}(\alpha)$ has at most one element. As $\alpha(1) = 1$, $\mathrm{ker}(\alpha)=\{1\}$.

    Conversely, suppose $\mathrm{ker}(\alpha) = 1$. If $\alpha(x)=\alpha(y)$ for $x,y\in G$, then
    \begin{align*}
        \alpha(xy^{-1}) = \alpha(x)(\alpha(y))^{-1}=1.
    \end{align*}
    As $\mathrm{ker}(\alpha)=1$, we obtain that $xy^{-1} = 1$ and $x=y$. Therefore $\alpha$ is injective.
\end{proof}

\begin{proposition}
    For all homomorphism $\alpha:G\rightarrow G'$, $\mathrm{ker}(\alpha)\le G$.
\end{proposition}
\begin{proof}
    First, if $x,y\in\ker(\alpha)$, then as $\alpha(xy)=\alpha(x)\alpha(y) = 1$, $xy\in\ker(\alpha)$. Hence $\ker(\alpha)$ is closed under operation. Also, as $\alpha(1)=1$, $1\in\ker(\alpha)$, and for all $x\in\ker(\alpha)$,
    \begin{align*}
        \alpha(x)\alpha(x^{-1}) = \alpha(xx^{-1}) = \alpha(1) = 1
    \end{align*}
    so $\alpha(x^{-1}) = (\alpha(x))^{-1}=1$ and $x^{-1}\in\ker(\alpha)$. Therefore $\ker(\alpha)=1$.
\end{proof}

\begin{definition}
    A subgroup $H\le G$ is called \emph{normal} if
    \begin{align*}
        g\in G,h\in H\implies ghg^{-1}\in H.
    \end{align*}
\end{definition}

\begin{proposition}
    A kernel of a homomorphism is a normal subgroup.
\end{proposition}
\begin{proof}
    Let $G, G'$ be groups. Let $\alpha$ be a homomorphism between $G$ and $G'$. Then for all $g\in G$ and $h\in\ker(\alpha)$,
    \begin{align*}
        \alpha(ghg^{-1}) = \alpha(g)\alpha(h)\alpha(g^{-1}) = \alpha(g)\alpha(g^{-1}) = \alpha(1) = 1.
    \end{align*}
    Therefore $\ker(\alpha)$ is a normal subgroup.
\end{proof}

\begin{example}
    Let $k$ be a field and let $G\le \mathrm{GL}_2(k)$ be the group of matrices whose bottom row is $(0,1)$. Let $G_1\subset G$ be the unipotent matrices of $G$ and $G_2\subset G$ be the diagonal matrices of $G$. It is clear that $G_1$ and $G_2$ are subgroups of $G$. Note that
    \begin{align*}
        \phi:G&\rightarrow \R^\times \\
        A &\mapsto \det(A)
    \end{align*}
    is a homomorphism $G\rightarrow (\R\setminus\{0\},\cdot)$. As $G_1=\ker \phi$, $G_1$ is a normal subgroup of $G$. However, as
    \begin{align*}
        \begin{pmatrix}
            1 & 1 \\ 0 & 1
        \end{pmatrix}^{-1} \begin{pmatrix}
            2 & 0 \\ 0 & 1
        \end{pmatrix} \begin{pmatrix}
            1 & 1 \\ 0 & 1
        \end{pmatrix} = \begin{pmatrix}
            1 & -1 \\ 0 & 1
        \end{pmatrix} \begin{pmatrix}
            2 & 0 \\ 0 & 1
        \end{pmatrix} \begin{pmatrix}
            1 & 1 \\ 0 & 1
        \end{pmatrix} = \begin{pmatrix}
            2 & 1 \\ 0 & 1
        \end{pmatrix}\notin G_2,
    \end{align*}
    $G_2$ is not a normal subgroup of $G$.
\end{example}


\section{Cosets, normality, and conjugation}
\begin{definition}
    Let $G$ be a group and let $H\le G$. A subset $A\subset G$ is called a \emph{left coset} (of $H$) if
    \begin{align*}
        A=aH=\{ah:h\in H\}
    \end{align*}
    for some $a\in G$. Such $a$ is called a \emph{coset representative}.
\end{definition}
\begin{remark}
    Right cosets can be defined in a similar way.
\end{remark}

\begin{proposition}\label{Proposition 1.4.2.}
    Group can be partitioned into its left cosets.
\end{proposition}
\begin{proof}
    First, for all $g\in G$, there exists a left coset $gH$ containing $g$. Now it suffices to show that for $g,g'\in G$, $gH$ and $gH'$ are the same or are disjoint. If $gH\cap gH'\neq \emptyset$, then there exists $a\in gH\cap gH'$. Then there exists $h,h'\in H$ such that
    \begin{align*}
        a = gh = g'h'.
    \end{align*}
    Then $g = g'h'h^{-1}$, so for all $gh''\in gH$,
    \begin{align*}
        gh'' = g'h'h^{-1}h''\in g'H.
    \end{align*}
    Hence $gH\subset gH'$. Similarly, we can obtain $gH'\subset gH$. Therefore $gH=gH'$ and we obtain the result.
\end{proof}

\begin{example}
    Let $H\subset S_4$ be the cyclic subgroup generated by $(1234)$. Then
    \begin{align*}
        (12)H &= \{(12),(12)(1234),(12)(13)(24),(12)(1432)\} \\
        &= \{(12),(234),(1324),(143)\}
    \end{align*}
    and
    \begin{align*}
        (34)H &= \{(34),(34)(1234),(34)(13)(24),(34)(1432)\} \\
        &= \{(34),(124),(1423),(132)\}.
    \end{align*}
    Therefore $(12)H\neq (34)H$. As
    \begin{align*}
        (13)H &= \{(13),(13)(1234),(13)(13)(24),(13)(1432)\} \\
        &= \{(13),(12)(34),(24),(14)(23)\}
    \end{align*}
    and
    \begin{align*}
        (24)H &= \{(24),(24)(1234),(24)(13)(24),(24)(1432)\} \\
        &= \{(24),(14)(23),(13),(12)(34)\},
    \end{align*}
    we obtain $(13)H=(24)H$.
\end{example}

\begin{example}
    Let $H\subset S_6$ be the cyclic subgroup generated by $(123456)$. Then
    \begin{alignat*}{2}
        (123)H &= \{ &&(123),(123)(123456),(123)(135)(246), \\
        & && (123)(14)(25)(36),(123)(153)(264),(123)(165432)\} \\
        &= \{ && (123), (134562), (24635), (142536), (15)(2643), (1654)\}
    \end{alignat*}
    and as $(456)\not\in (123)H$, $(123)H\neq (456)H$.
\end{example}

\begin{proposition}
    Let $G$ be a group and let $H\le G$. $H$ is normal if and only if
    \begin{align*}
        g\in G\implies gH=Hg.
    \end{align*}
\end{proposition}
\begin{proof}
    First, if $H$ is normal, then for all $g\in G$ and $h\in H$, $ghg^{-1}\in H$. Hence, for some $h'\in H$, $ghg^{-1}=h'$ and $gh=h'g$. Here we obtain $gh\in Hg$ and $h'g\in gH$. This holds for all $h,h'\in H$, so $gH\subset Hg$ and $Hg\subset gH$. Therefore $gH=Hg$. \\

    Conversely, if $gH=Hg$ for all $g\in G$, then for all $h\in H$ there exists $h'\in H$ such that $gh=h'g$ and $ghg^{-1}= h'$. This implies $ghg^{-1}\in H$, so $H$ is normal.
\end{proof}

\begin{definition}
    Let $G$ be a group and let $x,y\in G$. We say $x$ and $y$ are \emph{conjugate} if there exists $g\in G$ such that
    \begin{align*}
        gxg^{-1}=y.
    \end{align*}
\end{definition}

\begin{proposition}
    Let $G$ be a group and let $H\le G$. Then for all $g\in G$,
    \begin{align*}
        gHg^{-1} = \{ghg^{-1}:h\in H\}
    \end{align*}
    is a subgroup of $G$.
\end{proposition}
\begin{proof}
    For $h_1,h_2\in H$,
    \begin{align*}
        (gh_1g^{-1})(gh_2g^{-1}) = g(h_1h_2)g^{-1} \in gHg^{-1},
    \end{align*}
    so $gHg$ is closed under group operation. As $1\in H$,
    \begin{align*}
        1=g1g^{-1}\in gHg^{-1}
    \end{align*}
    so $gHg^{-1}$ contains the identity element. For all $ghg^{-1}\in gHg^{-1}$,
    \begin{align*}
        (ghg^{-1})^{-1}=gh^{-1}g^{-1}\in gHg^{-1}.
    \end{align*}
    Therefore $gHg^{-1}\le G$.
\end{proof}

\begin{definition}
    Let $G$ be a group and let $H_1,H_2\le G$. We say $H_1,H_2$ are \emph{conjugate} if
    \begin{align*}
        gH_1 g^{-1}=H_2
    \end{align*}
    for some $g\in G$.
\end{definition}

\begin{example}
    Let $H\subset S_4$ be the subgroup generated by $(123)$. We observe the subgroups of $S_4$ that are conjugate to $H$. Rather than computing $gHg^{-1}$ for all $g\in G$, we use the fact that $gHg^{-1}$ and $H$ have the same cardinality. (This comes directly from the cancellation law: if $gh_1g^{-1} = gh_2 g^{-1}$, then $h_1=h_2$.) There are four subgroups of $S_4$ with size $3$:
    \begin{align*}
        \langle(123)\rangle, \langle(124)\rangle, \langle(134)\rangle, \langle(234)\rangle.
    \end{align*}
    Note that
    \begin{align*}
        e\langle(123)\rangle e^{-1} &= \langle(123)\rangle, \\
        (34)\langle(123)\rangle(34)^{-1} &= \langle(124)\rangle, \\
        (234)\langle(123)\rangle(234)^{-1} &= \langle(134)\rangle, \\
        (1234)\langle(123)\rangle(1234)^{-1} &= \langle(234)\rangle.
    \end{align*}
    Therefore there are four subgroups of $S_4$ that are conjugate to $H$.
\end{example}

\begin{proposition}
    Let $G$ be a group and let $H\le G$. The set of all left cosets of $H$ has the same cardinality as the set of all right cosets of $H$.
\end{proposition}
\begin{proof}
    Let $\phi$ be a map mapping left cosets of $H$ to right cosets of $H$ such that
    \begin{align*}
        \phi(xH)=Hx^{-1}.
    \end{align*}
    If $xH=yH$, then $xy^{-1}\in H$ and $Hx^{-1} = Hy^{-1}$. So $\phi$ is well defined. It is clear that $\phi$ is surjective. If $Hx^{-1} = Hy^{-1}$, then $xy^{-1}\in H$ and $xH=yH$, so $\phi$ is also injective. As $\phi$ is bijective, we obtain the result.
\end{proof}

\begin{definition}
    The \emph{index} of $H$ in $G$, denoted by $(G:H)$, is defined to be the cardinality of the (left) cosets of $H$.
\end{definition}

\begin{proposition}
    Let $G$ be a group and let $H\le G$. If $(G:H)=2$, then $H$ is normal.
\end{proposition}
\begin{proof}
    Let $g\in G$ and $h\in H$. As $(G:H)=2$, there are two right cosets of $H$. If $g\in H$, then $ghg^{-1}\in H$. If $g\not\in H$, then we can write two right cosets of $H$ as $H$, $Hg$. As $gh\not\in H$, $gh\in Hg$. Hence there exists some $h'\in H$ such that $gh=h'g$, and $ghg^{-1}=h'\in H$. Therefore $ghg^{-1}\in H$ for all $g\in G$ and $h\in H$, and $H$ is normal.
\end{proof}

\begin{theorem}[Lagrange]
    If $G$ is a finite group and $H\le G$, then $|H|$ divides $|G|$.
\end{theorem}
\begin{proof}
    By Proposition \ref{Proposition 1.4.2.} it suffices to show that every left coset has the same cardinality as $H$. Let $gH$ be a left coset of $H$. Then
    \begin{align*}
        \phi:H &\rightarrow gH \\
        h&\mapsto gh
    \end{align*}
    is a bijection. Therefore $|H|=|gH|$, and we obtain the result.
\end{proof}

\begin{definition}
    Let $G$ be a group. The \emph{order} of $g\in G$ refers to the smallest positive integer $n$ such that $g^n=1$.
\end{definition}

\begin{corollary}
    If $G$ is a finite group and $g\in G$, then the order of $g$ divides $|G|$.
\end{corollary}
\begin{proof}
    This is direct from the fact $\langle g\rangle\le G$ and that $|\langle g\rangle|$ is the order of $g$.
\end{proof}

\begin{definition}
    Let $G$ be a group and $N$ be a normal subgroup of $G$. Let $G/N$ be the set of cosets of $N$ and define
    \begin{align*}
        (xN)(yN)=(xy)N
    \end{align*}.
\end{definition}
\begin{remark}
    We should first check that the operation is well defined. If $xN=x'N$, then $xx'^{-1}\in N$ and
    \begin{align*}
        (x'y)N = N(x'y) = N(xx'^{-1})(x'y) = N(xy) = (xy)N.
    \end{align*}
    Also, if $yN=y'N$, then $y'^{-1}y\in N$ and
    \begin{align*}
        (xy')N=(xy')(y'^{-1}y)N = (xy)N.
    \end{align*}
\end{remark}

\begin{proposition}
    The following holds.
    \begin{enumerate}
        \item $G/N$ is a group.
        \item $g\mapsto gN$ defines a homomorphism $\pi:G\rightarrow G/N$.
        \item $\ker \pi = N$.
    \end{enumerate}
\end{proposition}
\begin{proof}
    \begin{enumerate}
        \item As the operation is associative in $G$, the operation of $G/N$ is also associative. $G/N$ contains the identity element $N$, and for all $xN\in G/N$ there exists an inverse $x^{-1}N$. Therefore $G/N$ is a group.
        \item Note that
        \begin{align*}
            \pi(xy) = (xy)N = (xN)(yN)=\pi(x)\pi(y).
        \end{align*}
        \item $\pi(x) = xN = N$ if and only if $x\in N$. Therefore $\ker \pi = N$.
    \end{enumerate}
\end{proof}
\begin{remark}
    By the third assertion, we obtain that a subgroup of a group is normal if and only if it is a kernel of some homomorphism.
\end{remark}
\begin{remark}
    The quotient homomorphism $\pi:G\rightarrow G/N$ satisfies the following universal property: when $f(N)=1$, then
    \begin{figure}[H]
    \centering
    \begin{tikzcd}
        G \arrow[dr, "\pi"'] \arrow[rr, "f"] & & H \\
        & G/N \arrow[ur, dashed, "\exists!\phi"'] &
    \end{tikzcd}
    \end{figure}
    \noindent
    That is, if a homomorphism $f:G\rightarrow H$ satisfies $f(N)$, then $f$ uniquely factors through $\pi$.
\end{remark}

\begin{proposition}
    The universal property of the quotient homomorphism characterizes the isomorphism type of the quotient group $G/N$. That is, if another homomorphism $\pi':G\rightarrow G'$ satisfies the same universal property, then we have an isomorphism $G'\simeq G/N$.
\end{proposition}
\begin{proof}
    Note that $\pi(G/N) = \pi'(G')=1$. By the universal property, the homomorphisms
    \begin{align*}
        \phi:G/N\rightarrow G',\quad \phi':G'\rightarrow G/N
    \end{align*}
    satisfying $\pi'=\phi\circ\pi$ and $\pi=\phi'\circ\pi'$ are unique. Here, as $\pi'=\phi\circ\phi'\circ \pi'$ and as $\pi'$ is injective, $\phi\circ\phi'=i_{G'}$. Similarly, $\phi'\circ\phi = i_{G/N}$. Therefore $\phi$ and $\phi'$ are isomorphisms and we obtain $G'\simeq G/N$.
\end{proof}


\section{Group action}
\begin{definition}
    Let $G$ be a group and $S$ be a set. An \emph{action} of $G$ on $S$ is a homomorphism
    \begin{align*}
        G\rightarrow\mathrm{Perm}(S).
    \end{align*}
    Equivalently, an action is a map $G\times S\rightarrow S$ satisfying
    \begin{enumerate}
        \item $1s=s$ for all $s\in S$.
        \item $(xy)s=x(ys)$ for all $x,y\in G$ and $s\in S$. 
    \end{enumerate}
\end{definition}
\begin{remark}
    Since elements of $G$ act on the left side of $S$, such an action is often called a \emph{left action}. A right action can be similarly defined as $S\times G\rightarrow S$.
\end{remark}

\begin{example}
    For a group $G$, the map $G\times G\rightarrow G$, $(x,y)\mapsto xy$ defines a left action of $G$ on $G$. This can also be seen as a right action of $G$ on $G$.
\end{example}

\begin{example}
    For a group $G$, the map $G\times G\rightarrow G$, $(x,y)\mapsto xyx^{-1}$ defines a left action of $G$ on $G$. Similarly, the map $G\times G\rightarrow G$, $(x,y)\mapsto yxy^{-1}$ defines a right action of $G$ on $G$. Note that the first mapping cannot be seen as a right action of $G$ on $G$.
\end{example}

\begin{example}
    Let $G$ be a group and let $S_G$ be the set of all subgroups of $G$. If $H\in S_G$ and $g\in G$, then $gHg^{-1}\in S_G$. Therefore, the conjugation defines a map
    \begin{align*}
        G\times S_G\rightarrow S_G.
    \end{align*}
    This is an action of $G$ on $S_G$.
\end{example}

\begin{example}
    Observe an action $\phi:S_3\rightarrow \mathrm{Perm}(S_3)$ given by conjugation. To show that $\phi$ is injective, it suffices to show that $\ker(\phi) = \{e\}$. Let $g\in S_3$. If
    \begin{align*}
        gHg^{-1}=H,\quad \forall H\in S_{S_3},
    \end{align*}
     then
     \begin{align*}
         g\langle (12)\rangle g^{-1} = \langle(12)\rangle,
     \end{align*}
     so $g(12)g^{-1} = (12)$ and we obtain $g\in \{e,(12)\}$. Similarly, as
     \begin{align*}
         g\langle(13)\rangle g^{-1}=\langle (13)\rangle,
     \end{align*}
     $g(13)g^{-1}=(13)$ and we obtain $g\in \{e,(13)\}$. Therefore $g=e$ and we obtain $\ker(\phi)=\{e\}$.
\end{example}

\begin{example}
    Observe an action $\phi:\mathrm{GL}(2,\C)\rightarrow S_{\mathrm{GL}(2,\C)}$ given by conjugation. As for all diagonal matrices $D\in\mathrm{GL}(2,\C)$,
    \begin{align*}
        DMD^{-1} = M,\quad \forall M\in\mathrm{GL}(2,\C),
    \end{align*}
    $D\in \ker(\phi)$. Therefore $\phi$ is not injective.
\end{example}

\begin{example}
    Let $V$ be a vector space. We have a map
    \begin{align*}
        \mathrm{GL}(V) \times V & \rightarrow V \\
        (g,v) &\mapsto gv.
    \end{align*}
    With this map, $\mathrm{GL}(V)$ acts on $V$.
\end{example}

\begin{example}
    Let $G$ be a group and let $H\le G$. Let $G/H$ be the set of left cosets of $H$. We have a map
    \begin{align*}
        G\times G/H & \rightarrow G/H \\
        (g,xH) &\mapsto (gx)H.
    \end{align*}
    This action is well-defined, as if $xH=yH$, $x^{-1}y\in H$ and
    \begin{align*}
        (gx)H=(gx)(x^{-1}y)H = (gy)H.
    \end{align*}
\end{example}

\begin{example}
    Let $G=S_4$.
    \begin{enumerate}
        \item Let
        \begin{align*}
            H=\{e,(12)(34),(13)(24),(14)(23)\}.
        \end{align*}
        Consider the kernel of the action of $G$ on $G/H$. For $g\in G$, $g$ is in the kernel if and only if $(gx)H=xH$ for all $x\in G$. As $gH=H$, the kernel is a subset of $H$. As $H$ is normal, we obtain that the kernel is $H$.
        \item Let
        \begin{align*}
            H=\{e,(1234),(13)(24),(1432)\}.
        \end{align*}
        Consider the kernel of the action of $G$ on $G/H$. For $g\in G$, $g$ is in the kernel if and only if $(gx)H=xH$ for all $x\in G$. As $gH=H$, the kernel is a subset of $H$. Consider
        \begin{align*}
            (12)H = \{(12),(234),(1324),(143)\}.
        \end{align*}
        As
        \begin{align*}
            (1234)(12)&=(134) \not\in (12)H, \\
            (13)(24)(12)&=(1423)\not\in (12)H, \\
            (1432)(12)&=(243)\not\in (12)H,
        \end{align*}
        we obtain that the kernel is $\{e\}$.
    \end{enumerate}
\end{example}
\begin{remark}
    If $G$ is a group and $H$ is a normal subgroup of $G$, then the kernel of the action of $G$ on $G/H$ is $H$. In general cases, the kernel of the action of $G$ on $G/H$ is the largest subgroup of $H$ which is normal in $G$.
\end{remark}


\section{G-sets and their structures}
\begin{definition}
    Let $G$ be a group and $S$ be a $G$-set. Given an element $s\in S$, the map
    \begin{align*}
        G &\rightarrow S \\
        x &\mapsto xs
    \end{align*}
    is a \emph{$G$-map}. Its image $Gs$ is called the \emph{orbit} of $s$. A subset $A\subset S$ is called an \emph{orbit} if $A=Gs$ for some $s\in S$.

    If $S$ itself is an orbit, we say that the action of $G$ on $S$ is \emph{transitive} or that $S$ is a \emph{transitive $G$-set}.
\end{definition}

\begin{example}
    Let $X$ be a nonempty set. $X$ is a transitive $\mathrm{Perm}(X)$-set. To prove, fix $x_0\in X$. Then for all $x\in X$ there exists $\phi\in \mathrm{Perm}(X)$ such that $\phi(x_0) = x$. Therefore $X = (\mathrm{Perm}(X))x_0$.
\end{example}

\begin{example}
    Let $V$ be a vector space. $V$ is not a transitive $\mathrm{GL}(V)$-set. Let $v\in V$. If $v\neq \mathbf{0}$, then for all $M\in\mathrm{GL}(V)$ $MV\neq \mathbf{0}$. If $v=0$, then for all $M\in\mathrm{GL}(V)$, $MV=\mathbf{0}$. In fact, the map $\mathrm{GL(V)}\rightarrow V$ constructs two orbits in $V$: $\mathbf{0}$ and $V\setminus\mathbf{0}$.
\end{example}

\begin{proposition}
    Let $X$ be a $G$-set. The following statements are equivalent.
    \begin{enumerate}
        \item $X$ is transitive.
        \item For any $x,y\in X$, $gx=y$ for some $g\in G$.
    \end{enumerate}
\end{proposition}
\begin{proof}
    First, suppose that $X$ is transitive. Then for some $x_0\in X$, $Gx_0 = X$. Then for any $x,y\in X$, there exists $g_1, g_2\in G$ such that
    \begin{align*}
        g_1 x_0 = x,\quad g_2x_0 = y.
    \end{align*}
    Then
    \begin{align*}
        y = g_2x_0 = g_2g_1^{-1}x,
    \end{align*}
    so the second statement holds.

    Conversely, suppose the second statement holds. Then fixing $x$ we obtain $Gx=X$, so $X$ is transitive.
\end{proof}

\begin{example}
    Let $G$ be a group and let $H\le G$. Then we can view $G/H$ as a left $G$-set. As
    \begin{align*}
        (yx^{-1})(xH)=yH
    \end{align*}
    for all $x,y\in G$, by the proposition above we obtain that $G/H$ is transitive.
\end{example}

\begin{definition}
    If $X$ is a $G$-set and $x\in X$, then the \emph{isotropy group} is
    \begin{align*}
        G_x := \{g\in G:gx=x\}.
    \end{align*}
\end{definition}

\begin{proposition}
    If $X$ is a transitive $G$-set, then for a fixed $x\in X$,
    \begin{alignat*}{2}
        t: & X &&\rightarrow G/G_x \\
        & gx &&\mapsto gG_x
    \end{alignat*}
    is a bijection.
\end{proposition}
\begin{proof}
    First, we check that the map $t$ is well-defined. If $g_1x=g_2x$, then $x = g_1^{-1}g_2x$ and $g_1^{-1}g_2\in G_X$. Hence
    \begin{align*}
        g_1G_x = g_1(g_1^{-1}g_2)G_x = g_2G_x.
    \end{align*}
    For all $g\in G$, there exists $gx\in X$ such that $t(gx)=gG_x$, so the map $t$ is surjective. For all $g_1,g_2\in G$, if $g_1G_x = g_2G_x$, then $g_1^{-1}g_2\in G_x$, so
    \begin{align*}
        g_1x = g_1 (g_1^{-1}g_2 x) = g_2 x.
    \end{align*}
    Therefore $t$ is injective and we obtain the result.
\end{proof}
\begin{remark}
    The map $t$ is not only a bijection but also satisfies
    \begin{align*}
        t(gx)=g(t(x)).
    \end{align*}
    Such a map is called an \emph{equivalence} between $G$-sets.
\end{remark}

\begin{definition}
    For $G$-sets $X$ and $Y$, a map
    \begin{align*}
        r:X\rightarrow Y
    \end{align*}
    is called a \emph{morphism} of $G$-sets if
    \begin{align*}
        r(gx) = g(r(x))
    \end{align*}
    for all $g\in G$ and $x\in X$. A bijective morphism is called an \emph{equivalence}.
\end{definition}

\begin{definition}
    If $X$ is a $G$-set, then the \emph{orbit decomposition} refers to
    \begin{align*}
        X = \bigsqcup_{i\in I} O_i
    \end{align*}
    where $O_i$'s are the orbits in $X$.
\end{definition}
\begin{remark}
    Since each $O_i$ is equivalent to a $G$-set of the form $G/G_i$, $G_i$ being a subgroup, we conclude: for any $G$-set there is an equivalence
    \begin{align*}
        X\simeq \bigsqcup_{i\in I} G/G_i
    \end{align*}
    where $G_i\le G$.
\end{remark}

Conjugation action is one of the most natural group action acting on itself. The followings are examples of orbit decompositions of groups under the conjugation action.

\begin{example}
    Let $G=S_4$. The decomposition of $G$ into orbits under the conjugation action can be written as below:
    \begin{itemize}
        \item $\{e\}$,
        \item $\{(12),(13),(14),(23),(24),(34)\}$,
        \item $\{(12)(34),(13)(24),(14)(23)\}$,
        \item $\{(123),(132),(124),(142),(134),(143),(234),(243)\}$,
        \item $\{(1234),(1243),(1324),(1342),(1423),(1432)\}$.
    \end{itemize}
\end{example}

\begin{example}
    Let $G=\mathrm{GL}(2,\Z/2\Z)$. The decomposition of $G$ into orbits under the conjugation action can be written as below:
    \begin{itemize}
        \item $\left\{\begin{pmatrix}
            1 & 0 \\ 0 & 1
        \end{pmatrix}\right\}$,
        \item $\left\{\begin{pmatrix}
            1 & 0 \\ 1 & 1
        \end{pmatrix}, \begin{pmatrix}
            1 & 1 \\ 0 & 1
        \end{pmatrix}\right\}$,
        \item $\left\{\begin{pmatrix}
            0 & 1 \\ 1 & 0
        \end{pmatrix}\right\}$,
        \item $\left\{\begin{pmatrix}
            0 & 1 \\ 1 & 1
        \end{pmatrix}, \begin{pmatrix}
            1 & 1 \\ 1 & 0
        \end{pmatrix}\right\}$.
    \end{itemize}
\end{example}


\section{Orbit decomposition formula}
Let $G$ be a group and $S$ be a $G$-set. Consider the orbit decomposition
\begin{align*}
    S = \bigsqcup_{i\in I} Gs_i.
\end{align*}
Since
\begin{align*}
    |Gs_i|=(G:G_{s_i}),
\end{align*}
we have
\begin{align*}
    |S|=\sum_{i\in I} (G:G_{s_i}).
\end{align*}
The above equality is called the \emph{orbit decomposition formula}.\\

\begin{example}
    Let $G$ be a finite group, acting on $G$ via conjugation. Let $C$ be the set of representatives of conjugacy classes of $G$. Then by the orbit decomposition formula, we obtain
    \begin{align*}
        (G:1)=\sum_{x\in C}(G:G_x).
    \end{align*}
\end{example}

\begin{example}
    Let $G\subset \mathrm{GL}(2,\Z/3\Z)$ be the subgroup of upper triangular matrices. It is composed of $12$ elements. Now we verify the orbit decomposition formula. $G$ can be decomposed as
    \begin{align*}
        G =& \begin{pmatrix}
            1 & 0 \\ 0 & 1
        \end{pmatrix} G \sqcup \begin{pmatrix}
            1 & 0 \\ 0 & 2
        \end{pmatrix} G \sqcup \begin{pmatrix}
            2 & 0 \\ 0 & 1
        \end{pmatrix} G \sqcup \begin{pmatrix}
            2 & 0 \\ 0 & 2
        \end{pmatrix} G \sqcup \begin{pmatrix}
            1 & 1 \\ 0 & 1
        \end{pmatrix} G \sqcup \begin{pmatrix}
            2 & 1 \\ 0 & 2
        \end{pmatrix} G.
    \end{align*}
    Each representative has order $1$, $3$, $3$, $1$, $2$, $2$, and the orbit decomposition formula holds.
\end{example}


\section{\texorpdfstring{$p$-groups and Sylow $p$-subgroups}{p-groups and Sylow p-subgroups}}
\begin{definition}
    Let $G$ be a finite group and $p$ be a prime. $G$ is a \emph{$p$-group} if $|G|$ is a power of $p$. A subgroup $H\le G$ is called a \emph{$p$-subgroup} if $H$ is a $p$-group.
\end{definition}

\begin{definition}
    Let $G$ be a group. A $p$-subgroup $H\le G$ is called a \emph{$p$-Sylow subgroup} if $(G:H)$ is relatively prime to $p$.
\end{definition}

\begin{example}
    Let $H\le \mathrm{GL}(2,\Z/p\Z)$ be the subgroup of all upper triangular unipotent matrices. As
    \begin{align*}
        |\mathrm{GL}(2,\Z/p\Z)| &= (p^2-1)(p^2-p) = p(p-1)(p^2-1)
    \end{align*}
    and
    \begin{align*}
        |H|=p,
    \end{align*}
    $H$ is a $p$-Sylow subgroup of $\mathrm{GL}(2,\Z/p\Z)$.
\end{example}

\begin{lemma}
    Let $G$ be a finite abelian group. If $p$ is a prime divisor of $|G|$, then $G$ has a subgroup of order $p$.
\end{lemma}
\begin{proof}
    First, if $G$ is an abelian $p$-group, choose any $g\in G$ other than the identity element. As $\langle g\rangle\le G$, $|\langle g\rangle|=p^k$ for some $k$. Then consider $\langle g^{p^{k-1}}\rangle$; this is a subgroup of $G$ of order $p$.

    Let $G$ be an abelian group such that $p||G|$. If $G$ is a $p$-group, we obtain the result. If $G$ is not a $p$-group, choose any $g\in G$ other than the identity element. If $|\langle g\rangle| = |G|$, then $\langle \frac{|G|}{p}g\rangle$ is a subgroup of $G$ of order $p$. Otherwise, note that
    \begin{align*}
        (G:1)=(G:\langle g\rangle)(\langle g\rangle : 1).
    \end{align*}
    As $(G:\langle g\rangle),(\langle g\rangle)\neq 1$, at least one is a number dividing $p$ smaller than $|G|$. Hence at least one of $\langle g\rangle$ and $gG$ is a subgroup of $G$ whose order is smaller than $|G|$ and dividing $p$. We can continue this process; as this should terminate in finite steps, we conclude that there exists a subgroup of $G$ of order $p$.
\end{proof}
\begin{remark}
    This can also be directly obtained from the classification theorem of finitely generated abelian groups(Theorem \ref{Theorem 1.18.8.}). 
\end{remark}

\begin{theorem}
    Let $G$ be a finite group and $p$ be a prime factor of $|G|$. There is a $p$-Sylow subgroup of $G$.
\end{theorem}
\begin{proof}
    We assume the contrary, and suppose that $G$ is a counterexample with minimal order. First, we prove that the center $Z$ of $G$
    \begin{align*}
        Z:=\{z\in G:zg=gz\quad \forall g\in G\}
    \end{align*}
    has order divisible by $p$. Here we apply the orbit decomposition formula. Letting $C$ to be a set of representatives for the conjugacy classes of $G$, we get
    \begin{align*}
        (G:1)=\sum_{x\in C}(G:G_x).
    \end{align*}
    As $G=G_x$ if and only if $x\in Z$,
    \begin{align*}
        (G:1)=(Z:1)+\sum_{x\in C\setminus Z} (G:G_x).
    \end{align*}
    As $G$ does not have a $p$-Sylow subgroup, $(G:G_x)$ should be a multiple of $p$ for $x\in C\setminus Z$. Hence we get $p||Z|$.

    As $Z$ is a finite group with order divisible by $p$, there exists $H\le Z$ of order $p$. As $H$ is a subgroup of a center of $G$, $H$ is normal in $G$. By assumption, the quotient group $G/H$ has a $p$-Sylow subgroup $K_0$. Let
    \begin{align*}
        f:G\rightarrow G/H
    \end{align*}
    be the projection and let $G_0 = f^{-1}(K_0)$. Then as $K_0$ is a $p$-Sylow subgroup of $G/H$ and $|G_0|=p|K_0|$, $G_0$ is a $p$-Sylow subgroup of $G$. (Let $|G|=p^k m$ where $p$ and $m$ are coprime. Then $|G/H|=p^{k-1}m$, $|K_0|=p^{k-1}$ and $|G_0|=p^k$.) This contradicts our assumption, implying that every finite group with its order dividing $p$ has a $p$-Sylow subgroup.
\end{proof}


\section{More on Sylow subgroups}
\begin{definition}
    For a subgroup $M\le G$, the \emph{normalizer} of $M$ in $G$ refers to the set
    \begin{align*}
        N_G(M):=\{g\in G:gMg^{-1}=M\}.
    \end{align*}
\end{definition}

\begin{theorem}
    Let $G$ be a finite group and $P\subset G$ be a Sylow $p$-subgroup. If $H\le G$ is a nontrivial $p$-group, then $H\le gPg^{-1}$ for some $g\in G$.
\end{theorem}
\begin{proof}
    Let $Q$ be a Sylow $p$-subgroup and assume that $H$ normalizes $Q$, i.e.
    \begin{align*}
        hQh^{-1}=Q\quad \forall h\in H.
    \end{align*}
    Let $N$ be a normalizer of $Q$ in $G$. Then by definition, $H\le N$. Observe the natural map
    \begin{align*}
        H\rightarrow N/Q.
    \end{align*}
    As $H$ is a $p$-group, it has an order power of $p$. As $Q$ is a Sylow $p$-subgroup of $G$ and $N\le G$, $N/Q$ has an order coprime with $p$. Hence the map is a trivial map, and $H\le Q$. 

    Now, it is enough to show that $P$ has a conjugate normalized by $H$. Let $S$ be a collection of conjugates of $P$, and let $H$ act on $S$ by conjugation. Think of an orbit decomposition of $S$. As $H$ is a $p$-group, every non-singleton orbit has order divisible by $p$. As $p\nmid |S|$, there exists an orbit with order $1$. As the element $Q$ in that orbit is a fixed point of $H$-action, $H$ normalizes $Q$, and we obtain the result.
\end{proof}

\begin{corollary}
    Two Sylow $p$-subgroups in a finite group $G$ are conjugate.
\end{corollary}
\begin{proof}
    Let $P$, $Q$ be two Sylow $p$-subgroups of $G$. Then as $|P|=|Q|$ and $P\le gQg^{-1}$ for some $g\in G$, $P=gQg^{-1}$, and the two subgroups are conjugate.
\end{proof}


\section{Number of Sylow subgroups}
\begin{theorem}
    Let $G$ be a group of order dividing $p$ and let $n$ be the number of Sylow $p$-subgroups in $G$. Then $n\equiv 1 \mod p$.
\end{theorem}
\begin{proof}
    Let $P$ be a Sylow $p$-subgroup of $G$ and $S$ be the set of conjugates of $P$. We can consider $S$ as a $P$-set. It is clear that $P\in S$ is a fixed point. Now assume $Q\in S$ is a fixed point. Let $N$ be a normalizer of $Q$ in $G$. Then by definition and assumption, $P\subset N$. As $P,Q\le N\le G$ and $P,Q$ are Sylow $p$-subgroups of $G$, $P,Q$ are Sylow $P$-subgroups of $N$. As $Q$ is normal in $N$, every conjugate of $Q$ in $N$ is $Q$ itself, and hence $P=Q$. Therefore $P$ is a unique fixed point in $S$.

    Now we apply the orbit decomposition formula to $S$. Let $C$ be a set of representatives for the conjugacy classes of $S$. Then
    \begin{align*}
        (S:1) = \sum_{H\in C} (S:G_H)
    \end{align*}
    As $(S:G_P) = 1$ and $p|(S:G_H)$ for $H\neq P$, we obtain
    \begin{align*}
        (S:1)\equiv 1\mod p.
    \end{align*}
\end{proof}
\begin{remark}
    Also, note that the number of Sylow $p$-subgroups of a finite group should divide the order of the group. With this fact, we can efficiently reduce the candidates for the number of Sylow $p$-subgroups.
\end{remark}

The main theorems of section 8, 9, and 10 together are called \emph{Sylow theorems}.


\section{\texorpdfstring{$p$-groups as matrix groups}{p-groups as matrix groups}}

In this section, we observe that every $p$-group is isomorphic to a subgroup of a typical matrix group.

\begin{definition}
    Let $\mathrm{U}(n,\Z/p\Z)\le \mathrm{GL}(n,\Z/p\Z)$ be the set of upper triangular unipotent matrices.
\end{definition}

\begin{theorem}
    For $n\ge 2$, $\mathrm{U}(n,\Z/p\Z)$ is a Sylow $p$-subgroup of $\mathrm{GL}(n,\Z/p\Z)$.
\end{theorem}
\begin{proof}
    Note that
    \begin{align*}
        |\mathrm{GL}(n,\Z/p\Z)| &= (p^n-1)(p^n-p)\cdots(p^n-p^{n-1}) \\
        &= p^{1+2+\cdots+(n-1)}(p^n-1)(p^{n-1}-1)\cdots(p-1)
    \end{align*}
    and
    \begin{align*}
        |\mathrm{U}(n,\Z/p\Z)| = p^{(n-1)+(n-2)+\cdots+1}.
    \end{align*}
\end{proof}

\begin{theorem}\label{Theorem 1.11.3.}
    For any finite group $G$ and any field $k$, there exists an injective homomorphism
    \begin{align*}
        G\rightarrow \mathrm{GL}(n,k)
    \end{align*}
    for $n=|G|$.
\end{theorem}
\begin{proof}
    First, there exists an injective embedding $G\rightarrow S_n$, naturally induced by the regular action of $G$ on itself. There also exists an injective embedding $S_n\rightarrow \mathrm{GL}(n,k)$; think of permutation matrices. Therefore, there exists an injective homomorphism from $G$ to $\mathrm{GL}(n,k)$.
\end{proof}

\begin{theorem}
    Any $p$-group is isomorphic to a subgroup of $\mathrm{U}(n,\Z/p\Z)$, where $n$ is the order of a group.
\end{theorem}
\begin{proof}
    Let $G$ be a $p$-group. By Theorem \ref{Theorem 1.11.3.} there exists some $n$ with an injective homomorphism
    \begin{align*}
        G\rightarrow\mathrm{GL}(n,\Z/p\Z).
    \end{align*}
    By Sylow theorem, the image of the homomorphism is conjugate to a subgroup of $\mathrm{U}(n,\Z/p\Z)$. Such a subgroup is isomorphic to $G$.
\end{proof}

\begin{example}
    Observe an embedding
    \begin{align*}
        \Z/2\Z \times \Z/2\Z \rightarrow \mathrm{U}(n,\Z/2\Z).
    \end{align*}
    The theorem above guarantees that such an embedding exists for $n\ge 4$. The embedding also exists for $n=3$;
    \begin{align*}
        (a,b)\mapsto \begin{pmatrix}
            1 & a & b \\ 0 & 1 & 0 \\ 0 & 0 & 1
        \end{pmatrix}.
    \end{align*}
\end{example}

\begin{example}
    Think of a group $\mathrm{U}(3,\Z/p\Z)$. As this is a $p$-group, every element has an order a power of $p$. Note that
    \begin{align*}
        \begin{pmatrix}
            1 & a & b \\ 0 & 1 & c \\ 0 & 0 & 1
        \end{pmatrix} \begin{pmatrix}
            1 & a' & b' \\ 0 & 1 & c' \\ 0 & 0 & 1
        \end{pmatrix} = \begin{pmatrix}
            1 & a+a' & b+b'+ac' \\ 0 & 1 & c+c' \\ 0 & 0 & 1
        \end{pmatrix}.
    \end{align*}
    As
    \begin{align*}
        \begin{pmatrix}
            1 & * & * \\ 0 & 1 & * \\ 0 & 0 & 1
        \end{pmatrix}^{p^2} = \begin{pmatrix}
            1 & 0 & * \\ 0 & 1 & 0 \\ 0 & 0 & 1
        \end{pmatrix}^p = \begin{pmatrix}
            1 & 0 & 0 \\ 0 & 1 & 0 \\ 0 & 0 & 1
        \end{pmatrix},
    \end{align*}
    there does not exist an element with order larger than $p^2$. As for all $A\in \mathrm{U}(3,\Z/p\Z)$, $A$ is unipotent and therefore satisfies
    \begin{align*}
        (A-I)^3 = 0.
    \end{align*}
    Here, from the binomial expression we obtain
    \begin{align*}
        A^k = ((A-I)+I)^k = I + k(A-I) + \frac{k(k-1)}{2}(A-I)^2.
    \end{align*}
    for $k\in\N$. For $p\ge 3$, as
    \begin{align*}
        A^p = I + p(A-I) + \frac{p(p-1)}{2}(A-I)^2 = I,
    \end{align*}
    the largest order of an element in $\mathrm{U}(3,\Z/p\Z)$ is $p$. If $p=2$, note that
    \begin{align*}
        \begin{pmatrix}
            1 & 1 & 0 \\ 0 & 1 & 1 \\ 0 & 0 & 1
        \end{pmatrix}^2 = \begin{pmatrix}
            1 & 0 & 1 \\ 0 & 1 & 0 \\ 0 & 0 & 1
        \end{pmatrix}
    \end{align*}
    and there exists an element of order $p^2 = 4$.
\end{example}


\section{Classification of some finite groups}
\begin{definition}
    Let $G$ be a group. Suppose that $N,H\le G$ are disjoint and $N$ is normal in $G$. If any element $g\in G$ can be written uniquely in the form
    \begin{align*}
        g=nh,\quad n\in N,h\in H,
    \end{align*}
    we say $G$ is an \emph{(inner) semidirect product} of $N$ and $H$. 
\end{definition}

\begin{definition}
    Suppose that $N,H$ are groups and
    \begin{align*}
        \phi:H\rightarrow\mathrm{Aut}(N)
    \end{align*}
    is a group homomorphism. Let $\phi(h)=\phi_h$. The \emph{outer semidirect product} $N\rtimes_\phi H$ to be the set $N\times H$ equipped with
    \begin{align*}
        (n,h)(m,k) = (n(\phi_h(m)),hk),\quad (n,h),(m,k)\in N\rtimes_\phi H.
    \end{align*}
\end{definition}

\begin{theorem}
    The definitions of inner semidirect product and outer semidirect product are equivalent.
\end{theorem}
\begin{proof}
    First we consider the inner semidirect product $G$ of $N$ and $H$. Consider the projections
    \begin{align*}
        \pi_1:G&\rightarrow N,\quad \pi_2:G\rightarrow H
    \end{align*}
    where $\pi_1(nh)=n$ and $\pi_2(nh)=h$ for $n\in N$, $h\in H$. Note that by the definition of the inner semidirect product, the two projections are well-defined. Let
    \begin{align*}
        \phi_h : N &\rightarrow N \\
        n &\mapsto \pi_1(hn).
    \end{align*}
    As $N$ is normal, for $n\in N$ and $h\in H$, $hN=Nh$ and
    \begin{align*}
        hn = \pi_1(hn)h.
    \end{align*}
    Hence for $n_1,n_2\in N$ and $h_1,h_2\in H$,
    \begin{align*}
        (n_1h_1)(n_2h_2) &= n_1 \pi_1(h_1n_2)h_1h_2 \\
        &= n_1\phi_h(n_2)h_1h_2.
    \end{align*}
    Therefore there exists an isomorphism between $G$ and $N\rtimes_\phi H$.
\end{proof}

\begin{example}
    Let $G=\Z/7\Z \rtimes \Z/6\Z$ with the action $\phi_1(x) = tx$ for $t\in(\Z/7\Z)^\times$. First, we observe that
    \begin{align*}
        (n,h)(m,k)=(n+t^hm,h+k),\quad (n,h),(m,k)\in G,
    \end{align*}
    so
    \begin{align*}
        (m,k)^{-1} = (-t^{-k}m,-k),\quad (m,k)\in G.
    \end{align*}
    Now we observe the conjugate of $(n,h)\in G$. For $(m,k)\in G$,
    \begin{align*}
        (m,k)^{-1}(n,h)(m,k) &= (-t^{-k}m,-k)(n,h)(m,k) \\
        &= (-t^{-k} m + t^{-k}n,-k+h)(m,k) \\
        &= (t^{-k}n+t^{-k}(t^h-1)m,h).
    \end{align*}
    If $t=1$, there are $42$ conjugacy classes. If $t=6$, there are $24$ conjugacy classes. If $t=2,4$, there are $18$ conjugacy classes. If $t=3,5$, there are $6$ conjugacy classes. As the cases $t=2$ and $t=4$ are isomorphic by
    \begin{align*}
        (n,h)\mapsto (n,-h)
    \end{align*}
    and the cases $t=3$ and $t=5$ are isomorphic by
    \begin{align*}
        (n,h)\mapsto (n,-h),
    \end{align*}
    there are $4$ isomorphism classes.
\end{example}

\begin{example}
    We now observe the number of isomorphism classes of groups of order $n=6,7,8,9,10,11,12$. First, for $n=7,11$, as $n$ is prime, the only possible group of order $n$ is $\Z/n\Z$. \\
    
    For case $|G|=6$, note that there exists a $3$-Sylow subgroup $N\le G$. As $(G:N)=2$, $N$ is a normal subgroup, and we can write $G$ as an inner semidirect product of $N$ and $G/N$, where $G/N$ denotes the group of right cosets. As $|N|=3$ and $|G/N|=2$, $N\simeq \Z/3\Z$ and $G/N\simeq \Z/2\Z$. So it is sufficient to consider $\Z/3\Z \rtimes \Z/2\Z$. The only possibilities for the group homomorphism
    \begin{align*}
        \phi:\Z/2\Z\rightarrow \mathrm{Aut}(\Z/3\Z)
    \end{align*}
    are $\phi_1(n)=n$ and $\phi_1(n)=2n$. The first case corresponds to $\Z/6\Z$ and the second case corresponds to $S_3$. \\

    For case $|G|=10$, it is similar; there exists a $5$-Sylow subgroup $N\le G$, and the normality of $N$ comes from the fact $(G:N)=2$, so we can write $G$ as an inner semidirect product of $N$ and $G/N$. As $|N|=5$ and $|G/N|=2$, $N\simeq \Z/5\Z$ and $G/N\simeq \Z/2\Z$. So it is sufficient to consider $\Z/5\Z\rtimes \Z/2\Z$. The only possibilities for the group homomorphism
    \begin{align*}
        \phi:\Z/2\Z\rightarrow \mathrm{Aut}(\Z/5\Z)
    \end{align*}
    are $\phi_1(n)=n$ and $\phi_1(n)=4n$. The first case corresponds to $\Z/10\Z$ and the second case corresponds to $D_{10}$.\\

    For case $|G|=9$, every element is of order $1$, $3$, or $9$. If there exists an element of order $9$, then $G$ is cyclic and is isomorphic to $\Z/9\Z$. If every element is of order $1$ or $3$, then let $a\in G\setminus \{1\}$ and $b\in G\setminus \langle a\rangle$. Then by the law of cancellation,
    \begin{align*}
        1,a,a^2,b,ab,ab^2,b^2,ab^2,a^2b^2
    \end{align*}
    are distinct elements of $G$, and these exactly form the group $G$. This implies $G\simeq \Z/3\Z\times\Z/3\Z$.\\

    For case $|G|=12$, we observe that there exists a $3$-Sylow subgroup, which is of order $3$. By Sylow theorem, the number of the $3$-Sylow subgroup $n_3$ satisfies
    \begin{align*}
        n_3\equiv 1\mod 3,\quad n_3|4.
    \end{align*}
    Hence $n_3=1$ or $n_3=4$. First, if $n_3=1$, then there is a unique $3$-Sylow subgroup $N$, and it is normal in $G$. Then we can express $G$ as a semidirect product $N\rtimes (G/N)$. Here, $|N|=3$ implies $N\simeq \Z/3\Z$ and $|(G/N)|=4$ implies $G/N\simeq \Z/4\Z$ or $G/N\simeq \Z/2\Z\times\Z/2\Z$.
    \begin{itemize}
        \item If $G/N\simeq \Z/4\Z$, then $G\simeq \Z/3\Z\rtimes\Z/4\Z$. The only possibilities for the group homomorphism
        \begin{align*}
            \phi:\Z/4\Z\rightarrow \mathrm{Aut}(\Z/3\Z)
        \end{align*}
        are $\phi_1(n)=n$ and $\phi_1(n)=2n$. The first case corresponds to $\Z/12\Z$ and the second case corresponds to $\mathrm{Dic}_{3}$.
        \item If $G/N\simeq \Z/2\Z\times\Z/2\Z$, then $G\simeq \Z/3\Z\rtimes(\Z/2\Z\times\Z/2\Z)$. The only possibilities for the group homomorphism
        \begin{align*}
            \phi:\Z/2\Z\times\Z/2\Z\rightarrow \mathrm{Aut}(\Z/3\Z)
        \end{align*}
        are 
        \begin{alignat*}{3}
            \phi_{(1,0)}(n)&=n,\quad &&\phi_{(0,1)}(n) &&= n, \\
            \phi_{(1,0)}(n)&=n,\quad &&\phi_{(0,1)}(n) &&= 2n, \\
            \phi_{(1,0)}(n)&=2n,\quad &&\phi_{(0,1)}(n) &&= n, \\
            \phi_{(1,0)}(n)&=2n,\quad &&\phi_{(0,1)}(n) &&= 2n.
        \end{alignat*}
        The first case corresponds to $\Z/6\Z\times\Z/2\Z$. The second, third, and fourth cases lead to an isomorphic result, and the result corresponds to $D_{12}$.
    \end{itemize}
    If $n_3=4$, then note that $G$ has $8$ elements of order $3$. As there are $4$ elements not having $4$ as its order, and as a $2$-Sylow subgroup of $G$ is of order $4$, $G$ can have at most one $2$-Sylow subgroup $N$. The uniqueness of the $2$-Sylow subgroup gives the normality of $N$. Then we can express $G$ as a semidirect product $N\rtimes (G/N)$. As $|N|=4$, $N\simeq \Z/4\Z$ or $N\simeq\Z/2\Z\times\Z/2\Z$, and as $|G/N|=3$, $G/N\simeq \Z/3\Z$.
    \begin{itemize}
        \item If $N\simeq \Z/4\Z$, then $G\simeq \Z/4\Z\rtimes\Z/3\Z$. The only possibility for the group homomorphism
        \begin{align*}
            \phi:\Z/3\Z\rightarrow \mathrm{Aut}(\Z/4\Z)
        \end{align*}
        is $\phi_1(n)=n$. In fact, in this case we obtain $n_3=1$.
        \item If $N\simeq \Z/2\Z\times\Z/2\Z$, then $G\simeq (\Z/2\Z\times\Z/2\Z)\rtimes\Z/3\Z$. The only possibilities for the group homomorphism
        \begin{align*}
            \phi:\Z/3\Z\rightarrow \mathrm{Aut}(\Z/2\Z\times\Z/2\Z)
        \end{align*}
        are $\phi_1(m,n)=(m,n)$, $\phi_1(m,n)=(m+n,m)$ and $\phi_1(m,n)=(n,m+n)$. The first case leads to $n_3=1$. The second and the third cases lead to an isomorphic result $A_4$.
    \end{itemize}

    For case $|G|=8$, observe that every element is of order $1$, $2$, $4$, or $8$. If every element is of order $1$ or $2$, then consider a nontrivial element $a\in G$ and $b\in G\setminus \langle a\rangle$. Observe that
    \begin{align*}
        \langle a,b\rangle = \{1,a,b,ab,ba,aba,bab,abab,baba,\cdots\}.
    \end{align*}
    As we assumed that every element of $G$ is of order $1$ or $2$, $abab = (ab)^2 = 1$ and $baba=(ba)^2=1$. Hence $aba = b$ and $bab=a$, and we obtain
    \begin{align*}
        \langle a,b\rangle = \{1,a,b,ab,ba\}.
    \end{align*}
    As it is a subgroup of $G$, the order of $|\langle a,b\rangle|$ should divide $8$, and this implies $ab=ba$ and $\langle a,b\rangle \simeq \Z/2\Z\times\Z/2\Z$. As $(G:\langle a,b\rangle)=2$, $\langle a,b\rangle$ is a normal subgroup of $G$ and we can view $G$ as a semidirect product $(\Z/2\Z\times\Z/2\Z)\rtimes \Z/2\Z$. The only possibilities for the group homomorphism
        \begin{align*}
            \phi:\Z/2\Z\rightarrow \mathrm{Aut}(\Z/2\Z\times\Z/2\Z)
        \end{align*}
        are $\phi_1(m,n)=(m,n)$ and $\phi_1(m,n)=(n,m)$. The first case corresponds to $\Z/2\Z\times\Z/2\Z\times\Z/2\Z$ and the second case corresponds to $\mathrm{Dic}_2$. Otherwise, if there exists an element of order $4$, there exists a cyclic subgroup $N=\{1,a,a^2,a^3\}$ of order $4$ generated by that element. As $(G:N)=2$, $N$ is normal in $G$. Let the cosets of $N$ in $G$ be $N$ and $bN$. If $b^2=a$ or $b^2=a^3$, then $b$ is of order $8$ and $G\simeq \Z/8\Z$. If $b^2 = a^2$, we obtain
        \begin{align*}
            G=\{1,a,a^2,a^3,b,ab,a^2b=b^3,a^3b\}
        \end{align*}
        and consider the possibilities
        \begin{align*}
            ba=ab,\quad ba=a^2b,\quad ba=a^3b.
        \end{align*}
        The first case implies that $G$ is abelian, and there are no such groups satisfying $b\notin \langle a\rangle$ and $a^2=b^2$. The second case leads to $ba=b^3$ and $a=b^2=a^2$, which is a contradiction. The third case leads to $ba=ab^3=aba^2$ and $1=a^2$, which is also a contradiction. If $b^2=1$, the subgroups $N$ and $\{1,b\}$ are disjoint and we can write $G$ as a semidirect product $\Z/4\Z\rtimes\Z/2\Z$. The only possibilities for the group homomorphism
        \begin{align*}
            \phi:\Z/2\Z\rightarrow \mathrm{Aut}(\Z/4\Z)
        \end{align*}
        are $\phi_1(n)=n$ and $\phi_1(n)=3n$. The first case corresponds to $\Z/4\Z\times\Z/2\Z$ and the second case corresponds to $D_4$.
\end{example}


\section{Towers of groups}
\begin{definition}
    Let $G$ be a group. A \emph{commutator} is an element of the form
    \begin{align*}
        [g,h]:=ghg^{-1}h^{-1}.
    \end{align*}
    The \emph{commutator subgroup} $[G,G]$ is the subgroup generated by commutators.
\end{definition}
\begin{remark}
    The alternative definition
    \begin{align*}
        [g,h]':=g^{-1}h^{-1}gh
    \end{align*}
    is also possible. It is essentially the same as the original definition, as $[g^{-1},h^{-1}]'=[g,h]$.
\end{remark}
\begin{remark}
    Note that $[G,G]$ is the subgroup \emph{generated} by commutators. It is not guaranteed that every element of $[G,G]$ is a commutator.
\end{remark}

\begin{proposition}
    $[G,G]\le G$ is a normal subgroup. The quotient group $G/[G,G]$ is abelian.
\end{proposition}
\begin{proof}
    For $g\in G$,
    \begin{align*}
        g \left( \prod_{i=1}^n [g_i,h_i] \right)g^{-1} &= g \left( \prod_{i=1}^n g_{i}h_ig_i^{-1}h_i^{-1} \right)g^{-1} \\
        &= \prod_{i=1}^n gg_ih_ig_i^{-1}h_i^{-1}g^{-1} \\
        &= \prod_{i=1}^n [gg_ig^{-1},gh_ig^{-1}],
    \end{align*}
    so $[G,G]$ is normal in $G$. For $g,h\in G$,
    \begin{align*}
        [G,G]ghg^{-1}h^{-1} = [G,G]
    \end{align*}
    and
    \begin{align*}
        [G,G]gh = [G,G]hg,
    \end{align*}
    so $G/[G,G]$ is abelian.
\end{proof}

\begin{theorem}
    If $A$ is an abelian group and $f:G\rightarrow A$ is a homomorphism, then there uniquely exists a homomorphism $\phi:G/[G,G]\rightarrow A$. This gives us a universality:
    \begin{figure}[H]
    \centering
    \begin{tikzcd}
        G \arrow[dr, "i"'] \arrow[rr, "\pi"] & & A \\
        & G/[G,G] \arrow[ur, dashed, "\phi"'] &
    \end{tikzcd}
    \end{figure}
\end{theorem}
\begin{remark}
    With this reason, we also call $\pi:G\rightarrow [G,G]$ the \emph{abelianization} of $G$.
\end{remark}

\begin{definition}
    The \emph{derived series} is a sequence of subgroups
    \begin{align*}
        G \ge G^1 \ge G^2 \ge G^3 \ge \cdots
    \end{align*}
    with $G^{n+1} := [G^n,G^n]$.
\end{definition}
\begin{remark}
    $[G,G]$ is often called the \emph{derived subgroup} of $G$.
\end{remark}

\begin{definition}
    A group $G$ is \emph{solvable} if $G^n=\{1\}$ for some $n$.
\end{definition}
\begin{remark}
    Such groups are called solvable as it is related to solving polynomials (with radicals). The related contents will be introduced later.
\end{remark}

\begin{theorem}\label{Theorem 1.13.6.}
    $G$ is solvable if and only if for some $n$, there exists a series
    \begin{align*}
        G=G_0\ge G_1\ge G_2\ge \cdots\ge G_n = 1
    \end{align*}
    such that $G_{n+1}$ is normal in $G_i$ and $G_i/G_{i+1}$ is abelian for all $i$.
\end{theorem}
\begin{proof}
    If $G$ is solvable, the derived series satisfies the given condition, as $[G,G]$ is normal in $G$ and $G/[G,G]$ is abelian for all $G$. Conversely, assume that there exists a finite sequence $\{G_i\}$ satisfying the given condition. By the universality of commutators, we obtain $G/G_1\le G/[G,G]$ and $G^1 \le G_1$. Inductively, if $G^n\le G_n$, then
    \begin{align*}
        G_{n+1} \ge [G_n,G_n] \ge [G^n,G^n]=G^{n+1}.
    \end{align*}
    As $G_n=1$ for some $n$, $G^n=1$ for some $n$ and $G$ is solvable.
\end{proof}
\begin{remark}
    This is often used as an alternative definition for a solvable group.
\end{remark}

\begin{example}
    Let $k$ be a field and let $k^\times$ act on $k$ by multiplication. We now think of the semidirect product $G=k\rtimes k^\times$. Observe the sequence
    \begin{align*}
        k\rtimes k^\times \ge k \ge 1.
    \end{align*}
    By the definition of the semidirect product, $k$ is normal in $k\rtimes k^\times$. As $k\rtimes k^\times / k \simeq k^\times$ and $k^\times$ is abelian, $k\rtimes k^\times / k$ is abelian. It is clear that $1$ is normal in $k$ and $k/1\simeq k$ is abelian. By Theorem \ref{Theorem 1.13.6.}, $G$ is solvable.
\end{example}

\begin{definition}
    Let $G$ be a group. The \emph{lower central series} is defined by $G=G^{(1)}$ and
    \begin{align*}
        G^{(n+1)}=[G,G^{(n)}].
    \end{align*}
\end{definition}

\begin{definition}
    A group $G$ is \emph{nilpotent} if $G^{(n)}=1$ for some $n$.
\end{definition}

\begin{example}
    Above we have shown that for a field $k$, $G=k\rtimes k^\times$ is solvable. However, it is not nilpotent. The lower central series of $G$ is written as
    \begin{align*}
        G=k\rtimes k^\times \ge k \ge k \ge \cdots
    \end{align*}
    unless $k^\rtimes$ is trivial.
\end{example}
\begin{remark}
    Every nilpotent group is solvable by definition.
\end{remark}

\begin{example}
    We observe the lower central series of $U(n,k)$. Here, observe that
    \begin{align*}
        \begin{pmatrix}
            1 & a_1 & * & * & \cdots & * \\
            0 & 1 & a_2 & * & \cdots & * \\
            \vdots & \vdots & \ddots & & & \vdots \\
            0 & 0 & \cdots & & 1 & a_{n-1} \\
            0 & 0 & \cdots & & 0 & 1
        \end{pmatrix} \begin{pmatrix}
            1 & b_1 & * & * & \cdots & * \\
            0 & 1 & b_2 & * & \cdots & * \\
            \vdots & \vdots & \ddots & & & \vdots \\
            0 & 0 & \cdots & & 1 & b_{n-1} \\
            0 & 0 & \cdots & & 0 & 1
        \end{pmatrix} \\= \begin{pmatrix}
            1 & a_1 + b_1 & * & * & \cdots & * \\
            0 & 1 & a_2 + b_2 & * & \cdots & * \\
            \vdots & \vdots & \ddots & & & \vdots \\
            0 & 0 & \cdots & & 1 & a_{n-1} + b_{n-1} \\
            0 & 0 & \cdots & & 0 & 1
        \end{pmatrix}.
    \end{align*}
    Denote the diagonals on the right of the diagonal of the matrix as the first, second, third, ... diagonals. By the observation above, we know that $[G,G]$ consists of unipotent matrices having zeros on the first diagonal. With further computation, we can inductively show that $[G,G^{(i)}]$ consists of unipotent matrices having zeros on the first, second, ... $(i-1)$-th diagonals. Hence $G^{(n)}=[G,G^{(n-1)}]=1$ and $G$ is nilpotent.
\end{example}

\begin{example}
    As every $p$-group is isomorphic to a subgroup of $U(n,\Z/p\Z)$, by the example above we obtain that every $p$-group is nilpotent.
\end{example}

\begin{theorem}
    A finite nilpotent group is a direct product of its Sylow subgroups.
\end{theorem}
\begin{proof}
    Let $G$ be a finite nilpotent group. First, we show that for every proper subgroup $H$, $N_G(H)\neq H$. Suppose there exists a proper subgroup $H$ such that $N_G(H)=H$. Then there exists $g\in G$ such that
    \begin{align*}
        gHg^{-1} \ge H
    \end{align*}
    and hence for all $K\ge H$, $[G,K]\ge H$. However, as $G$ is nilpotent, there exists some $n$ such that $G^{(n)}\ge H$ and $G^{(n+1)}\not\ge H$. This leads to a contradiction.

    Now, suppose that $P$ is a Sylow $p$-subgroup of $G$. Consider
    \begin{align*}
        N_G(P),N_G(N_G(P)),\cdots.
    \end{align*}
    By definition of normalizers, $P$ is normal in $N_G(P)$. If $x\in N_G(N_G(P))$, then
    \begin{align*}
        xPx^{-1}\le xN_G(P)x^{-1} = N_G(P).
    \end{align*}
    As $P$ and $xPx^{-1}$ are Sylow $p$-subgroups of $N_G(P)$ and as every Sylow $p$-subgroups are conjugate, there exists $y\in N_G(P)$ such that
    \begin{align*}
        xPx^{-1} = yPy^{-1}=P.
    \end{align*}
    Therefore $x\in N_G(P)$, which implies $N_G(N_G(P))=N_G(P)$. As $N_G(P)=N_G(N_G(P))$, $N_G(P)$ is not a proper subgroup of $G$ and $N_G(P)=G$. Hence $P$ is normal in $G$.

    Let $|G|=p_1^{e_1}p_2^{e_2}\cdots p_n^{e^n}$ and let $P_i$ be a Sylow $p_i$-subgroup. As $P_i$'s are normal and are all disjoint, if $i\neq j$ and $a\in P_i$, $b\in P_j$, then
    \begin{align*}
        aba^{-1}b^{-1}\in P_i\cap P_j = 1
    \end{align*}
    and $ab=ba$. Therefore $G$ is equivalent to $P_1\times P_2\times\cdots\times P_n$.
\end{proof}
\begin{remark}
    The converse also holds, as every $p$-group is nilpotent and the direct product of nilpotent groups is nilpotent.
\end{remark}


\section{Jordan-H\"older}
\begin{definition}
    A series
    \begin{align*}
        G=G_0\ge G_1\ge G_2\ge \cdots
    \end{align*}
    is called \emph{subnormal} if $G_{i+1}$ is normal in $G_i$ for all $i$.
\end{definition}

\begin{definition}
    A group $H$ is called \emph{simple} if there are exactly two normal subgroups $\{1\}$ and $H$.
\end{definition}

\begin{definition}
    A subnormal series is called a \emph{composition series} if $G_n=\{1\}$ for some $n$ and $G_i/G_{i+1}$ is simple for all $i$.
\end{definition}

\begin{example}
    We find a composition series for $S_3$. The normal subgroups of $S_3$ are $1, \langle (123)\rangle,S_3$, and the only subgroup making its quotient simple is $\langle(123)\rangle$. So the composition series is formed as
    \begin{align*}
        S_3\ge \langle(123)\ge 1.
    \end{align*}
\end{example}

\begin{example}
    There exists two composition series for $\Z/10\Z$;
    \begin{align*}
        \Z/10\Z \ge 2(\Z/5\Z)\simeq \Z/5\Z \ge 1
    \end{align*}
    and
    \begin{align*}
        \Z/10\Z\ge 5(\Z/2\Z) \simeq \Z/2\Z\ge 1.
    \end{align*}
\end{example}

\begin{example}
    As $\Z$ is abelian, every subgroup is normal. Every subgroup of $\Z$ can be expressed in the form $\Z/n\Z$. However, for all $n$, $\Z/(\Z/n\Z)\simeq \Z$ is not simple. Therefore $\Z$ does not admit a composition series.
\end{example}

\begin{definition}
    Given a subnormal series $(G_i)_i$ of length $n$, the isomorphism classes of
    \begin{align*}
        \{G_i/G_{i+1}:i=0,1,\dots,n-1\}
    \end{align*}
    form a multiset. Two subnormal series are called \emph{equivalent} if the associated multisets are the same.
\end{definition}

\begin{theorem}[Jordan-H\"older]
    Any two composition series of a group are equivalent.
\end{theorem}
\begin{proof}
    It suffices to show that two subnormal series of a group have equivalent refinements (or subsequences). Let
    \begin{align*}
        &G = G_0\ge G_1\ge\cdots\ge G_n = 1, \\
        &G = H_0\ge H_1\ge\cdots\ge H_n = 1
    \end{align*}
    are two subnormal series of $G$. Let
    \begin{align*}
        G_{i,j} = G_{j+1}(H_i\cap G_j),\quad H_{i,j} = H_{j+1}(G_i\cap H_j).
    \end{align*}
    As
    \begin{align*}
        G_j = G_{j+1}(H_0\cap G_j) \ge G_{j+1}(H_1\cap G_j) \ge \cdots \ge G_{j+1}(H_n\cap G_j) = G_{j+1}
    \end{align*}
    forms a subnormal series of $G_j$, we can find a refinement of the first series:
    \begin{align*}
        G &= G_{0,0} \ge G_{1,0} \ge \cdots \ge G_{n,0} \\
        &= G_{0,1} \ge G_{1,1} \ge \cdots \ge G_{n,1} \\
        &\ \ \vdots \\
        &= G_{0,n-1} \ge G_{1,n-1} \ge \cdots \ge G_{n,n-1}=1
    \end{align*}
    We can similarly find a refinement of the second series. As
    \begin{align*}
        G_{j+1}(H_i \cap G_j) / G_{j+1}(H_{i+1}\cap G_j) \simeq H_{i+1}(G_j\cap H_i)/H_{i+1}(G_{j+1}\cap H_i),
    \end{align*}
    the refinements are equivalent, and we obtain the desired result.
\end{proof}


\clearpage
\section{Symmetric and alternating groups}

\begin{definition}
    Let $n\ge 1$ and $S_n$ be the symmetric group on $\{1,2,\dots,n\}$. A \emph{transposition} is $(ij)$ with $i\neq j$.
\end{definition}

\begin{theorem}
    There is a unique group homomorphism sign $\mathrm{sign}:S_n\rightarrow\{\pm1\}$ such that $\mathrm{sign}(\tau)=-1$ for every transposition $\tau$.
\end{theorem}
\begin{proof}
    Let $\Delta$ be the function
    \begin{align*}
        \Delta(x_1,x_2,\dots, x_n)=\prod_{i<j}(x_j-x_i).
    \end{align*}
    Let $\tau=(rs)$ be a transposition, where $r<s$. Consider
    \begin{align*}
        \tau\Delta(x_1,\dots,x_n)=\prod_{i<j}(x_{\tau(j)}-x_{\tau(i)}).
    \end{align*}
    Note that $\tau$ changes the factor $(x_s-x_r)$ to $-(x_s-x_r)$. The factors
    \begin{align*}
        (x_k-x_s)(x_k-x_r),&\quad k>s, \\
        (x_s-x_k)(x_k-x_r),&\quad r<k<s, \\
        (x_s-x_k)(x_r-x_k),&\quad k<r
    \end{align*}
    are preserved. Hence $\tau\Delta = -\Delta$. 

    For $\sigma\in S_n$, let $\mathrm{sign}(\sigma)\in\{\pm 1\}$ be the one satisfying
    \begin{align*}
        \sigma\Delta = \mathrm{sign}(\sigma)\Delta.
    \end{align*}
    As for $\sigma_1,\sigma_2\in S_n$, $\sigma_1(\sigma_2\Delta)=(\sigma_1\sigma_2)\Delta$, so $\mathrm{sign}$ is a homomorphism. This is the homomorphism satisfying the desired result. The uniqueness comes from the fact that every element in $S_n$ can be expressed as a product of transpositions.
\end{proof}
\begin{remark}
    There exists a simple proof: note that $\mathrm{sign}(\sigma)$ is a determinant of a permutation matrix of $\sigma$.
\end{remark}

\begin{definition}
    The \emph{alternating group} on $n$ letters is defined as
    \begin{align*}
        A_n := \ker (\mathrm{sign}).
    \end{align*}
\end{definition}

\begin{definition}
    A \emph{partition} of $n$ refers to a finite sequence of non-increasing positive integers
    \begin{align*}
        m_1 \ge m_2 \ge \cdots \ge m_k \ge 1
    \end{align*}
    such that $m_1+\cdots+m_k = n$. Let $p(n)$ be the number of partitions of $n$.
\end{definition}

\begin{example}
    Every $\sigma\in S_n$ is the product of disjoint cycles. For example, see
    \begin{align*}
        (123)(46)(57)(8)\in S_8.
    \end{align*}
    This cycle decomposition gives rise to a bijection
    \begin{align*}
        \{\mathrm{conjugacy~classes~in~}S_n\}\simeq \{\mathrm{partitions~of~}n\}.
    \end{align*}
    Note that for $\sigma\in S_n$ and a cycle $(a_1a_2\cdots a_m)$,
    \begin{align*}
        \sigma(a_1a_2\cdots a_m)\sigma^{-1}=(\sigma(a_1)\sigma(a_2)\cdots\sigma(a_m)).
    \end{align*}
    Hence the conjugacy preserves the partition.
\end{example}

\begin{proposition}
    The set of all cycles of length $n$ forms a conjugacy class, which is an $A_n$-set under the conjugation action. This $A_n$-set is transitive if and only if $n$ is even. If $n$ is odd, it has exactly two orbits.
\end{proposition}
\begin{proof}
    As $(S_n:A_n)=2$, Each conjugacy class viewed as an $A_n$-set has at most two orbits. If $n$ is even, consider two cycles of length $n$
    \begin{align*}
        \sigma_1 = (a_1\dots a_n), \sigma_2 = (b_1\dots b_n).
    \end{align*}
    Think of $\tau_1, \tau_2\in S_n$ such that
    \begin{align*}
        \tau_1(a_i)=b_i,\quad \tau_2(a_i)=b_{i+1}
    \end{align*}
    where $b_{n+1}=b_1$. As $\tau_2 = \sigma_2\tau_1$ and $\sigma_2\in S_n\setminus A_n$, at least one of $\tau_1$ and $\tau_2$ is in $A_n$. As
    \begin{align*}
        \tau_1^{-1}\sigma_1\tau_1 = \tau_2^{-1}\sigma_1\tau_2 = \sigma_2,
    \end{align*}
    we obtain that every two cycles of length $n$ is in the same orbit, and the given $A_n$-set is transitive.

    If $n$ is odd, note that
    \begin{align*}
        (12\dots (n-1)n),(12\dots n(n-1))
    \end{align*}
    are not conjugate. Hence there exists exactly two orbits.
\end{proof}

\begin{proposition}\label{Proposition 1.15.7.}
    If $n\ge 6$, the conjugacy class containing a non-trivial element $\sigma\in S_n$ has at least $2n$ elements.
\end{proposition}
\begin{proof}
    If $\sigma$ is a cycle of length $k\ge 2$, then the number of conjugacy class is
    \begin{align*}
        {n\choose k}\times (k-1)! = \frac{n!}{(n-k)!k}= \frac{n(n-1)\cdots(n-k+1)}{k}\ge 2n.
    \end{align*} 
    If $\sigma$ contains a cycle of length $k$ longer than $n/2$, then there should be exactly one cycle of length $k$ and the number of conjugacy class is larger than
    \begin{align*}
        {n\choose k}\times (k-1)! = \frac{n!}{(n-k)!k}= \frac{n(n-1)\cdots(n-k+1)}{k}\ge 2n.
    \end{align*}
    Now it suffices to show that the assertion holds for all $\sigma=\sigma_1\sigma_2$ where
    \begin{align*}
        1\neq \sigma_1 = \mathrm{Perm}(\{1,2,\dots,n_1\})
    \end{align*}
    and
    \begin{align*}
        1\neq \sigma_2 \in \mathrm{Perm}(\{n_1+1,n_1+2,\dots,n_1+n_2\})
    \end{align*}
    with $n_1+n_2=n$ and $n_1\ge n_2\ge 3$. Here we use induction; assume the assertion holds for $i<n$. Then the conjugacy class containing $\sigma_1$ has at least $2n_1\ge n$ elements. If $\sigma_2$ is a cycle $(a_1\dots a_{n-1} a_n)$, we can find an element in the conjugation class $(a_1\dots a_n a_{n-1})$, and if $\sigma_2$ is not a cycle, then we can find an element in the conjugation class, exchanging one element from two different partitions. Hence the conjugacy class containing $\sigma_2$ has at least $2$ elements. As $\sigma_1$ and $\sigma_2$ are disjoint, we conclude that the conjugation class containing $\sigma=\sigma_1\sigma_2$ has at least $2n$ elements.
\end{proof}

\begin{proposition}
    If $n\ge 5$, the conjugacy class of a non-trivial element in $A_n$ contains at least $n$ elements.
\end{proposition}
\begin{proof}
    Let $C\subset S_n$ be a conjugacy class of a non-trivial element in $S_n$. As $C$ is transitive with respect to $S_n$ and $(S_n:A_n)=2$, $C$ has at most two orbits under the $A_n$-action. If $n\ge 6$, by Proposition \ref{Proposition 1.15.7.} $C$ has at least $2n$ elements, and each orbit of $C$ with respect to $A_n$ has at least $n$ elements. Note that each orbit of $C$ as an $A_n$-set is a conjugacy class in $A_n$. For $n=5$, we can directly compute the (nontrivial) conjugacy classes:
    \begin{itemize}
        \item $(a_1a_2)(a_3a_4)$ - $15$ elements
        \item $(a_1a_2a_3)$ - $20$ elements
        \item $(a_1a_2a_3a_4a_5)$ - two conjugacy classes, each of $12$ elements
    \end{itemize}
    Therefore we obtain the result.
\end{proof}

\begin{theorem}
    For $n\ge 5$, $A_n$ is simple.
\end{theorem}
\begin{proof}
    Suppose that $N\le A_n$ is a normal subgroup of size smaller than or equal to $n$. Then as a normal subgroup is a union of conjugacy classes, and as every nontrivial conjugacy class contains at least $n$ elements, $N$ is trivial. (Note that $N$ also contains a trivial conjugacy class.)

    Now we use induction on $n$ Suppose that $A_k$ is simple for $5\le k\le n-1$. Let $N\le A_n$ be a normal subgroup, and let $H_i\le A_n$ be the isotropy subgroup of $i$ (the set of even permutations fixing $i$). If $\sigma\in N\cap H_i$, then for all $\tau\in H_i$ $\tau^{-1}\sigma\tau\in N$. So $N\cap H_i$ is a normal subgroup of $H_i$. Note that
    \begin{align*}
        H_i=A_n\cap \mathrm{Perm}\{1,2,\dots,i-1,i+1,\dots,n\}\simeq A_{n-1}
    \end{align*}
    and as $A_{n-1}$ is even, we obtain that $N\cap H_i=1$ or $N\cap H_i=H_i$. If $N\cap H_i=H_i$ for some $H_i$, then as for all $\sigma\in A_n$ $H_{\sigma(i)}$ is a conjugate of $H_i$, $N$ contains all $H_1,\dots,H_n$. As these subgroups generates $A_n$, $N=A_n$. On the other hand, if $N\cap H_i=1$ for all $i$, then consider the embedding
    \begin{alignat*}{2}
        f:\ &N &&\rightarrow \{1,2,\dots,n\} \\
        &\sigma &&\mapsto \sigma(1).
    \end{alignat*}
    If $\sigma,\tau\in N$ and $f(\sigma)=f(\tau)$, then $\sigma(\tau^{-1})(1)=1$ and $\sigma\circ \tau^{-1}=1$. Hence $f$ is injective, and we conclude that $|N|\le n$. By the first statement of our proof, $N$ is trivial.

    Therefore $A_n$ is simple, and with induction we obtain the result.
\end{proof}


\section{\texorpdfstring{Simplicity of $\mathrm{PSL}(2,\Z/p\Z)$}{Simplicity of PSL(2)}}
\begin{definition}
    Let $k$ be a field and $n\ge 2$ an integer. The \emph{projective special linear group} $\mathrm{PSL}(n,k)$ is defined as
    \begin{align*}
        \mathrm{PSL}(n,k):=\mathrm{SL}(n,k)/Z
    \end{align*}
    where $\mathrm{SL}(n,k)$ is the set of matrices of determinant $1$ and $Z$ consists of scalar matrices in $\mathrm{SL}(n,k)$.
\end{definition}

In this section, we prove that $\mathrm{PSL}(2,\Z/p\Z)$ is simple. Before proof, we introduce two useful propositions.

\begin{definition}
    Let $G$ be a group and $X$ be a $G$-set. $X$ is said to be \emph{doubly transitive} if for all $x_1,x_2,y_1,y_2\in X$ satisfying $x_1\neq x_2$ and $y_1\neq y_2$, there exists $g\in G$ such that $gx_1=y_1$ and $gx_2=y_2$.
\end{definition}

\begin{proposition}\label{Proposition 1.16.3.}
    Let $G$ be a group which acts doubly transitively on a set $X$. For any $x\in X$, $G_x\le G$ is a maximal proper subgroup.
\end{proposition}
\begin{proof}
    Let $x\in X$ and $h\in G-G_x$. We first show that
    \begin{align*}
        G=G_x \cup G_x h G_x.
    \end{align*}
    As $G$ is doubly transitive on $X$, $G_xhx = X-\{x\}$. For all $g\not\in G_X$, there exists some $u\in G_x$ such that $gx = uhx$, and there exists some $v\in G$ such that $v = u^{-1}h^{-1}g$ and $g=uhv$. As $v$ should fix $x$, we obtain that $v\in G_x$. Hence $G=G_x \cup G_x h G_x$.

    Now, let $G_x < K \le G$. We can choose $h\in K-G_x$. Then as
    \begin{align*}
        G = G_x\cup G_x h G_x \le K,
    \end{align*}
    we obtain that $K=G$. Therefore $G_x$ is a maximal proper subgroup.
\end{proof}

\begin{proposition}\label{Proposition 1.16.4.}
    Let $G$ be a group which acts doubly transitively on a set $X$. Let $X$ be a faithful $G$-set. Assume for some $x\in X$ the group $G_x$ has an abelian normal subgroup $A$ whose conjugate subgroups generate $G$, and $G=[G,G]$. Then $G$ is simple.
\end{proposition}
\begin{proof}
    Let $N$ be a proper normal subgroup of $G$. Let $x\in X$ and let $G_x$ be the isotropy subgroup. We first show that $N\le G_x$ or $NG_x = G$. If $N\not\le G_x$, then there exists some $h\in N-G_x$. Then by Proposition \ref{Proposition 1.16.3.} we know that
    \begin{align*}
        G = G_x \cup G_x h G_x
    \end{align*}
    and for all $g\in G$, $g\in G_x\le NG_x$, or $g\in G_xhG_x$ and for some $u,v\in G_x$,
    \begin{align*}
        g = uhv \in Nuv \le NG_x.
    \end{align*}
    Hence $G=NG_x$.
    
    If $N\le G_x$, note that $gN=Ng$ for all $g\in G$, and for all $x\in X$,
    \begin{align*}
        Ngx = gNx = gx.
    \end{align*}
    As $G$ is transitive on $X$, $gx$ can be arbitrary element in $X$, and this implies that every element of $N$ fixes every element of $X$. As $G$ is faithful, we obtain $N=1$.

    If $G=NG_x$, then for all $g\in G$ there exists some $g'\in G_x$ such that
    \begin{align*}
        g \in Ng'Ag'^{-1} = NAg'g'^{-1} = NA.
    \end{align*}
    Hence $G=NA$. This implies that $A=G/N$ is abelian, and by the universality of commutators we obtain $G/[G,G]\ge G/N$ and $N\ge [G,G]$. AS $[G,G]=G$, we obtain $N=G$. This contradicts with the fact that $N$ is a proper subgroup of $G$.

    Therefore every proper normal subgroup of $G$ is trivial, and $G$ is simple.
\end{proof}

\begin{theorem}
    For $p\ge 5$, $G=\mathrm{PSL}(2,\Z/p\Z)$ is simple.
\end{theorem}
\begin{proof}
    Let $\mathbb{P}^1(\Z/p\Z)$ denotes the set of $1$-dimensional subspaces in $(\Z/p\Z)^2$. Then $G$-action on $\mathbb{P}^1(\Z/p\Z)$ is doubly transitive; if $x_1,x_2,y_1,y_2\in \mathbb{P}^1(\Z/p\Z)$ satisfies $x_1\neq x_2$ and $y_1\neq y_2$, then there exists $g\in G$ such that $gx_1=y_1$ and $gx_2=y_2$. To prove, let $u_1, u_2,v_1,v_2$ be the basis vectors for $x_1,x_2,y_1,y_2$. Then the matrix
    \begin{align*}
        M = 
        \begin{pmatrix}
            v_1 & v_2
        \end{pmatrix}
        \begin{pmatrix}
            u_1 & u_2
        \end{pmatrix}^{-1} \in \mathrm{SL}(2,\Z/p\Z)
    \end{align*}
    sends $u_1$ to $v_1$ and $u_2$ to $v_2$. Note that for any scalar matrices $\lambda I$, $M(\lambda I)$ sends $u_1$ to $\lambda v_1$ and $u_2$ to $\lambda v_2$, and thus sends $x_1$ and $x_2$ to $y_1$ and $y_2$. Therefore there exists such a $g\in G$.

    Also, $\mathbb{P}^1(\Z/p\Z)$ is a faithful $G$-set, as if $M,N\in \mathrm{SL}(2,\Z/p\Z)$ satisfies
    \begin{align*}
        Mx=Nx\quad \forall x\in \mathbb{P}^1(\Z/p\Z),
    \end{align*}
    then $MN^{-1}=kI$ for some $k\in\R\setminus\{0\}$ and $MZ=NZ$.

    Note that Jordan decomposition tells us that for every matrix $M\in \mathrm{GL}(2,\Z/p\Z)$ there exists an upper triangular matrix $U\in\mathrm{U}(2,\Z/p\Z)$ and an invertible matrix $P$ such that
    \begin{align*}
        M=PUP^{-1}.
    \end{align*}
    Hence $G$ is generated by the conjugates of $\mathrm{U}(2,\Z/p\Z)$, which is a Sylow $p$-subgroup of $G$. Also, note that
    \begin{align*}
        \begin{pmatrix}
            a & 0 \\ 0 & 1/a
        \end{pmatrix} \begin{pmatrix}
            1 & b \\ 0 & 1
        \end{pmatrix} \begin{pmatrix}
            a & 0 \\ 0 & 1/a
        \end{pmatrix}^{-1} \begin{pmatrix}
            1 & b \\ 0 & 1
        \end{pmatrix}^{-1} = \begin{pmatrix}
            1 & b(a^2 - 1) \\ 0 & 1
        \end{pmatrix}.
    \end{align*}
    As $p\ge 5$, we can pick $a\in \Z/p\Z$ such that $a^2-1\neq 0$. Then $\mathrm{U}(2,\Z/p\Z)\le [G,G]$ under isomorphism. As conjugates of $\mathrm{U}(2,\Z/p\Z)$ generate $G$, we obtain $G=[G,G]$.

    Now we know that $G$ and $\mathbb{P}^1(\Z/p\Z)$ satisfies every condition of Proposition \ref{Proposition 1.16.4.}. We therefore conclude that $G$ is simple.
\end{proof}


\clearpage
\section{Direct sum and product}
\begin{proposition}
    Let $X$, $Y$ be sets. The direct product satisfies the universal property:
    \begin{figure}[H]
    \centering
    \begin{tikzcd}
        & X\times Y \arrow[dl, "\pi_X"'] \arrow[dr, "\pi_Y"] & \\
        X & & Y \\
        & Z \arrow[ul, "f_X"] \arrow[ur, "f_Y"'] \arrow[uu, dashed, "\exists !\phi"] &
    \end{tikzcd}
    \end{figure}
    \noindent
    where $f_X$ and $f_Y$ are arbitrary functions and $\phi$ satisfies
    \begin{align*}
        f_X=\pi_X\circ \phi,\quad f_Y=\pi_Y\circ \phi.
    \end{align*}
\end{proposition}

\begin{proposition}
    Let $X$, $Y$ be disjoint sets. The disjoint union satisfies the universal property:
    \begin{figure}[H]
    \centering
    \begin{tikzcd}
        & X\times Y \arrow[dd, dashed, "\exists !\phi"'] & \\
        X \arrow[ur, "i_X"] \arrow[dr, "f_X"'] & & Y \arrow[ul, "i_Y"'] \arrow[dl, "f_Y"] \\
        & Z &
    \end{tikzcd}
    \end{figure}
    \noindent
    where $f_X$ and $f_Y$ are arbitrary functions and $\phi$ satisfies
    \begin{align*}
        f_X=\phi\circ i_X,\quad f_Y=\phi\circ i_Y.
    \end{align*}
\end{proposition}
\begin{remark}
    Direct products and disjoint unions of families of sets also satisfy the universal properties.
\end{remark}

\begin{theorem}
    Let $G_1$, $G_2$ be groups. The direct product satisfies the universal property: there exists a unique $\phi$ satisfying
    \begin{align*}
        f_X=\pi_X\circ \phi,\quad f_Y=\pi_Y\circ \phi,
    \end{align*}
    where $f_1$ and $f_2$ are arbitrary group homomorphisms from a group $Z$.
    \begin{figure}[H]
    \centering
    \begin{tikzcd}
        & G_1\times G_2 \arrow[dl, "\pi_X"'] \arrow[dr, "\pi_Y"] & \\
        G_1 & & G_2 \\
        & Z \arrow[ul, "f_X"] \arrow[ur, "f_Y"'] \arrow[uu, dashed, "\exists !\phi"] &
    \end{tikzcd}
    \end{figure}
\end{theorem}
\begin{remark}
    Same holds for abelian groups. Here we can condition $Z$ to be abelian.
\end{remark}
\begin{proof}
    Let $\phi$ be a map obtained from the universal property of the direct product of sets. It suffices to show that $\phi$ is an isomorphism. For $z_1,z_2\in Z$,
    \begin{align*}
        \phi(z_1)\phi(z_2) &= (\pi_X(\phi(z_1)\phi(z_2)),\pi_Y(\phi(z_1)\phi(z_2))) \\
        &= (\pi_X(\phi(z_1))\pi_X(\phi(z_2)),\pi_Y(\phi(z_1))\pi_Y(\phi(z_2))) \\
        &= (f_X(z_1)f_X(z_2),f_Y(z_1)f_Y(z_2)) \\
        &= (f_X(z_1z_2),f_Y(z_1z_2)) \\
        &= (\pi_X(\phi(z_1z_2)),\pi_Y(\phi(z_1z_2))) \\
        &= \phi(z_1z_2).
    \end{align*}
\end{proof}

\begin{theorem}
    Let $A_1$, $A_2$ be groups. The direct sum satisfies the universal property: there exists a unique $\phi$ satisfying
    \begin{align*}
        f_1=\phi\circ i_1,\quad f_2=\phi\circ i_2,
    \end{align*}
    where $f_1$ and $f_2$ are arbitrary group homomorphisms from an abelian group $Z$.
    \begin{figure}[H]
    \centering
    \begin{tikzcd}
        & A_1\oplus A_2 \arrow[dd, dashed, "\exists !\phi"'] & \\
        A_1 \arrow[ur, "i_1"] \arrow[dr, "f_1"'] & & A_2 \arrow[ul, "i_2"'] \arrow[dl, "f_2"] \\
        & Z &
    \end{tikzcd}
    \end{figure}
    \noindent
    Here, $A_1\oplus A_2$ is the set of formal sums
    \begin{align*}
        a_1+a_2,\quad a_1\in A_1,a_2\in A_2
    \end{align*}
    such that for $a_1,b_1\in A_1$ and $a_2,b_2\in A_2$,
    \begin{align*}
        a_1+b_1=a_2+b_2 \Leftrightarrow (a_1,a_2)=(b_1,b_2).
    \end{align*}
\end{theorem}

\begin{definition}
    Let $\{A_i\}_{i\in I}$ be a family of abelian groups. The direct sum $\bigoplus_{i\in I} A_i$ is defined as the set of finite formal sums
    \begin{align*}
        a_{i_1}+\cdots+a_{i_k},\quad a_{i_j}\in A_{i_j}.
    \end{align*}
\end{definition}
\begin{remark}
    This is isomorphic a subset of $\prod_{i\in I} A_i$ consisting of elements having all but finitely many coordinates as zero.
\end{remark}

\begin{theorem}
    The direct sum $\bigoplus_{i\in I A_i}$ equipped with $i_k:A_k\rightarrow \bigoplus _{i\in I} A_i$ satisfies the universal property: for a given family of homomorphisms
    \begin{align*}
        f_k:A_k\rightarrow Z
    \end{align*}
    there exists a unique
    \begin{align*}
        \phi:\bigoplus_{i\in I}A_i\rightarrow Z
    \end{align*}
    such that
    \begin{align*}
        \phi=f_k\circ i_k.
    \end{align*}
\end{theorem}


\section{Free and finitely generated abelian groups}
\begin{definition}
    We say $A$ is \emph{free} if it is isomorphic to $\bigoplus_{i\in I} \Z$ for some set $I$.
\end{definition}

\begin{proposition}
    Let $T$ be an abelian group. For a given $t\in T$, there exists a unique homomorphism
    \begin{align*}
        f_t:\Z\rightarrow T
    \end{align*}
    such that $f_t(1)=t$.
\end{proposition}
\begin{proof}
    $f_t(1)=t$ gives $f_t(a)=at$. This is indeed a unique homomorphism.
\end{proof}

\begin{definition}
    Let $M\in\mathrm{M}_{m\times n}(\Z)$ be an $m\times n$ matrix over $\Z$. A \emph{Smith decomposition} of $M$ is a factorization
    \begin{align*}
        M=XDY
    \end{align*}
    where $X\in \mathrm{GL}_m(\Z)$, $Y\in\mathrm{GL}_n(\Z)$ and $D$ is in the Smith normal form(SNF); it is defined recursively as follows:
    \begin{enumerate}
        \item A square diagonal $n\times n$ matrix is a SNF if its diagonal entries $d_1,\dots,d_n$ are positive and satisfy $d_i|d_{i+1}$ for $i=1,2,\dots,n-1$.
        \item If you add a zero column to the right of a SNF, then it is a SNF.
        \item If you add a zero row to the bottom of a SNF, then it is a SNF.
    \end{enumerate}
\end{definition}

\begin{example}
    Think of a matrix
    \begin{align*}
        \begin{pmatrix}
            1 & 2 & 3 \\ 2 & -1 & 0 \\ 4 & 5 & 6
        \end{pmatrix}.
    \end{align*}
    we can find a Smith decomposition as
    \begin{align*}
        \begin{pmatrix}
            1 & 0 & 0 \\ 0 & 1 & 0 \\ 0 & 0 & 1
        \end{pmatrix} \begin{pmatrix}
            1 & 2 & 3 \\ 2 & -1 & 0 \\ 4 & 5 & 6
        \end{pmatrix} \begin{pmatrix}
            1 & 0 & 0 \\ 0 & 1 & 0 \\ 0 & 0 & 1
        \end{pmatrix} &= \begin{pmatrix}
            1 & 0 & 0 \\ 2 & 1 & 0 \\ 4 & 0 & 1
        \end{pmatrix} \begin{pmatrix}
            1 & 2 & 3 \\ 0 & -5 & -6 \\ 0 & -3 & -6
        \end{pmatrix} \begin{pmatrix}
            1 & 0 & 0 \\ 0 & 1 & 0 \\ 0 & 0 & 1
        \end{pmatrix} \\
        &= \begin{pmatrix}
            1 & 0 & 0 \\ 2 & 1 & 0 \\ 4 & 0 & 1
        \end{pmatrix} \begin{pmatrix}
            1 & 0 & 0 \\ 0 & -5 & -6 \\ 0 & -3 & -6
        \end{pmatrix} \begin{pmatrix}
            1 & 2 & 3 \\ 0 & 1 & 0 \\ 0 & 0 & 1
        \end{pmatrix} \\
        &= \begin{pmatrix}
            1 & 0 & 0 \\ 2 & 1 & 0 \\ 4 & 0 & 1
        \end{pmatrix} \begin{pmatrix}
            1 & 0 & 0 \\ 0 & 1 & -6 \\ 0 & 3 & -6
        \end{pmatrix} \begin{pmatrix}
            1 & 2 & 3 \\ 0 & 1 & 0 \\ 0 & 1 & 1
        \end{pmatrix} \\
        &= \begin{pmatrix}
            1 & 0 & 0 \\ 2 & 1 & 0 \\ 4 & 0 & 1
        \end{pmatrix} \begin{pmatrix}
            1 & 0 & 0 \\ 0 & 1 & 0 \\ 0 & 3 & 12
        \end{pmatrix} \begin{pmatrix}
            1 & 2 & 3 \\ 0 & -5 & -6 \\ 0 & 1 & 1
        \end{pmatrix} \\
        &= \begin{pmatrix}
            1 & 0 & 0 \\ 2 & 1 & 0 \\ 4 & 3 & 1
        \end{pmatrix} \begin{pmatrix}
            1 & 0 & 0 \\ 0 & 1 & 0 \\ 0 & 0 & 12
        \end{pmatrix} \begin{pmatrix}
            1 & 2 & 3 \\ 0 & -5 & -6 \\ 0 & 1 & 1
        \end{pmatrix}.
    \end{align*}
\end{example}

\begin{theorem}
    Every matrix has a unique SNF.
\end{theorem}

\begin{theorem}
    Any subgroup of $\Z^n$ is finitely generated and free.
\end{theorem}
\begin{proof}
    We view $\Z^n$ as a subgroup of $\Q^n$ and let $M\le \Z^n\le \Q^n$. Let $V=\mathrm{span}(M)\le \Q^n$. As $\Q^n$ is a vector space, we can find a $\Q$-basis
    \begin{align*}
        m_1,m_2,\dots,m_k\in M
    \end{align*}
    of $V$. Let
    \begin{align*}
        M' = \Z m_1 + \Z m_2 + \cdots + \Z m_k.
    \end{align*}
    It is clear that $M'\le M$. Consider a $k$-dimensional closed box (or a cell of a lattice) having $m_1,\dots,m_k$ as edges. As $m_1,\dots,m_k$ are vectors of finite length, the box contains finitely many integer points. Hence $[M:M']<\infty$.

    As $M'$ is finitely generated and $[M:M']<\infty$, $M$ is also finitely generated. Let
    \begin{align*}
        e_1,\dots,e_m
    \end{align*}
    be the generators of $M$. Then we know that
    \begin{align*}
        M=\{n_1e_1+\cdots+n_me_m:n_1,\dots,n_m\in \Z\},
    \end{align*}
    and as $M\le \Z^n$,
    \begin{align*}
        n_1e_1+\cdots+n_me_m=n_1'e_1+\cdots+n_m'e_m\iff (n_1,\dots,n_m)=(n_1',\dots,n_m').
    \end{align*}
    Therefore $M\simeq \Z^m$ and $M$ is free.
\end{proof}

\begin{example}
    $\Q$ is not a free abelian group. If $\Q$ is free and if $\Q\simeq \bigoplus_{i\in I} \Z$ for some $I$, then there exists an isomorphism $f$ mapping $\Q$ to $\bigoplus_{i\in I} \Z$ and
    \begin{align*}
        f(1) = n_{i_1} e_{i_1}+\cdots+n_{i_k}e_{i_k}
    \end{align*}
    for some $i_1,\dots,i_k\in I$ and $n_{i_1},\dots,n_{i_k}\in \Z$. For all $n\in \N$, $n\cdot(1/n)=1$, but there does not exist $x\in \bigoplus_{i\in I} \Z$ such that
    \begin{align*}
        nx = n_{i_1} e_{i_1}+\cdots+n_{i_k}e_{i_k}.
    \end{align*}
    This contradicts with the fact that $f$ is an isomorphism.
\end{example}

\begin{theorem}\label{Theorem 1.18.8.}
    A finitely generated abelian group $A$ is isomorphic to
    \begin{align*}
        \Z^r \oplus \Z/d_1\Z \oplus \cdots\oplus \Z/d_n\Z
    \end{align*}
    for some $r\ge 0$ and integers $d_i\ge 2$ satisfying $d_i|d_{i+1}$ for $i=1,2,\dots,n-1$.
\end{theorem}
\begin{proof}
    Let $\{a_1,\dots,a_N\}$ be the generators of $A$. Then there exists a homomorphism
    $\pi:\Z^N\rightarrow A$ such that
    \begin{align*}
        \pi(n_1,\dots,n_N) = n_1a_1 + \cdots+n_Na_N.
    \end{align*}
    As $\ker \pi$ is a subgroup of $\Z^N$, it is finitely generated and free. Let $M$ be a matrix having basis vectors of $\ker\pi$ as its columns. Consider a Smith decomposition $M=XDY$ of $M$. Here, $X$ and $Y$ can be considered as change of bases, so $\ker\pi$ is isomorphic to a subspace $S$ of $\Z^N$ having columns of $D$ as a basis. We may write
    \begin{align*}
        D = \mathrm{diag}(1,1,\dots,1,d_1,d_2,\dots,d_n,\underbrace{0,0,\dots,0}_r)
    \end{align*}
    where $r\ge 0$ and $d_i\ge 2$, $d_i|d_{i+1}$ for $i=1,2,\dots,n-1$. Then
    \begin{align*}
        A\simeq \Z^N/D\simeq \Z^N/S \simeq \Z^r \oplus \Z/d_1\Z \oplus \cdots\oplus \Z/d_n\Z
    \end{align*}
    and we obtain the desired result.
\end{proof}

\begin{theorem}
    Let $d_i$, $d_i'\ge 2$ and $d_i|d_{i+1}$, $d_i'|d_{i+1}'$ for $i=1,2,\dots$. Then
    \begin{align*}
        \Z^r\oplus\Z/d_1\Z \oplus \cdots\oplus \Z/d_n\Z\simeq\Z^{r'}\oplus\Z/{d_1'}\Z \oplus \cdots\oplus \Z/d_n'\Z
    \end{align*}
    if and only if $r=r'$, $n=n'$ and $(d_1,\dots,d_n)=(d_1',\dots,d_n')$.
\end{theorem}
\begin{proof}
    If the factors are the same, the isomorphism is clear. We prove the opposite direction. First, consider the torsion subgroup
    \begin{align*}
        M_\mathrm{tor}=\{m\in M:am=0\mathrm{~for~some~}a\in \Z\}.
    \end{align*}
    As the torsion subgroups of each groups are $\Z^r$ and $\Z^{r'}$, we obtain $r=r'$.

    Now we reduce the problem to the case
    \begin{align*}
        \Z/d_1\Z \oplus \cdots\oplus \Z/d_n\Z\simeq\Z/{d_1'}\Z \oplus \cdots\oplus \Z/d_n'\Z
    \end{align*}
    To prove $(d_1,\dots,d_n)=(d_1',\dots,d_n')$, we use induction. Assume $d_i=d_i'$ for $i<k$. Without loss of generality, let $d_{k}\le d_{k}'$. If $d_k\neq d_k'$, there exists an element in the left group having order $d_k$, but there does not exist such an element in the right group. Therefore $d_k=d_k'$.
\end{proof}


\section{Torsion abelian groups}
\begin{definition}
    Let $M$ be a group. A \emph{torsion abelian group} is a group defined as
    \begin{align*}
        M_\mathrm{tor}=\{m\in M:am=0\mathrm{~for~some~}a\in \Z\}.
    \end{align*}
\end{definition}

\begin{proposition}
    $\Q/\Z$ is not finitely generated.
\end{proposition}
\begin{proof}
    As $\Q/\Z$ contains $\frac{1}{n}\Z$ for all $n\in\N$, it is infinite. However, every finitely generated subgroup of $\Q/\Z$ is finite: if
    \begin{align*}
        \frac{p_1}{q_1},\dots,\frac{p_r}{q_r}
    \end{align*}
    are the generators and if $L$ is the least common multiple of $q_1,\dots,q_r$, then the finitely generated subgroup is the subgroup of the finite group $\langle L^{-1}\rangle$. Therefore $\Q/\Z$ is not finitely generated.
\end{proof}
\begin{remark}
    Note that $\Q/\Z$ is a torsion of itself. This is an example of a torsion abelian group which is not finitely generated.
\end{remark}

\begin{definition}
    Let $A$ be a torsion abelian group and $p$ be a prime. An element $a\in A$ is called \emph{$p$-primary} if $p^k a = 0$ for some $k\ge 1$. The set $A(p)$ of all $p$-primary elements is called the \emph{$p$-primary subgroup} of $A$.
\end{definition}
\begin{remark}
    $A(p)$ is a subgroup of $A$, as it is closed under group operation, contains an identity element, and contains every inverses of its elements.
\end{remark}

\begin{theorem}
    A torsion subgroup is the direct sum of its $p$-primary subgroups.
\end{theorem}
\begin{proof}
    It suffices to show that the family of inclusions $i_p:A(p)\rightarrow A$ satisfies the universal property of a direct sum. Let $Z$ be an abelian group and let
    \begin{align*}
        f_p : A(p)\rightarrow Z
    \end{align*}
    be homomorphisms. Let $\phi:A\rightarrow Z$ be a homomorphism satisfying
    \begin{align*}
        f_p = \phi\circ i_p
    \end{align*}
    for all prime $p$. For $a\in A$, let $n$ be the smallest positive integer satisfying $na=0$ and write
    \begin{align*}
        n=p_1^{e_1}\cdots p_k^{e_k}.
    \end{align*}
    Then by Chinese remainder theorem, there exists $x\in \Z$ and $m_i\equiv \delta_{ij} \mod p_j^{e_j}$ such that
    \begin{align*}
        1 = nx + m_1+\cdots+m_k,
    \end{align*}
    and from $na=0$ we obtain
    \begin{align*}
        a = m_1 a + \cdots + m_k a.
    \end{align*}
    Here, as $m_i p_i^{e^i} = n$, $m_i a \in A(p_i)$ and
    \begin{align*}
        \phi(a) = \phi(m_1a+\cdots+m_k a) = f_{p_1}(m_1 a)+\cdots+f_{p_k}(m_k a).
    \end{align*}
    Therefore such homomorphism exists and is unique, and we obtain the result.
\end{proof}

\begin{example}
    $\Q/\Z(p)$ contains $\frac{1}{p^k}\Z$ for all $k\in \N$, so it is infinite.
\end{example}

\begin{example}
    $\Q/\Z$ is isomorphic to the torsion subgroup of $\C^\times$ by the isomorphism
    \begin{align*}
        \phi(q+\Z) = e^{2\pi iq}.
    \end{align*}
\end{example}


\section{Duality for finite abelian groups}
\begin{definition}
    Let $A$ be a group. By its \emph{dual}, we mean the set of all group homomorphisms
    \begin{align*}
        \chi:A\rightarrow \C^\times.
    \end{align*}
\end{definition}

If $G$ is a torsion, then $\chi$ necessarily factors through the torsion subgroup of $\C^\times$, which is isomorphic to $\Q/\Z$. For a (torsion) abelian group $A$, put
\begin{align*}
    A^\vee = \{\chi|A\xrightarrow{\chi}\Q/\Z\}.
\end{align*}

Consider
\begin{align*}
    A\times A^\vee \rightarrow \Q/\Z
\end{align*}
which maps $(a,\chi)\mapsto \chi(a)$. Or more generally, for two finite abelian groups, we may consider a $\Z$-bilinear map
\begin{align*}
    A\times B &\rightarrow \Q/\Z \\
    (a,b) &\mapsto \phi(a,b).
\end{align*}
It gives rise to two maps $\phi_L:A\rightarrow B^\vee$ and $\phi_R : B\rightarrow A^\vee$.

\begin{definition}
    The bilinear map $\phi$ is called \emph{perfect} if both $\phi_L$ and $\phi_R$ are isomorphisms. It is called \emph{nondegenerate} if both $\phi_L$ and $\phi_R$ are injective. 
\end{definition}

\begin{theorem}
    Let $k$ be a field. For a bilinear pairing $\phi:V\times W\rightarrow k$ on a finite dimensional vector spaces $V$ and $W$, the following statements are equivalent.
    \begin{enumerate}
        \item $\phi$ is perfect.
        \item $\phi_L$ is an isomorphism.
        \item $\phi_R$ is an isomorphism.
        \item $\phi$ is nondegenerate.
    \end{enumerate}
\end{theorem}
\begin{proof}
    If $\phi$ is perfect, it is clear that the other assertions hold. Let
    \begin{align*}
        \{v_1,\dots,v_n\},\{w_1,\dots,w_m\}
    \end{align*}
    be the bases of $V$ and $W$. Consider the matrices
    \begin{align*}
        A = ((\phi_L(v_i))(w_j))_{i,j},\quad B=((\phi_R(w_i))(v_j))_{i,j}.
    \end{align*}
    $\phi_L$ is an isomorphism if and only if $A$ is square and is of full rank. Similarly, $\phi_R$ is an isomorphism if and only if $B$ is square and is of full rank. Note that
    \begin{align*}
        (\phi_L(v))(w) = \phi(v,w) = (\phi_R(w))(v),\quad v\in V,w\in W.
    \end{align*}
    Hence $A$ and $B$ has the same rank, and $\phi_L$ is an isomorphism if and only if $\phi_R$ is an isomorphism. Moreover, $\phi_L$ and $\phi_R$ being injective implies that $\phi_L$ and $\phi_R$ are isomorphisms. Therefore if $\phi$ is nondegenerate, then $\phi$ is perfect.
\end{proof}

\begin{example}
    For infinite dimensional vector spaces, the fact that $\phi$ is nondegenerate does not imply that $\phi$ is perfect. For example, let $V$ be the space of continuous compactly-supported $\R$-valued functions on $\R$ and let
    \begin{align*}
        \phi:(f,g)\mapsto \int_{-\infty}^\infty f(x)g(x)dx.
    \end{align*}
    $\phi_L$ and $\phi_R$ are injective; if $f_1,f_2\in V$  and $f_1(x)\neq f_2(x)$ for some $x\in \R$, then consider a sequence of $g$ converging to a Dirac-delta function $\delta_x$. We can observe that the value of $\phi$ converges to different values $f_1(x)$ and $f_2(x)$. However, $\phi_L$ and $\phi_R$ are not isomorphisms, as there are no $f\in V$ satisfying
    \begin{align*}
        (\phi_L(f))(g) = g(0).
    \end{align*}
    Therefore $\phi$ is not perfect.
\end{example}

\begin{theorem}
    Let $A$ be a finite abelian group. Then the pairing $\phi:A\times A^\vee\rightarrow \Q/\Z$ is perfect.
\end{theorem}
\begin{proof}
    First, let $A_1$ and $A_2$ be abelian groups and consider the mapping
    \begin{align*}
        \phi:A_1^\vee \times A_2^\vee &\rightarrow (A_1\oplus A_2)^\vee 
    \end{align*}
    where for all $a_1\in A_1$ and $a_2\in A_2$,
    \begin{align*}
        (\phi(\chi_1,\chi_2))(a_1+a_2)=\chi_1(a_1) \chi_2(a_2),\quad \chi_1\in A_1^\vee,\chi_2\in A_2^\vee.
    \end{align*}
    For all $\chi_1\in A_1^\vee$ and $\chi_2\in A_2^\vee$, $\phi(\chi_1,\chi_2):A_1\oplus A_2\rightarrow \Q/\Z$ is indeed a group homomorphism. Moreover, $\phi$ itself is an isomorphism as it is bijective by the universal property and for all $\chi_1,\chi_1'\in A_1^\vee$ and $\chi_2,\chi_2'\in A_2^\vee$,
    \begin{align*}
        (\phi(\chi_1\chi_1',\chi_2\chi_2'))(a_1+a_2) &= (\chi_1\chi_1')(a_1) (\chi_2\chi_2')(a_2) \\
        &= \chi_1(a_1)\chi_1'(a_1)\chi_2(a_2)\chi_2'(a_2) \\
        &= (\phi(\chi_1,\chi_2)\phi(\chi_1',\chi_2'))(a_1,a_2).
    \end{align*}
    Hence $(A_1\oplus A_2)^\vee \simeq A_1^\vee \times A_2^\vee$.

    Now we prove $|A|=|A^\vee|$. Note that $A$ is a direct sum of $\Z/d\Z$'s. By the result above, it suffices to show for the cyclic case $A\simeq \Z/d\Z$, where $d\ge 2$. Here $|\Z/d\Z|=d$, and the number of homomorphisms
    \begin{align*}
        \Z/d\Z\rightarrow \Q/\Z
    \end{align*}
    is also $d$, as we can map $1$ to $\frac{0}{d},\frac{1}{d},\dots,\frac{d-1}{d}$. Hence $|A|=|A^\vee|$. (In fact, $A\simeq A^\vee$.)

    Note that $\chi(a) = \phi(a,\chi)=(\phi_L(a))(\chi)$. If $\phi_L(a) = \phi_L(b)$, then $\chi(a) = \chi(b)$ for all $\chi\in A^\vee$, so $a=b$. This implies that $\phi_L:A\rightarrow (A^\vee)^\vee$ is injective, and as $|A|=|A^\vee|=|(A^\vee)^\vee|$, $\phi_L$ is an isomorphism.

    Conversely, $\chi(a) = \phi(a,\chi) = (\phi_R(\chi))(a)$. This holds for all $a\in A$, so $\phi_R(\chi) = \chi$ and $\phi_R$ is an identity isomorphism. Therefore we conclude that $\phi$ is perfect.
\end{proof}

\begin{theorem}
    Let $A$ and $B$ be finite abelian groups. For a bilinear pairing $\phi:A\times B\rightarrow \Q/\Z$, the following statements are equivalent.
    \begin{enumerate}
        \item $\phi$ is perfect.
        \item $\phi_L$ is an isomorphism.
        \item $\phi_R$ is an isomorphism.
        \item $\phi$ is nondegenerate.
    \end{enumerate}
\end{theorem}
\begin{proof}
    If $\phi$ is perfect, it is clear that the other assertions hold. If $\phi_L$ is an isomorphism, $A\simeq B^\vee$ and $\phi$ is perfect, Similarly, if $\phi_R$ is an isomorphism, $B\simeq A^\vee$ and $\phi$ is perfect.

    If $\phi_L$ and $\phi_R$ are both injective, we obtain that $|B|=|B^\vee|\le |A|$ and $|A|=|A^\vee|\le |B|$. Hence $|A|=|B|$, and $\phi_L$ and $\phi_R$ are both isomorphisms.
\end{proof}

\begin{example}
    For infinite abelian groups, the fact that $\phi$ is nondegenerate does not imply that $\phi$ is perfect. For example, let
    \begin{align*}
        \phi:\Z\times \Q/\Z &\rightarrow \Q/\Z \\
        (a,b) &\mapsto ab.
    \end{align*}
    Here, $(\phi_L(a))(b) = ab = (\phi_R(b))(a)$. It is clear that $\phi$ is nondegenerate. However, $\phi_L$ is not an isomorphism; it does not contain the group homomorphism
    \begin{align*}
        \Q/\Z &\rightarrow \Q/\Z \\
        b &\mapsto b/2.
    \end{align*}
\end{example}


\section{Inverse limit and completion}
The imperfect pairing
\begin{align*}
    \Z\times \Q/\Z\rightarrow \Q/\Z
\end{align*}
can be perfected by taking a `limit'. In finite terms, we have a perfect pairing
\begin{align*}
    \Z/n\Z\times \frac{1}{n}\Z/\Z \rightarrow \Q/\Z.
\end{align*}
They stack as
\begin{align*}
    \begin{matrix}
        \Z/nm\Z & \times & \frac{1}{mn}\Z/\Z & \rightarrow & \Q/\Z \\
        \downarrow & & \uparrow & & || \\
        \Z/n\Z & \times & \frac{1}{n}\Z/\Z & \rightarrow & \Q/\Z
    \end{matrix}
\end{align*}
A correct notion is
\begin{align*}
    \hat{\Z}:= \{(a_n)_{n\in \N}:a_n\in \Z/n\Z,a_{mn}\equiv a_n\mod n\ \forall m\in \N\}.
\end{align*}
which is an example of an \emph{inverse limit}. \\

Note that there exists a natural homomorphism
\begin{align*}
    \phi:\hat{\Z}\rightarrow (\Q/\Z)^\vee
\end{align*}
such that
\begin{align*}
    \phi((a_n)_{n\in \N})\left(\frac{m}{n}\right) = \frac{ma_n}{n}.
\end{align*}
As
\begin{align*}
    m\phi((a_n)_{n\in \N})\left(\frac{1}{mn}\right) = \frac{a_{mn}}{n} = \frac{a_n}{n} = \phi((a_n)_{n\in \N})\left(\frac{1}{n}\right),
\end{align*}
it is well-defined. Also, $\phi$ is a homomorphism, as
\begin{align*}
    \phi((a_n)_{n\in\N} + (b_n)_{n\in\Z})\left(\frac{1}{n}\right) = \frac{a_n+b_n}{n} = \phi((a_n)_{n\in \Z})\left(\frac{1}{n}\right) + \phi((b_n)_{n\in \Z})\left(\frac{1}{n}\right).
\end{align*}

We now generally define the inverse limit using directed systems.

\begin{definition}
    A \emph{directed system} is a poset $(I,\ge)$ such that for every pair $i,j\in I$ there exists $k\in I$ such that
    \begin{align*}
        k\ge i,\quad k\ge j.
    \end{align*}
\end{definition}

\begin{example}
    Every totally ordered set is a directed system, as $i\ge j$ or $j\ge i$ holds.
\end{example}

\begin{example}
    For any set $X$, $(2^X,\supset)$ and $(2^X,\subset)$ are directed systems; consider the union and the intersection.
\end{example}

\begin{example}
    Any poset with two distinct maximal elements is not a directed system. It is clear that the two maximal elements do not satisfy the condition.
\end{example}

\begin{definition}
    A \emph{diagram} of abelian groups over a directed set $(I,\ge)$ is a family of homomorphisms
    \begin{align*}
        \phi_{ji}:A_i\rightarrow A_j
    \end{align*}
    for each relation $i\ge j$ subject to the compatibility condition
    \begin{align*}
        \phi_{kj}\circ \phi_{ji} = \phi_{ki},\quad i\ge j\ge k.
    \end{align*}
    The maps $\phi_{ji}$ are called the \emph{transition maps}. If the transition maps are clear from the context, we often denote a diagram by $(A_i)_{i\in I}$.
\end{definition}

\begin{definition}
    The \emph{inverse limit} of a given diagram $(A_i)_i$ is an abelian group $A$ equipped with homomorphisms
    \begin{align*}
        \pi_i:A\rightarrow A_i
    \end{align*}
    which satisfies the universal property:
    \begin{figure}[H]
    \centering
    \begin{tikzcd}
        A \arrow[rr, "\pi_i"] & & A_i \\
        & B \arrow[ul, dashed, "\exists!\phi"] \arrow[ur, "q_i"'] &
    \end{tikzcd}
    \end{figure}
    \noindent
    We denote the inverse limit by
    \begin{align*}
        A = \lim_{i\in I} A_i = \lim_{\leftarrow} A_i.
    \end{align*}
\end{definition}
\begin{remark}
    An inverse limit is also known as a \emph{projective limit}. The terminology \emph{pro-finite} stands for `projective limit of finite objects'.
\end{remark}

\begin{example}
    Consider $(\Z/n\Z)_{n\in \N}$. We denote $\hat{\Z}$ as the inverse limit:
    \begin{align*}
        \hat{\Z} = \lim_{n\ge 1} \Z/n\Z.
    \end{align*}
\end{example}

\begin{example}
    We can also define the inverse limit of $\Z/n\Z$ where $n$ are the powers of a fixed prime $p$. The inverse limit
    \begin{align*}
        \Z_p = \lim_{n\ge 1} \Z/p^n\Z
    \end{align*}
    is called the \emph{pro-p completion of $\Z$}. $\Z_p$ viewed as a ring is also known as the \emph{ring of $p$-adic integers}.
\end{example}

\begin{definition}
    Let $G$ be a group. The \emph{profinite completion} of $G$ is defined as
    \begin{align*}
        \hat{G}:=\lim_{N} G/N
    \end{align*}
    where $N\trianglelefteq G$ runs over the set of all finite index normal subgroups.
\end{definition}
\begin{remark}
    It is consistent with our earlier notation $\hat{\Z}$.
\end{remark}

\begin{proposition}
    Let $G$ be a finite group. Then $G\simeq \hat{G}$.
\end{proposition}
\begin{proof}
    As $(G:1)$ is finite, $G=G/1\le \hat{G}$ up to isomorphism. For all normal subgroup $N\le G$ (of finite index), there exists a homomorphism
    \begin{align*}
        \phi : G &\rightarrow G/N \\
        g &\mapsto gN.
    \end{align*}
    Hence $\hat{G}\le G$ up to isomorphism. Therefore $G\simeq \hat{G}$.
\end{proof}

\begin{definition}
    Let $A$ be an abelian group and $p$ a prime. The \emph{$p$-adic} (or \emph{pro-}$p$) \emph{completion} of $A$ is defined to be
    \begin{align*}
        A_p := \lim_{n\ge 1} A/p^n A.
    \end{align*}
\end{definition}

\begin{proposition}
    Let $A$ be a finite abelian group. Then the natural map $A(p)\rightarrow A_p$ is an isomorphism.
\end{proposition}
\begin{proof}
    Note that if $|A|=p^km$ where $m$ and $p$ are coprime, then
    \begin{align*}
        A/pA \le A / p^2 A \le \cdots \le A / p^{k-1}A \le A / p^k A = A / p^{k+1}A = \cdots
    \end{align*}
    and
    \begin{align*}
        A / p^k A \simeq A(p).
    \end{align*}
    We can give a homomorphism $A(p)\rightarrow A/p^n A$ naturally from
    \begin{align*}
        A(p) \rightarrow A \rightarrow A/p^n A.
    \end{align*}
\end{proof}

\begin{example}
    Let $k[x]$ be the ring of polynomials in a variable $x$. The inverse limit
    \begin{align*}
        \lim_{n\ge 1} k[x] / x^n k[x]
    \end{align*}
    completes $k[x]$ $x$-adically. In fact, the elements of the inverse limit can be written as
    \begin{align*}
        \lim_{n\ge 1} k[x] / x^n k[x] &\rightarrow k[x] / x^n k[x] \\
        \sum_{i=0}^\infty a_i x^i &\mapsto \sum_{i=0}^{n-1} a_i x^i + x^nk[x].
    \end{align*}
    Therefore the inverse limit can be viewed as the ring of formal power series with variable $x$.
\end{example}


\section{Free groups}
\begin{definition}
    Let $X$ be a set. A free group on $X$ is a group $F$ equipped with a map
    \begin{align*}
        i:X\rightarrow F
    \end{align*}
    satisfying the universal property: for every map $u$ mapping $X$ to some group $G$, there exists a unique group homomorphism $f$.
    \begin{figure}[H]
    \centering
    \begin{tikzcd}
        X \arrow[rr, "u"] \arrow[dr, "i"'] & & G \\
        & F \arrow[ur, dashed, "\exists!f"'] &
    \end{tikzcd}
    \end{figure}
\end{definition}

We next discuss about the existence of free groups.

\begin{definition}
    A \emph{word} in a set $S$ refers to a string of finite length in $S$. A \emph{reduced word} in $X\sqcup X$ refers to a word in $X\sqcup X$ in which no two adjacent letters are $xx^{-1}$ or $x^{-1}x$.
\end{definition}

\begin{example}
    Let $X=\{a,b\}$. $aabba^{-1}ba^{-1}$ is reduced but $abb^{-1}ab$ is not.
\end{example}

\begin{definition}
    Let $w(X)$ be the set of all reduced words in $X\sqcup X$. Then, define an operation
    \begin{align*}
        w_1\circ w_2 = \mathrm{reduction~of~}w_1w_2
    \end{align*}
    where the reduction refers to the unique reduced word
    \begin{align*}
        w_1'w_2'
    \end{align*}
    where $w_1=w_1'v$ and $w_2=v^{-1}w_2'$.
\end{definition}

\begin{theorem}
    For any set $X$, $(w(X),\circ)$ is a group. It is a free group with respect to the injection $X\rightarrow w(X)$ mapping $x$ to $\text{`}x\text{'}$.
\end{theorem}
\begin{proof}
    Let $w_1,w_2,w_3\in w(X)$ such that $w_1\circ w_2=w_1'w_2'$, $w_2\circ w_3 = w_2''w_3'$ where
    \begin{align*}
        w_1 = w_1'u,\quad w_2 = u^{-1}w_2' = w_2''v,\quad w_3=v^{-1}w_3.
    \end{align*}
    Then 
    \begin{align*}
        (w_1\circ w_2)\circ w_3 = (w_1'w_2')\circ w_3 = (w_1'uw_2)\circ w_3 = w_1'uw_2''v
    \end{align*}
    and
    \begin{align*}
        w_1\circ (w_2\circ w_3) = w_1\circ(w_2''w_3') = w_1\circ (w_2v^{-1}w_3') = w_1'w_2'v^{-1}w_3'.
    \end{align*}
    As $uw_2''=w_2'v^{-1}$,
    \begin{align*}
        (w_1\circ w_2)\circ w_3 = w_1\circ (w_2\circ w_3)
    \end{align*}
    and $(w(X),\circ)$ is associative.

    There exists an identity element $1$, which refers to an empty string. Also, for all $w\in w(X)$, write
    \begin{align*}
        w=x_1^{e_1}x_2^{e_2}\cdots x_r^{e_r},\quad x_i\in X,e_i\in \Z
    \end{align*}
    and we can find the inverse
    \begin{align*}
        w^{-1} = x_r^{-e_r}\cdots x_1^{-e_1}.
    \end{align*}

    We now verify the universal property. Let $u:X\rightarrow G$ be a map to some group $G$ and let $i:X\rightarrow w(X)$ mapping $x$ to `$x$'. Then there exists a homomorphism $\phi:w(X)\rightarrow G$ satisfying $\phi\circ i = u$;
    \begin{align*}
        \phi(x_1^{e_1}x_2^{e_2}\cdots x_r^{e_r}) = \phi(x_1)^{e_1}\cdots \phi(x_r)^{e_r}.
    \end{align*}
\end{proof}

\begin{example}
    Free objects may not exist. For example, the notion of a free field is vacuous. Note that $\Q(x,y)$ is not a free field on $x$ and $y$. Think of a universality
    \begin{figure}[H]
    \centering
    \begin{tikzcd}
        \{x,y\} \arrow[rr, "u"] \arrow[dr, "i"'] & & G \\
        & \Q(x,y) \arrow[ur, dashed, "\exists!f"'] &
    \end{tikzcd}
    \end{figure}
    \noindent
    $u$ may map $x$ and $y$ to the same element in $G$. In the case, as $x-y$ is not invertible in $G$, there does not exist any field homomorphism $f:\Q(x,y)\rightarrow G$.
\end{example}

\begin{example}
    Let $G\le \mathrm{SL}(2,\Z)$ generated by
    \begin{align*}
        a = \begin{pmatrix}
            1 & 2 \\ 0 & 1
        \end{pmatrix},b=\begin{pmatrix}
            1 & 0 \\ 2 & 1
        \end{pmatrix}.
    \end{align*}
    Then $G$ is a free group on $\{a,b\}$. To prove, we observe how the elements of $G$ act on $\R^2$. Let
    \begin{align*}
        X_1 = \{(x,y)\in\R^2:|x|<|y|\}, \\
        X_2 = \{(x,y)\in\R^2:|x|>|y|\}.
    \end{align*}
    Then for $n\in\Z\setminus \{0\}$,
    \begin{align*}
        a^n\begin{pmatrix}
            x \\ y
        \end{pmatrix} = \begin{pmatrix}
            1 & 2n \\ 0 & 1
        \end{pmatrix} \begin{pmatrix}
            x \\ y
        \end{pmatrix} = \begin{pmatrix}
            x+2ny \\ y
        \end{pmatrix}, \\
        b^n\begin{pmatrix}
            x \\ y
        \end{pmatrix} = \begin{pmatrix}
            1 & 0 \\ 2n & 1
        \end{pmatrix} \begin{pmatrix}
            x \\ y
        \end{pmatrix} = \begin{pmatrix}
            x \\ 2nx + y
        \end{pmatrix}.
    \end{align*}
    Hence, $a^n$ maps $X_1$ to (a subset of) $X_2$ and $b^n$ maps $X_2$ to (a subset of) $X_1$. Now, observe the words. Words of the form
    \begin{align*}
        a^{n_1}b^{m_1}\cdots a^{n_r}
    \end{align*}
    maps $X_1$ to $X_2$, and words of the form
    \begin{align*}
        b^{m_1}a^{n_1}\cdots b^{n_r}
    \end{align*}
    maps $X_2$ to $X_1$, so they cannot be an identity. Also, if
    \begin{align*}
        a^{n_1}b^{m_1}\cdots a^{n_r}b^{m_r}=1,
    \end{align*}
    then
    \begin{align*}
        a^{2n_1}b^{m_1}\cdots a^{n_r}b^{m_r}a^{-n_1}=1,
    \end{align*}
    and we know that this leads to a contradiction, Similarly, words of the form
    \begin{align*}
        b^{m_1}a^{n_1}\cdots b^{n_r}a^{m_r}
    \end{align*}
    cannot be an identity. Therefore $G$ is a free group on $a,b$.
\end{example}
\begin{remark}
    This kind of argument can be generalized in a form of \emph{ping-pong lemma}.
\end{remark}

A free group is free of relations other than those forced by the axioms. When additional relations are imposed, we get a notion of a group presentation.

\begin{definition}
    Let $R\subset w(X)$ is a subset and denote by $N(R)$ the normal subgroup generated by $R$ in $F(X)$. A \emph{group presentation} is the quotient
    \begin{align*}
        \langle X|R\rangle := F(X)/N(R).
    \end{align*}
\end{definition}

\begin{example}
    If $R=\emptyset$, then we recover free groups. If $R$ consists of commutators, then $N(R)=[F(X),F(X)]$.
\end{example}

\begin{theorem}
    Group presentations satisfy the universal property: if `f(r)=1' for all $r\in R$, then $f$ factors uniquely as
    \begin{figure}[H]
    \centering
    \begin{tikzcd}
        X \arrow[rr, "f"] \arrow[dr, "i"'] & & G \\
        & \langle X|R\rangle \arrow[ur, dashed, "\exists!\phi"'] &
    \end{tikzcd}
    \end{figure}
\end{theorem}
\begin{proof}
    Note that
    \begin{align*}
        N(R) = \bigcap_{N\supset R}N
    \end{align*}
    where the intersection runs over every normal subgroups containing $R$. As $G\simeq F(X)/N$ for some normal subgroup $N\le F(X)$, the universal property holds.
\end{proof}

\begin{example}
    Let $C_n$ be a cyclic group of order $n$ and let $g$ be a generator. Then $g^n=1$ is a relation for $C_n$, and we have a map
    \begin{align*}
        \rho:\langle x|x^n\rangle &\rightarrow C_n \\
        x \mapsto g
    \end{align*}
    as $\rho$ is an isomorphism, we conclude that $\langle x|x^n\rangle$ is a group representation for $C_n$.
\end{example}

\begin{example}
    Consider $S_3$. Let
    \begin{align*}
        g=(12),\quad h=(123).
    \end{align*}
    Then $S_3$ is generated by $g$ and $h$, and we have a surjective map
    \begin{align*}
        F(g,h)\rightarrow S_3.
    \end{align*}
    We have relations
    \begin{align*}
        g^2 = 1,\quad h^3=1,\quad gh=h^2g,\quad\cdots
    \end{align*}
    In fact, we can write every words in the form $h^i g^j$ with the relation $gh=h^2 g$, as we can move $g$ in the left of $h$ to the right by the relation. Then with relations $g^2=1$ and $h^3=1$, we get the words
    \begin{align*}
        1,g,h,hg,h^2,h^2g
    \end{align*}
    which are all different. This implies that
    \begin{align*}
        \langle g,h|g^2,h^3,gh=h^2 g\rangle
    \end{align*}
    is the group representation for $S_3$.
\end{example}

\begin{definition}
    A group $G$ is \emph{finitely presented} if $\langle X|R\rangle\simeq G$ for some finite $X$ and $R$.
\end{definition}

\begin{definition}
    The \emph{word problem} for $G$ refers to finding an algorithm which gets a reduced word $w\in F(X)$ as an input and outputs whether $w=1$ in $G$ or not.
\end{definition}
\begin{remark}
    There exists a finitely presented group which has no solution for the word problem. (Novikov \& Boone)
\end{remark}


\clearpage
\section{Coproduct of groups}

In this section, we observe how to combine two groups in a `free' way.

\begin{definition}
    Let $G$ and $H$ be groups. The \emph{coproduct} $G*H$ is defined to satisfy the universal property
    \begin{figure}[H]
    \centering
    \begin{tikzcd}
        & G*H \arrow[dd, dashed, "\exists!\phi"'] & \\
        G \arrow[ur, "i"] \arrow[dr] & & H \arrow[ul, "i"'] \arrow[dl] \\
        & K &
    \end{tikzcd}
    \end{figure}
\end{definition}

\begin{example}
    We observe one of the most simple examples. Let $G=F(\{x\})$ and $H=F(\{y\})$, which are the free groups with one generator. Then $F(\{x,y\})\simeq F(\{x\})*F(\{y\})$; observe the universality
    \begin{figure}[H]
    \centering
    \begin{tikzcd}
        & F(\{x,y\}) \arrow[dd, dashed, "\exists!\phi"'] & \\
        F(\{x\}) \arrow[ur, "i"] \arrow[dr] & & F(\{y\}) \arrow[ul, "i"'] \arrow[dl] \\
        & K &
    \end{tikzcd}
    \end{figure}
    \noindent
    Here, we may write $F(\{x,y\})\simeq \Z*\Z$.
\end{example}

\begin{proposition}
    $G*H$ can be written as the set of reduced words of the form
    \begin{align*}
        x_1*x_2*\cdots*x_n,\quad x_i \in (G-\{1\})\sqcup (H-\{1\})
    \end{align*}
    where the consecutive pairs do not belong to the same group.
\end{proposition}
\begin{proof}
    First, we prove that the set of reduced words form a group. First we check for the associativity. Let
    \begin{align*}
        w_1 = x_1*\cdots * x_l,\quad w_2=y_1*\cdots * y_m,\quad w_3 = z_1*\cdots * z_n
    \end{align*}
    be the reduced words. If $(x_l,y_1)$ or $(y_m,z_1)$ belongs to different groups, it is clear that $(w_1w_2)w_3 = w_1(w_2w_3)$. Now assume $(x_l,y_1)$ and $(y_m,z_1)$ belongs to the same group. Then if $m=1$, then
    \begin{align*}
        (w_1w_2)w_3 = x_1 * \cdots * x_{l-1} * (x_ly_1z_1) * z_2*\cdots*z_n = w_1(w_2w_3).
    \end{align*}
    If $m\neq 1$, then
    \begin{align*}
        (w_1w_2)w_3 = x_1 * \cdots * x_{l-1} * (x_ly_1) * y_2 * \cdots * (y_mz_1) * z_2*\cdots*z_n = w_1(w_2w_3).
    \end{align*}
    Hence the operation is associative.
    
    Now it suffices to show that the group satisfies the universal property. Let there exists homomorphisms
    \begin{align*}
        g:G\rightarrow K,\quad h:H\rightarrow K
    \end{align*}
    to some group $K$. Then there exists a unique homomorphism $\phi:G*H\rightarrow K$ satisfying the universality conditions, which is
    \begin{align*}
        \phi(x_1*\cdots * x_n) = \phi_1(x_1)\cdots \phi_n(x_n)
    \end{align*}
    where $\phi_i=g$ if $x_i\in G$ and $\phi_i=h$ if $x_i\in H$.
\end{proof}

\begin{definition}
    Let $G_1,G_2$ be groups equipped with homomorphisms
    \begin{align*}
        f_i:H\rightarrow G_i
    \end{align*}
    from a group $H$. An \emph{amalgamated free product} (or a \emph{pushout}), denoted by $G_1*_H G_2$, satisfies the universal property:
    \begin{figure}[H]
    \centering
    \begin{tikzcd}
        & H \arrow[dl, "f_1"'] \arrow[dr, "f_2"] & \\
        G_1 \arrow[dr, "i_1"] \arrow[ddr, "\alpha_1"'] & & G_2 \arrow[dl, "i_2"'] \arrow[ddl, "\alpha_2"] \\
        & G_1 *_H G_2 \arrow[d, dashed, "\exists!\phi"] & \\
        & K &
    \end{tikzcd}
    \end{figure}
    \noindent
    if $\alpha_1\circ f_1 = \alpha_2\circ f_2$, then there exists a unique homomorphism $\phi:G_1*_H G_2\rightarrow K$ such that
    \begin{align*}
        \phi\circ i_1 = \alpha_1,\quad \phi\circ i_2 = \alpha_2.
    \end{align*}
\end{definition}
\begin{remark}
    For existence, take $G_i=\langle X_i|R_i\rangle$ and put
    \begin{align*}
        G_1*_H G_2 = \langle X_1\cup X_2 | R_1\cup R_2\cup \{h\in H:f_1(h)=f_2(h)\}\rangle.
    \end{align*}
\end{remark}

\begin{example}
    In topology, Seifert-van Kampen theorem asserts that
    \begin{align*}
        \pi_1(X\cup Y,b)\simeq \pi_1(X,b)*_{\pi_1(B,b)}\pi_1(Y,b)
    \end{align*}
    where $X,Y$ are subsets of a topological space with intersection $B=X\cap Y$, subject to certain conditions.
\end{example}

\begin{proposition}
    Let $\Gamma_n$ be the bouquet of $n$ circles, joined at $b\in\Gamma_n$. Then $\pi_1(\Gamma_n,b)$ is isomorphic to the free group on $n$ generators.
\end{proposition}
\begin{proof}
    The fundamental group of a circle is a free group on one generator. Applying Seifert-van Kampen theorem we obtain the result.
\end{proof}

\begin{theorem}[Schreier index formula]
    Let $F_n$ be a free group on $n$ generators. If $H\subset F_n$ is a subgroup of index $k$, then $H$ is free on $1+k(n-1)$ generators.
\end{theorem}
\begin{proof}
    The covering of $\Gamma_n$ corresponding to $H$ has $nk$ edges and $k$ vertices. Contracting a spanning tree, the Euler characteristic is preserved and we obtain a bouquet of circles with Euler characteristic $k-nk$. Since a bouquet has exactly one vertex, it has $nk-k+1$ circles.
\end{proof}




\chapter{Rings}
\section{Product of rings and idempotents}

We let all rings to be unital(have multiplicative identity).

\begin{definition}
    A \emph{ring homomorphism} between rings $R_1$ and $R_2$ is a function $f$ such that
    \begin{align*}
        f:(R_1,+)\rightarrow (R_2,+)
    \end{align*}
    is an abelian group homomorphism and
    \begin{align*}
        f:(R_1,\cdot)\rightarrow (R_2,\cdot)
    \end{align*}
    is a monoid homomorphism.
\end{definition}

\begin{definition}
    Let $R_1$ and $R_2$ be rings. The \emph{product} of $R_1$ and $R_2$ is a ring $R_1\times R_2$ having component-wise addition and multiplication.
\end{definition}
\begin{remark}
    $R_1\times R_2$ with projections $R_1\times R_2\rightarrow R_1$ and $R_1\times R_2\rightarrow R_2$ is universal among rings which map into $R_1$ and $R_2$: the universality comes from the universality of abelian groups.
    \begin{figure}[H]
    \centering
    \begin{tikzcd}
        & R_1\times R_2 \arrow[dl] \arrow[dr] &  \\
        R_1 & & R_2 \\ 
        & A \arrow[ul] \arrow[ur] \arrow[uu, dashed]
    \end{tikzcd}
    \end{figure}
    \noindent
    However, it does not give the other universal property for abelian groups; note that $R_1\rightarrow R_1\times R_2$, $x\mapsto (x,0)$ is not a ring homomorphism as it does not map multiplicative identity to multiplicative identity.
\end{remark}

We use idempotents to understand a product of rings.
\begin{definition}
    An \emph{idempotent} is an element $e$ with $e^2=e$.
\end{definition}
\begin{remark}
    In $R_1\times R_2$, $(1,0)$ and $(0,1)$ are idempotents.
\end{remark}

\begin{definition}
    Let $R$ be a ring and $a\in R$. Then we define
    \begin{align*}
        Ra := \{ra:a\in R\}.
    \end{align*}
\end{definition}
\begin{remark}
    $Ra$ is a left ideal.
\end{remark}

Let $R$ be a ring. Consider a decomposition
\begin{align*}
    1=e_1+e_2
\end{align*}
such that
\begin{enumerate}
    \item each $e_i$ is an idempotent,
    \item each $e_i$ is central; $e_ir=re_i$ for all $r\in R$,
    \item $e_1$ and $e_2$ are orthogonal; $e_1e_2=e_2e_1=0$.
\end{enumerate}
Then, the centrality implies that $J_i:Re_i$ is two sided ($J_i=Re_i R$) and we can consider the quotient. So we have two ring homomorphisms
\begin{align*}
    R\rightarrow R_1:=R/Re_2,\quad R\rightarrow R_2:=R/Re_1.
\end{align*}
We obtain a homomorphism $f:R\rightarrow R_1\times R_2$.

To show that $f$ is an isomorphism, it suffices to show that $f$ is a bijection. To show that $f$ is injective, note that
\begin{align*}
    \ker f = Re_1 \cap R e_2 = 1.
\end{align*}
Also, $f$ is surjective, as for all $r_1+Re_2\in R_1$ and $r_2+Re_1\in R_2$,
\begin{align*}
    f(r_1e_1+r_2e_2) = (r_1e_1+Re_2,r_2e_2+Re_1) = (r_1+Re_2,r_2+Re_1).
\end{align*}

\begin{example}
    Chinese remainder theorem asserts a ring isomorphism
    \begin{align*}
        \Z/N\Z\simeq \prod_{p|N} \Z/p^{t_p}\Z
    \end{align*}
    where $p^{t_p}\| N$. For example, when $N=30$, we can decompose $1\in \Z/30\Z$ as
    \begin{align*}
        1 = 15 + 10 + 6,\quad 15\in 15(\Z/2\Z),10\in 10(\Z/3\Z),6\in 6(\Z/5\Z).
    \end{align*}
\end{example}

\begin{example}
    Let $A\in M_n(\C)$ be a diagonalizable matrix. In terms of eigenspaces, we have a decomposition
    \begin{align*}
        \C^n = \bigoplus_{i=1}^k V_{\lambda_k}
    \end{align*}
    where $\lambda_1,\dots,\lambda_k$ denote the distinct eigenvalues of $A$. Define
    \begin{align*}
        R = \{B\in M_n(\C):AB=BA\}.
    \end{align*}
    Then we can decompose $R$ as
    \begin{align*}
        R\simeq \prod_{i=1}^k M_{\dim V_{\lambda_k}}(\C).
    \end{align*}
\end{example}

In terms of ideals, one can ask when
\begin{align*}
    R/I\rightarrow R/I_1\times\cdots R/I_k
\end{align*}
is an isomorphism, provided that the ideals are all two-sided and satisfy $I\subset I_t$ for $t=1,2,\cdots,k$. The injectivity is equivalent to
\begin{align*}
    I=I_1\cap \cdots\cap I_k,
\end{align*}
which can be easily seen by considering the kernel. The mapping is surjective if and only if
\begin{align*}
    I_t+I_s = R
\end{align*}
for all $t<s$. If $k=2$, note that the mapping maps
\begin{align*}
    (r_1+r_2) + I\mapsto (r_1 + I_1,r_2 + I_2),\quad r_2\in I_1,r_1\in I_2.
\end{align*}
If $I_t+I_s \subsetneq R$, then there exists $r\in R\setminus (I_t+t_s)$ and no element can map to $(r+I_t,r+I_s)$. The rest can be shown by induction; if it holds for $k\le i$, then note that the mapping
\begin{align*}
    R/I\rightarrow R/I_j\times R/I_{i+1}
\end{align*}
should also be surjective. \\

We end with a useful proposition which can be used when $R$ is commutative.
\begin{proposition}
    Assume that $R$ is commutative. Let $I$ and $J$ be two-sided ideals of $R$. If $I+J=R$, then $I\cap J=IJ$, where $IJ$ is the ideal generated by $xy$ with $x\in I$, $y\in J$.
\end{proposition}
\begin{proof}
    If $x\in I$ and $y\in J$, then $xy\in I$ and $xy\in J$ by the definition of ideals. Hence $IJ\subset I\cap J$. As $I+J=R$, there exists $x\in I$ and $y\in J$ such that $1=x+y$. So if $z\in I\cap J$, then $z=zx+zy\subset IJ$ as $zx,zy\subset IJ$. Therefore $I\cap J=IJ$.
\end{proof}


\clearpage
\section{Monoid rings}

\begin{definition}
    Let $R$ be a commutative ring. An \emph{$R$-module} is an abelian group $M$ with a mapping $R\times M\rightarrow M$ satisfying
    \begin{itemize}
        \item $r(m_1+m_2)=rm_1+rm_2$ for all $r\in R$, $m_1,m_2\in M$
        \item $(r_1+r_2)m=r_1m+r_2m$ for all $r_1,r_2\in R$, $m\in M$
        \item $(r_1r_2)m=r_1(r_2m)$ for all $r_1,r_2\in R$, $m\in M$
        \item $r\cdot 0 = 0$, $1\cdot m = m$, $0\cdot m = 0$ for all $r\in R$, $m\in M$.
    \end{itemize}
    The $R$-module
    \begin{align*}
        \bigoplus_{x\in X} R\cdot x
    \end{align*}
    is called a \emph{free $R$-module} on $X$.
\end{definition}

\begin{example}
    Every abelian group is a $\Z$-module.
\end{example}

\begin{example}
    Let $F$ be a field. Every vector space on $F$ is a $F$-module.
\end{example}

\begin{definition}
    Let $R$ be a commutative ring and $M$ be a monoid. A \emph{monoid ring} $R[M]$ is the free $R$-module on $M$ having formal sums
    \begin{align*}
        \sum_{m\in M} r_m m,\quad r_m\in R
    \end{align*}
    with $r_m=0$ for all but finitely many $m\in M$ as elements. The multiplication is defined as
    \begin{align*}
        (r_m m)(r_n n) = (r_mr_n)mn.
    \end{align*}
\end{definition}
\begin{remark}
    If $M=G$ is a group, we call $R[G]$ a group ring.
\end{remark}

\begin{example}
    The polynomial ring $R[x_1,\cdots,x_k]$ is $R[M_k]$ where $M_k\simeq (\Z_{\ge 0},+)^k$ is the free commutative monoid in $k$ generators.
\end{example}

The monoid ring satisfies the universal property:
\begin{figure}[H]
    \centering
    \begin{tikzcd}
        M \arrow[rr] \arrow[dr] & & A & \\
        & R[M] \arrow[ur, dashed, "\exists !\phi"] & & R \arrow[ul] \arrow[ll] 
    \end{tikzcd}
\end{figure}
\noindent
Here the ring homomorphisms are $R$-algebra homomorphisms. Combining with the universal property of free commutative monoid $\Z_{\ge 0}^X$ on a set $X$ we obtain
\begin{figure}[H]
    \centering
    \begin{tikzcd}
        & X \arrow[dl] \arrow[dr] & \\
        \Z_{\ge 0}^X \arrow[rr] \arrow[dr] & & A \\
        & R[\Z_{\ge 0}^X] \arrow[ur] &
    \end{tikzcd}
\end{figure}


\section{Localization of rings}
\begin{definition}
    Let $A$ be a commutative ring and $a,b\in A$. If $ab=0$ and $a,b\neq 0$, then $a$ and $b$ are \emph{zero divisors} of $A$.
\end{definition}

\begin{definition}
    A commutative ring with no zero-divisors is called an \emph{integral domain}.
\end{definition}
\begin{remark}
    It is often called an \emph{entire ring}.
\end{remark}

\begin{definition}
    Let $A$ be an integral domain. The field of fractions is defined as
    \begin{align*}
        \mathrm{Frac}(A)=\left\{\frac{a}{b}:a,b\in A,b\neq 0\right\}
    \end{align*}
    where two fractions $\frac{a}{b}$ and $\frac{c}{d}$ are identified when $ad=bc$.
\end{definition}
\begin{remark}
    The field of fractions satisfies the following universality
    \begin{figure}[H]
        \centering
        \begin{tikzcd}
            A \arrow[rr, "f"] \arrow[dr] & & B \\
            & \mathrm{Frac(A)} \arrow[ur, dashed] &
        \end{tikzcd}
    \end{figure}
    \noindent
    provided that $f(A-\{0\})\subset B^\times$.
\end{remark}

\begin{definition}
    Let $S\subset A$ be an arbitrary subset. A \emph{localization} of $A$ at $S$ is defined as a commutative ring $A_S$ equipped with a homomorphism
    \begin{align*}
        A\rightarrow A_S
    \end{align*}
    which has the universal property
    \begin{figure}[H]
        \centering
        \begin{tikzcd}
            A \arrow[rr, "f"] \arrow[dr] & & B \\
            & A_S \arrow[ur, dashed] &
        \end{tikzcd}
    \end{figure}
    \noindent
    provided that $f(S)\subset B^\times$.
\end{definition}

Note that the construction of $A_S$ is similar to that of $\mathrm{Frac}(A)$ except for the possible zero-divisors in $A$. Assume that $S$ is multiplicatively closed. Let
\begin{align*}
    I_S:=\{b\in A:sb=0\mathrm{~for~some~}s\in S\}.
\end{align*}
Elements in $I_S$ must be zero if elements in $S$ were invertible. That is, if $f:A\rightarrow B$ satisfies $f(S)\subset B^\times$, then $I_S\subset \ker (f)$. Since $S$ is multiplicatively closed, $I_S$ is an ideal. \\

Consider the set of pairs $A\times S$ with the equivalence relation
\begin{align*}
    (a,s)\sim (b,t) \iff at-bs\in I_s.
\end{align*}
Define the ring structure as
\begin{align*}
    (a,s)+(b,t)=(at+bs,st),\quad (a,s)(b,t)=(ab,st).
\end{align*}
Then it becomes a localization of $A$ at $S$. Alternatively, one could consider the product $A/I_S\times \overline{S}$ where $\overline{S}$ is the image of $S$ in $A/I_S$.

\begin{example}
    Let $A$ be an integral domain and let $S=A-\{0\}$. Then the localization $A_S$ is the field of fractions.
\end{example}

\begin{example}
    The mapping
    \begin{align*}
        \Z/p^2\Z\rightarrow \Z/p\Z
    \end{align*}
    cannot be a localization. If it is a localization for some $S\subset \Z/p^2\Z$, then $(p,m)$ for arbitrary $m\in S$ is nilpotent in the localization $A\times S$, as $(p,m)$ is nonzero in $A\times S$ and $(p,m)(p,m)=(0,m^2)$ is zero in $A\times S$. However, $\Z/p\Z$ is a field and does not contain any nilpotent element.
\end{example}

\begin{example}
    Let $R$ be a localization of $\C[t]/(t(t-1)^2)$. The ring has two prime ideals $(t)$ and $(t-1)$. For prime ideal $(t)$, note that the localization $(\C[t]/(t(t-1)^2))_{(t)}$ can be seen as $(\C[t]/(t(t-1)^2))\times (\C[t]/(t(t-1)^2) - (t))$ with the equivalence relation
    \begin{align*}
        (f_1,g_1)\sim (f_2,g_2)\iff f_1 g_2 - f_2g_1\in(t).
    \end{align*}
    So if $f=a_0+a_1t+a_2t^2+(t(t-1)^2)$ and $g=b_0+b_1t+b_2t^2+(t(t-1)^2)$, then
    \begin{align*}
        (f,g)\sim (a_0+(t(t-1)^2),b_0+(t(t-1)^2)).
    \end{align*}
    Therefore the localization is isomorphic to $\C$. For the prime ideal $(t-1)$, the localization $(\C[t]/(t(t-1)^2))_{(t-1)}$ can be seen as $(\C[t]/(t(t-1)^2))\times (\C[t]/(t(t-1)^2) - (t-1))$ with the equivalence relation
    \begin{align*}
        (f_1,g_1)\sim (f_2,g_2)\iff f_1 g_2 - f_2g_1\in((t-1)^2).
    \end{align*}
    So if $f=a_0+a_1(t-1)+a_2(t-1)^2+(t(t-1)^2)$ and $g=b_0+b_1(t-1)+b_2(t-1)^2+(t(t-1)^2)$,
    \begin{align*}
        (f,g)\sim (a_0+a_1(t-1)+(t(t-1)^2),b_0+(t(t-1)^2))
    \end{align*}
    where $(t-1,0)^2=0$. Hence we obtain $\C[\epsilon]$ with $\epsilon^2=0$.
\end{example}


\section{Local rings}
\begin{definition}
    Let $A$ be a commutative ring with a prime ideal $P$. Write $A_P$ as the localization at $A-P$.
\end{definition}
\begin{remark}
    Note that the localization at $P$ becomes $0$, as $0\in P$.
\end{remark}

\begin{definition}
    An ideal $M\subset A$ is called \emph{maximal} if it is maximal among proper ideals. A ring $A$ is \emph{local} if it has a unique maximal ideal.
\end{definition}

\begin{proposition}
    Let $I\subset A$ be an ideal and $\pi:A\rightarrow A/I$ be the quotient map. If $J\subset A/I$ is an ideal, then $\pi^{-1}(J)\subset A$ is an ideal. The association $J\mapsto \pi^{-1}(J)$ induces a bijection
    \begin{align*}
        \{J\subset A/I\}\simeq \{K\subset A:K\supset I\}.
    \end{align*}
\end{proposition}
\begin{proof}
    As $(J,+)$ is an abelian group, it is closed under addition, contains additive identity, and contains every inverse of its elements. As $\pi$ is a homomorphism, $\pi^{-1}(J)$ is also closed under addition, contains additive identity, and contains every inverse of its elements. Hence $(\pi^{-1}(J),+)$ is a group. Also, for all $g\in A$ and $x\in \pi^{-1}(J)$, $\pi(x)=y+A/I$ for some $y\in J$, so $\pi(gx)=\pi(g)(y+I) = \pi(g)y + I\in J$. Therefore $\pi^{-1}(J)$ is an ideal.

    Now we show that $J\mapsto \pi^{-1}(J)$ is a bijection. If $J\neq J'$ are two distinct ideals, then there exists some $x+A/I$ in only one of the ideals. Without loss of generality, let $x+A/I\in J-J'$. Then $x\in \pi^{-1}(J)-\pi^{-1}(J')$, so the mapping is injective. Conversely, for all ideals $K$ such that $K\subset A$ and $K\supset I$, $\pi^{-1}(\pi(K))=K$ and $\pi(K)$ is an ideal in $A/I$, so $\pi^{-1}$ is surjective.
\end{proof}

If $M$ is a maximal ideal of a ring $A$, we know that $A/M$ has no ideals other than the zero ideal and the unit ideal. They are distinct because $M\neq A$. The next proposition shows that $A/M$ is a field.

\begin{proposition}
    Suppose that a commutative ring $R$ has exactly two ideals. Then, $R$ is a field.
\end{proposition}
\begin{proof}
    If $R$ is not a field, choose a non-zero $a\in R-R^\times$. Then $0,aR,R$ are three distinct ideals.
\end{proof}
\begin{remark}
    The converse also holds. (The zero ring is not regarded a field.) 
\end{remark}

\begin{proposition}
    Let $P\subset A$ be a prime ideal. Then $A_P$ is a local ring whose maximal ideal is generated by the image of $P$ under $A\rightarrow A_P$.
\end{proposition}
\begin{proof}
    Let $M=PA_P$. Viewing $A_P$ as $A\times(A-P)$, let $(a,s)+M\in A_P/M$. if $a\in P$, then $(a,s)\in(1,s)P\subset M$ and $(a,s)+M=M$. Otherwise, $(s,a)\in A_P$ and
    \begin{align*}
        ((a,s)+M)((s,a)+M)=1+M.
    \end{align*}
    Hence $A_P/M$ is a field, and $M$ is a maximal ideal of $A_P$.
    \end{proof}

\begin{example}\label{Example 2.4.6.}
    $f:A\rightarrow A_P$ produces a local ring, and if $J\subset A_P$ is an ideal, then $f^{-1}(J)\subset A$ is also an ideal. However, in this case, we do not obtain a bijection of the form
    \begin{align*}
        \{J\subset A_P\}\simeq \{I\subset A:I\supset P\}.
    \end{align*}
    For counterexample, take $\Z_{(p)}$. Its ideals are of the form $p^k\Z_{(p)}$ for $k\ge 0$. On the other hand, $I\subset \Z$ with $I\subset (p)$ is not necessarily of the form $I=(p^k)$. If $q\neq p$ is a prime, then $(pq)\subset (p)$ but $(pq)\Z_{(p)}=p\Z_{(p)}$.
\end{example}

\begin{example}
    Let $A=\C[t]$ and $P=(t)$. Then $A/P\simeq\C$. As $A$ has no zero divisors,
    \begin{align*}
        A_P=\{\frac{f(x)}{g(x)}:f,g\in \C[t],g(0)\neq 0\}.
    \end{align*}
    The $t$-adic completion of $A$ at $P$ is $\C[[t]]$, which is a ring of formal power series. Non of them are isomorphic to each other.
\end{example}


\section{Localization and prime ideals}
In example \ref{Example 2.4.6.}, we observed that we cannot obtain a bijection between ideals in local rings and ideals in original ring containing the prime ideal. However, in this section, we prove that there exists a bijection between prime ideals. First, we characterize prime ideals in terms of integral domains.

\begin{proposition}
    Let $A$ be a ring. An ideal $I\subset A$ is prime if and only if $A/I$ is an integral domain.
\end{proposition}
\begin{proof}
    Let $I$ be a prime ideal. If $(a+I)(b+I)=ab+I=0$ for $a+I,b+I\in A/I$, then $ab\in I$. As $I$ is prime, we obtain $a\in I$ or $b\in I$. Hence $A/I$ has no zero divisor, and it is an integral domain.

    Conversely, if $A/I$ is an integral domain, it does not have a zero divisor. So if $a,b\notin I$, then $a+I,b+I\neq I$ and $(a+I)(b+I)=ab+I\neq I$ and $ab\notin I$. Hence $I$ is prime.
\end{proof}

Next, we prove that if $J$ is prime, then $\pi^{-1}(J)$ is also prime.

\begin{proposition}
    If $\phi:R_1\rightarrow R_2$ is a ring homomorphism and $P\subset R_2$ is a prime ideal, then $\phi^{-1}(P)$ is a prime ideal of $R_1$.
\end{proposition}
\begin{proof}
    Note that the mapping
    \begin{align*}
        R_1/\phi^{-1}(P)\rightarrow R_2/P
    \end{align*}
    is injective. As $P$ is prime in $R_2$, $R_2/P$ is an integral domain. As $R_1/\phi^{-1}(P)$ is isomorphic to a subring of $R_2/P$, it is an integral domain. Therefore $\phi^{-1}(P)$ is prime in $R_1$.
\end{proof}

Now we obtain the main result.

\clearpage
\begin{theorem}
    The map $J\mapsto f^{-1}(J)$ induces a bijection
    \begin{align*}
        \{J\subset A_S:J\mathrm{~is~a~prime~ideal~of~}A_S\}\simeq \{I\subset A:I\mathrm{~is~a~prime~ideal~with~}I\cap S=\emptyset\}.
    \end{align*}
\end{theorem}
\begin{proof}
    As we know that $f^{-1}(J)$ is prime when $J$ is prime, we have a map
    \begin{align*}
        f^*:\{J\subset A_S:J\mathrm{~is~a~prime~ideal~of~}A_S\}\simeq \{I\subset A:I\mathrm{~is~a~prime~ideal~with~}I\cap S=\emptyset\}.
    \end{align*}
    It suffices to find an inverse of $f^*$.

    If $Q\subset A$ is a prime ideal of $A$ with $Q\cap S=\emptyset$, then consider a mapping
    \begin{align*}
        A/Q\rightarrow A_S/QA_S.
    \end{align*}
    If $(a,s)(b,t)=(ab,st)\in QA_S$, $ab=stq$ for some $q\in Q$. As $Q\cap S=\emptyset$, $s,t\notin Q$ and $ab=0$. As $Q$ is prime, $a=0$ or $b=0$, and $(a,s)\in QA_S$ or $(b,t)\in QA_S$. Hence $QA_S$ is a prime ideal in $A_S$. Here we obtain a mapping $\sigma:Q\mapsto QA_S$.
    
    Let $Q\subset A$ be a prime ideal satisfying $Q\cap S=\emptyset$. Note that the mapping $A/Q\rightarrow A_S/QA_S$ is injective, as it is a localization of an integral domain. Then we obtain
    \begin{align*}
        f^*\circ \sigma(Q) = f^*(QA_S)=A_S.
    \end{align*}

    Conversely, let $P\subset A_S$ be a prime ideal. Note that $a/s\in P$ if and only if $a\in f^*(P)$, and if and only if $a/s\in \sigma\circ f(P)$. Hence $\sigma\circ f^*(P)=P$.

    Therefore we obtain that the mapping is a bijection.
\end{proof}

\begin{example}
    The ideals in $\Z_{(p)}$ are of the form $p^k\Z_{(p)}$. Among them, there is only one prime ideal $p\Z_{(p)}$. Note that there is only one prime ideal in $\Z$ containing $(p)$: it is $(p)$ itself.
\end{example}

\begin{example}
    For a ring $A$, denote $A[f^{-1}]$ as the localization of $A$ at $S=\{1,f,f^2,\dots\}$. For $N\in \N$, the prime ideals in $\Z[N^{-1}]$ are of the form
    \begin{align*}
        p\Z[N^{-1}]
    \end{align*}
    where $p$ is a prime number not dividing $N$. (Note that $(p)\cap \{1,N,N^2,\cdots\}=\emptyset$.)
\end{example}


\clearpage
\section{Factorization and Noetherian Property}
\begin{definition}
    A ring $R$ is called \emph{left Noetherian} if it satisfies the ascending chain condition; if
    \begin{align*}
        I_1\subset I_2\subset I_3\subset\cdots\subset R
    \end{align*}
    is a sequence of left ideals, then for some $N$ we have $I_N=I_{N+k}$ for all $k>0$.

    If $R$ is commutative, we simply say $R$ is \emph{Noetherian}.
\end{definition}

\begin{proposition}
    Let $R$ be Noetherian. Then every ideal of $R$ is finitely generated. Conversely, if every ideal of $R$ is finitely generated, then $R$ is Noetherian.
\end{proposition}
\begin{proof}
    Let $R$ be Noetherian. If there exists an ideal $I$ of $R$ which is not finitely generated, we can infinitely choose $f_n\in I$ satisfying $f_n\notin I-(f_1,\dots,f_{n-1})$. Then we obtain a strictly increasing sequence of ideals
    \begin{align*}
        (f_1)\subsetneq (f_1,f_2) \subsetneq (f_1,f_2,f_3)\subsetneq \cdots,
    \end{align*}
    which contradicts with the fact that $R$ is Noetherian.

    Conversely, let every ideal of $R$ is finitely generated. Let
    \begin{align*}
        I_1\subset I_2\subset\cdots\subset R
    \end{align*}
    is a sequence of ideals in $R$. Note that $I=\bigcup_n I_n$ is also an ideal of $R$. Let $I=(f_1,\dots,f_k)$ and let $N_i$ be the smallest integer such that $f_i\in I_{N_i}$. Let $N=\max\{N_1,\dots,N_k\}$. Then we obtain
    \begin{align*}
        I_N=I_{N+1}=\cdots=I.
    \end{align*}
    Therefore $R$ is Noetherian.
\end{proof}

\begin{definition}
    An element $r\in R$ in an integral domain $R$ is called \emph{reducible} if there exists $a,b\in R-R^\times$ such that $r=ab$. Otherwise, $r$ is called \emph{irreducible}. We denote the set of irreducible elements in $R$ as $\mathrm{Irr}(A)$.
\end{definition}

\begin{proposition}
    If $A$ is a Noetherian integral domain, every nonzero element $a\in A$ can be written as a product of irreducible elements. 
\end{proposition}
\begin{proof}
    Consider a chain of ideals
    \begin{align*}
        1 \subset (a) \subset A.
    \end{align*}
    If $a$ is irreducible, we can find $f,g\notin A^\times$ such that $a=fg$ and extend a chain to
    \begin{align*}
        1 \subset (a) \subset (f) \subset A.
    \end{align*}
    If there exists adjacent pair of ideals $(f)$, $(g)$ such that $f/g$ is not an irreducible, or if the generator of the last ideal is reducible, then we can continue the extension. As $A$ is Noetherian, this cannot continue infinitely, and we get a final result
    \begin{align*}
        1 \subset (a) \subset (a_1)\subset (a_2)\subset\cdots\subset (a_k)\subset R. 
    \end{align*}
    Let $a_0=a$ and let $f_i = a_{i-1}/a_i$. Then for all $i$, $f_i$ is irreducible and we get a factorization
    \begin{align*}
        a = f_1f_2\cdots f_k.
    \end{align*}
\end{proof}

Now, we observe the conditions for the factorization to be unique. First we define what does it mean that the factorization is unique.

\begin{definition}
    An integral domain $A$ is called a \emph{unique factorization domain(UFD)} if
    \begin{align*}
        a = u\prod_{f\in \mathrm{Irr}(A)/A^\times} f^{e_f(a)}
    \end{align*}
    for a unique $u\in A^\times$ and a unique set of exponents $e_f(a)\in \Z_{\ge 0}$ which are zero except for finitely many $f$'s.
\end{definition}
\begin{remark}
    $A$ is a UFD if and only if
    \begin{align*}
        A-\{0\} \simeq A^\times \times \{\mathrm{a~free~commutative~monoid}\}.
    \end{align*}
\end{remark}

\begin{definition}
    An element $p\in A$ is called \emph{prime} if
    \begin{align*}
        p|ab \implies p|a \mathrm{~or~} p|b.
    \end{align*}
\end{definition}

The next proposition gives a condition for a Noetherian integral domain to be a UFD.

\begin{proposition}\label{Proposition 2.6.7.}
    Suppose that $A$ is a Noetherian integral domain in which every irreducible is prime. Then $A$ is a UFD.
\end{proposition}
\begin{proof}
    Let
    \begin{align*}
        a = u\prod_{f\in \mathrm{Irr}(A)/A^\times} f^{e_f(a)} = u'\prod_{f\in \mathrm{Irr}(A)/A^\times} f^{e_f(a)'}.
    \end{align*}
    As every irreducible is prime, for all $f,g\in\mathrm{Irr}(A)$ we know $f\nmid g$. If there exists some $f\in\mathrm{Irr}(A)$ such that $e_f(a)<e_f(a)'$, then consider $a\mapsto af^{-e_f(a)}$. Then the first factorization maps to an element not divisible by $f$, but the second factorization maps to an element divisible by $f$. This is a contradiction, and hence $e_f(a) = e_f(a)'$ for all irreducible $f$. Then we also obtain $u=u'$. Therefore the factorization is unique and $A$ is a UFD.
\end{proof}


\clearpage
\section{Principal ideal domains, unique factorization domains, and Euclidean domains}
\begin{definition}
    A \emph{principal ideal domain (PID)} is a commutative ring in which every ideal is principal; generated by a single element.
\end{definition}

\begin{proposition}\label{Proposition 2.7.2.}
    Every PID is Noetherian.
\end{proposition}
\begin{proof}
    Let
    \begin{align*}
        I_1\subset I_2\subset I_3\subset\cdots
    \end{align*}
    be an ascending chain of ideals. Then as it is a PID, there exists an element $a$ such that
    \begin{align*}
        \bigcup_{k\in\N} I_k = (a).
    \end{align*}
    As there exists $n$ such that $a\in I_n$, the chain is finite. Therefore every PID is Noetherian.
\end{proof}

Alternatively, one can use the equivalence below.

\begin{proposition}[Restated]
    Let $A$ be a commutative ring. Then $A$ is Noetherian if and only if every ideal of $I$ is finitely generated.
\end{proposition}

\begin{proposition}\label{Proposition 2.7.4.}
    In every PID, irreducible implies prime.
\end{proposition}
\begin{proof}
    Let $p$ be irreducible. If $(p)\subset (q)$, then $p=qa$ and $q$ or $a$ should be a unit, which implies $(q)=(p)$ or $(q)$ is a PID itself. Hence $(p)$ is maximal. As every maximal ideal is a prime ideal, $(p)$ is a prime ideal. If $ab|p$, then $ab\in (p)$ and this implies $a\in (p)$ or $b\in(p)$. Therefore $a|p$ or $b|p$, and $p$ is prime.
\end{proof}

From Propositions \ref{Proposition 2.6.7.}, \ref{Proposition 2.7.2.} and \ref{Proposition 2.7.4.}, we know that every PID is a UFD. To find PID's, we define a Euclidean function as below.

\begin{definition}
    We say $\phi:A-\{0\}\rightarrow\Z_{\ge 0}$ with $\phi(0)=-\infty$ is a \emph{Euclidean function} if the following condition holds.
    \begin{enumerate}
        \item If $a,b\in A-\{0\}$ satisfies $\phi(a)\ge \phi(b)$, then there exists a $r\in A$ such that $\phi(a-rb)<\phi(b)$.
        \item If $a,b\in A-\{0\}$, then $\phi(a)\le \phi(ab)$.
    \end{enumerate}
    An integral domain with a Euclidean function is called a \emph{Euclidean domain}.
\end{definition}
\begin{remark}
    If $a-rb = 0$, then $-\infty < \phi(b)$, so $\phi(a-rb)<\phi(b)$. This justifies the definition $\phi(0)=-\infty$.
\end{remark}

Below are some examples of a Euclidean domain with Euclidean functions.

\begin{example}
    $\Z$ with absolute value is a Euclidean domain.
\end{example}

\begin{example}\label{Example 2.7.7.}
    If $k$ is a field, then $k[t]$ with degree is a Euclidean domain.
\end{example}

\begin{example}
    $\Z[\sqrt{-1}]$ with the square of the absolute value $|a+b\sqrt{-1}|^2=a^2+b^2$ is a Euclidean domain.
\end{example}

\begin{proposition}\label{Proposition 2.7.9.}
    Every Euclidean domain is a principal ideal domain(PID).
\end{proposition}
\begin{proof}
    Choose a Euclidean function $\phi$. Take an ideal $I\subset A$. Since the zero ideal is principal, we may assume $I\neq 0$. Now choose $a\in I-\{0\}$ such that
    \begin{align*}
        \phi(a) = \min_{x\in I-\{0\}} \phi(x).
    \end{align*}
    For all $b\in I$, then $\phi(a)\le \phi(b)$, and there exists $r\in A$ such that $\phi(b-ra)<\phi(a)$. As $a$ is defined to have a minimal $\phi$ value, we obtain that $b-ra=0$. This implies that $I=(a)$ and that $I$ is principal.
\end{proof}

The next proposition gives us why the second condition of the definition of Euclidean functions is useful (and that the additional condition does not make the constraint stronger).

\begin{proposition}
    Suppose that $\phi:A-\{0\}\rightarrow\Z_{\ge 0}$ satisfies the first condition; if $a,b\in A-\{0\}$ satisfies $\phi(a)\ge \phi(b)$, then there exists a $r\in A$ such that $\phi(a-rb)<\phi(b)$. Let
    \begin{align*}
        \psi(a) = \min_{x\in A-\{0\}} \phi(xa).
    \end{align*}
    Then $\psi$ is a Euclidean function.
\end{proposition}
\begin{proof}
    If $a,b\in A-\{0\}$ satisfies $\psi(a)\ge \psi(b)$, then there exists $x,y\in A-\{0\}$ such that $\psi(a)=\phi(xa)$ and $\psi(b) = \phi(yb)$. Then $\phi(yb) = \psi(b)\le \psi(a)\le \phi(ya)$, and there exists $r\in A$ such that $\phi(ya-ryb) < \phi(yb)$. As
    \begin{align*}
        \psi(a-rb) \le \phi(ya-ryb) < \phi(yb) = \psi(b),
    \end{align*}
    $\psi$ satisfies the first condition. Also, if $a,b\in A-\{0\}$, then
    \begin{align*}
        \psi(a) = \min_{x\in A-\{0\}} \phi(xa) \le \min_{x\in A-\{0\}} \phi(xab) = \psi(ab)
    \end{align*}
    and we obtain that $\psi$ satisfies the second condition.
\end{proof}

\begin{definition}
    We say $\phi:A-\{0\}\rightarrow\Z_{\ge 0}$ is a \emph{Dedekind-Hasse function} if the following condition holds.
    \begin{enumerate}
        \item If $a,b\in A-\{0\}$ satisfies $\phi(a)\ge \phi(b)$, then either $b|a$ or there exists a $r,s\in A$ such that $\phi(ra+sb)<\phi(b)$.
        \item If $a,b\in A-\{0\}$, then $\phi(a)\le \phi(ab)$.
    \end{enumerate}
\end{definition}
\begin{remark}
    An integral domain with a Dedekind-Hasse function is not always a UFD; one example is $\Z+\Z\frac{1+\sqrt{-19}}{2}$.
\end{remark}


\clearpage
\section{Hilbert basis theorem}
The Hilbert basis theorem gives us a method to find Noetherian rings.

\begin{theorem}
    If a commutative ring $A$ is Noetherian, so is $A[x]$.
\end{theorem}
\begin{proof}
    Let $I\subset A[x]$ be an ideal. Define $J_k\subset A$ as the ideal generated by the leading coefficients of all polynomials of the form
    \begin{align*}
        a_0+a_1x+\cdots+a_k x^k \in A[x]\cap I.
    \end{align*}
    then we have an ascending chain of ideals
    \begin{align*}
        J_0\subset J_1\subset J_2\subset\cdots\subset A.
    \end{align*}
    As $A$ is Noetherian, there exists some $N\in\N$ such that $J_N=J_{N+1}=\cdots$.

    Choose a large $M$ and choose $f_{ij}\in I$ with leading coefficient $a_{ij}$ for each $1\le i\le N$ and $1\le j\le M$ so that
    \begin{align*}
        J_k = \langle a_{k1},a_{k2},\dots,a_{kM}\rangle.
    \end{align*}
    Let $I'$ be the ideal generated by $f_{ij}$'s. It suffices to show that $I=I'$. It is clear that $I'\subset I$. If $I-I'\neq\emptyset$, let $f\in I-I'$ with minimal degree $d$ and the leading coefficient $a$. If $d\le N$, then $a\in I_d$ and we can write
    \begin{align*}
        a=b_{d1}a_{d1}+\cdots+b_{dM}a_{dM},\quad b_{dj}\in A.
    \end{align*}
    Then
    \begin{align*}
        f' = f-\sum_{j=1}^M b_{dj}x^{d-\deg f_{dj}}f_{dj}\in I-I',
    \end{align*}
    which contradicts with the minimality of $f$. If $d>N$, we can write
    \begin{align*}
        a = \sum_{i,j} b_{ij}a_{ij},\quad b_{ij}\in A
    \end{align*}
    and
    \begin{align*}
        f'=f-\sum_{i,j}b_{ij}x^{d-\deg f_{ij}}f_{ij}\in I-I'
    \end{align*}
    contradicts with the minimality of $f$. Therefore $I=I'$ and $A[x]$ is Noetherian.
\end{proof}

\begin{corollary}
    Let $k$ be a field. For every $n\ge 1$, $k[x_1,\dots,x_n]$ is Noetherian.
\end{corollary}
\begin{proof}
    Note that $k[x_1,\dots,x_n]=k[x_1][x_2]\cdots[x_n]$.
\end{proof}

\begin{corollary}
    Let $k$ be a field and $n\ge 1$ be an integer. For every ideal $I\subset k[x_1,\dots,x_n]$, the quotient $k[x_1,\dots,x_n]/I$ is Noetherian.
\end{corollary}
\begin{proof}
    Note that the quotient of the Noetherian ring is Noetherian.
\end{proof}


\clearpage
\section{Gauss' lemma}

\begin{definition}
    For $f\in A[x]-\{0\}$, its \emph{content} is defined to be the gcd of its coefficients. We denote it by $c(f)$.
\end{definition}

\begin{lemma}
    For $f_1,f_2\in A[x]-\{0\}$, $c(f_1f_2)=c(f_1)c(f_2)$.
\end{lemma}
\begin{proof}
    It suffices to show for an irreducible element $p$ of $A$. Let
    \begin{align*}
        f_1(x) &= a_mx^m + a_{m-1}x^{m-1}+\cdots+a_1x+a_0, \\
        f_2(x) &= b_nx^n + b_{n-1}x^{n-1}+\cdots+b_1x+b_0.
    \end{align*}
    Let $p^{e_1}\|c(f_1)$, $p^{e_1}\| a_s$ and $p^{e_2}\|c(f_2)$, $p^{e_2}\| b_t$. It is clear that $p^{e_1+e_2}|f_1f_2$. As the degree $(s+t)$ term of $f_1f_2$ has exponent of $p$ exactly $e_1+e_2$, we conclude that $p^{e_1+e_2}\|f_1f_2$ and $c(f_1f_2)=c(f_1)c(f_2)$.
\end{proof}

\begin{proposition}[Gauss]
    For every UFD $A$, the mapping
    \begin{align*}
        \alpha:(A[x]-\{0\})/A^\times &\rightarrow ((A-\{0\})/A^\times)\times ((K[x]-\{0\})/K^\times) \\
        f &\mapsto (c(f),f)
    \end{align*}
    is an isomorphism.
\end{proposition}
\begin{proof}
    By the preceding lemma, it is clear that $\alpha$ is a homomorphism. As the kernel of a mapping is trivial, $\alpha$ is injective. As for all $(a,f)$ in the image, we may assume that $c(f)=1$ and we obtain that $\alpha(af)=(a,f)$, so the mapping is surjective.
\end{proof}

\begin{example}
    Gauss' lemma is useful when studying the multiplicative structure of polynomials. For example, if $f\in \Z[x]$ and $c(f)=1$, then $f(x)$ is irreducible in $\Z[x]$ if and only if it is irreducible in $\Q[x]$.
\end{example}

\begin{example}
    To examine the irreducibility of $x^4+x+1$ in $\Q[x]$, we only have to test on $\Z[x]$. Suppose $x^4+x+1=fg$ for some $f,g\in \Z[x]$. If one of the polynomials are linear, than it should be $(x-a)$ where $a$ divides $1$. This is impossible, and hence $f$ and $g$ should be quadratic. Using the Lagrange interpolation and the values on $-1$, $0$, and $1$, we obtain that there does not exist such $f$ and $g$.
\end{example}


\clearpage
\section{Irreducibility criterion}
\begin{theorem}[Eisenstein's criterion]
    Let $A$ be a UFD and $f(x)\in A[x]$ be a polynomial
    \begin{align*}
        f(x) = a_0 + a_1x + \cdots + a_n x^n.
    \end{align*}
    Assume that $p\in A$ is prime in $A$. If $p\|a_0$, $p|a_1,\dots,a_{n-1}$ and $p\nmid a_n$, then $f(x)$ is irreducible.
\end{theorem}

\begin{example}
    Applying Eisenstein's criterion, we see that $x^n + pax +p\in\Z[x]$ is irreducible for any $n\ge 2$ and any $a\in\Z$.
\end{example}

\begin{example}
    Let $p$ be an odd prime. Eisenstein's criterion cannot be applied to
    \begin{align*}
        f(x)=x^{p-1}+x^{p-2}+\cdots+1\in \Z[x],
    \end{align*}
    but it does to
    \begin{align*}
        f(x+1) = \frac{(x+1)^p - 1}{(x+1)-1} = x^{p-1} + {p\choose 1}x^{p-2} + \cdots + {p\choose p-2} x + {p\choose p-1}.
    \end{align*}
\end{example}

\begin{example}
    We examine a general version of the previous example. The $n$-th cyclotomic polynomial is defined recursively by
    \begin{align*}
        \Phi_n(x) := \frac{x^n-1}{\prod_{d|n,0<d<n}\Phi_d(x)}.
    \end{align*}
    We first claim that form $n=p^k$,
    \begin{align*}
        \Phi_{p^k}(x) = \Phi_p(x^{p^{k-1}}).
    \end{align*}
    It clearly holds for $k=1$, and if it holds for $k<m$, then for $k=m$ we obtain
    \begin{align*}
        \Phi_{p^m}(x) &= \frac{x^{p^m}-1}{\prod_{d|n,0<d<p^m}\Phi_d(x)} \\
        &= \frac{x^{p^m}-1}{\prod_{0<k<m}\Phi_{p^k}(x)} \\
        &= \frac{x^{p^m}-1}{\prod_{0<k<m}\Phi_{p^k}(x)} \\
        &= \frac{\left(x^{p^{m-1}}-1\right)\left(x^{p^{m-1}(p-1)}+x^{p^{m-1}(p-2)}+\cdots+x^{p^{m-1}}+1\right)}{\prod_{0<k<m}\Phi_{p^k}(x)} \\
        &= x^{p^{m-1}(p-1)}+x^{p^{m-1}(p-2)}+\cdots+x^{p^{m-1}}+1 \\
        &= \Phi_{p}(x^{p^{m-1}}).
    \end{align*}
    As $\Phi_{p}$ is irreducible, $\Phi_{p^k}$ is irreducible for all $k\in\N$. In fact, $\Phi_n$ is irreducible for all $n\in\N$; this will be discussed in the future.
\end{example}

\begin{example}
    Eisenstein's criterion can also be applied to bivariate polynomials. For example, consider
    \begin{align*}
        f(x,y)=x^2+y^2-1\in F[x,y].
    \end{align*}
    As $x-1$ is irreducible in $F[x]$, we can apply Eisenstein criterion to $F[x][y]$ about $x-1$, which implies that $f$ is irreducible.
\end{example}

In the view of combinatorics, Eisenstein's criterion can be generalized using Newton polygons.

\begin{definition}
    Let $p\in A$ be prime and $f(x)\in A[x]$. Let $U\subset \Z_{\ge 0}$ be the support of $f(x)$, meaning that
    \begin{align*}
        f(x) = \sum_{i\in U} a_ix^i,\quad a_i\neq 0.
    \end{align*}
    The \emph{Newton polygon} $\Delta (f)$ of $f(x)$ is the convex hull of $(i,v_p(a_i))$ where $p^{v_p(a_i)}\|a_i$.
\end{definition}
\begin{remark}
    When Eisenstein's criterion applies to $f(x)$ of degree $n$, its Newton polygon has an edge connecting $(0,1)$ and $(n,0)$.
\end{remark}

Using Newton polygon, we can re-prove the Eisenstein criterion. The idea is that if $f=f_1f_2$, then
\begin{align*}
    \Delta(f) = \Delta(f_1)+\Delta(f_2).
\end{align*}
The example below gives another argument we can obtain from the Newton polygon.

\begin{example}
    Suppose that $f(x)$ has an odd degree $2k+1$, and assume that $\Delta (f)$ has edge connecting $(0,2)$ and $(2k+1,0)$. If $f(x)$ is not irreducible, there exists $f_1,f_2$ such that $f(x)=f_1(x)f_2(x)$ and $\Delta (f) = \Delta(f_1) + \Delta(f_2)$. The point $(0,2)$ can be decomposed as $(0,1)+(0,1)$ or $(0,0)+(0,2)$.
    
    If it is decomposed as $(0,1)+(0,1)$, let $n = \deg f_1 < \deg f_2$. Then $\Delta(f_1)$ has an edge connecting $(0,1)$ and $(n,0)$, and hence $\Delta(f)$ contains an edge connecting $(0,2)$ and $(n,1)$. Note that $n\le k$. As the edge is below the edge connecting $(0,2)$ and $(2k+1,0)$, it is a contradiction.

    If it is decomposed as $(0,0)+(0,2)$, let the polynomial containing $(0,2)$ has degree $n<2k+1$. Then $\Delta(f)$ contains an edge connecting $(0,0)$ and $(n,0)$, which is impossible. Therefore $f$ is irreducible.
\end{example}


\chapter{Fields}

\section{Finite and algebraic extensions}

\begin{definition}
    If $F$ is a field and $E$ is another field containing as a subfield, we say that $E$ is an \emph{extension} of $F$. We often write $E/F$ when $E$ is a field extension of $F$.
\end{definition}

\begin{example}\label{Example 3.1.2.}
    Below are some examples of extensions of field.
    \begin{itemize}
        \item $\C$ is an extension of $\R$.
        \item $\C$ is an extension of $\Q$.
        \item $\R$ is an extension of $\Q$.
        \item $\Q(\sqrt{-1})$ is an extension of $\Q$.
        \item For an indeterminate $x$, $\Q(x)$ is an extension of $\Q$.
    \end{itemize}
\end{example}

\begin{definition}
    Let $E/F$ be an extension. We say $\alpha\in E$ is \emph{algebraic over} $F$ if there is a nonzero polynomial $f(x)\in F[x]$ such that $f(\alpha)=0$. We say $\alpha$ is \emph{transcendental} if it is not algebraic.
\end{definition}

\begin{proposition}
    $\alpha\in E$ is algebraic if and only if the homomorphism
    \begin{align*}
        e_\alpha:F[x]\rightarrow E
    \end{align*}
    sending $x$ to $\alpha$ has a nonzero kernel.
\end{proposition}
\begin{proof}
    If $\alpha\in E$ is algebraic, there exists $f(x)\in F[x]$ such that $f(\alpha)=0$ and $e_\alpha(f)=0$. Hence $e_\alpha$ has a nonzero kernel. Conversely, if $e_\alpha$ has a nonzero kernel, let $f$ be a nonzero element in $\ker(e_\alpha)$. Then $f(\alpha)=0$, so $\alpha$ is algebraic.
\end{proof}

\begin{definition}
    Suppose that $\alpha\in E$ is algebraic over $F$. The \emph{minimal polynomial} of $\alpha\in E$ is defined to be the monic polynomial $m(x)\in F[x]$ which generates $\ker(e_\alpha)$.
\end{definition}
\begin{remark}
    From Example \ref{Example 2.7.7.} and Proposition \ref{Proposition 2.7.9.}, we know that $\ker(e_\alpha)$ is principal. Moreover, if $m(x)$ is reducible, say $m(x)=f(x)g(x)$ for nonconstant polynomials $f(x),g(x)\in F[x]$, then it contradicts with the fact that $m(x)$ is the polynomial in $\ker(e_\alpha)$ with minimal degree. Therefore $m(x)$ is irreducible. In fact, $m(x)$ is the minimal polynomial of $\alpha$ if and only if it is irreducible and $m(\alpha)=0$.
\end{remark}

\begin{example}
    The minimal polynomial of $\sqrt{2}$ over $\Q$ is $x^2-2$. (Note that $x^2-2$ is irreducible over $\Q$ by Eisenstein's criterion.)
\end{example}

\begin{example}
    Let $a,b\in \R$. The minimal polynomial of $a+b\sqrt{-1}\in \C$ over $\R$ is
    \begin{align*}
        (x-a)^2+b^2=x^2-2ax+a^2+b^2.
    \end{align*}
\end{example}

\begin{example}
    We find a minimal polynomial of $\sqrt{2}+\sqrt{3}$ over $\Q$. Note that
    \begin{align*}
        x=\sqrt{2}+\sqrt{3} &\implies x^2 = 5+2\sqrt{6} \\
        &\implies x^2-5 = 2\sqrt{6} \\
        &\implies x^4-10x^2+25=24 \\
        &\implies x^4-10x^2+1=0.
    \end{align*}
    As $x^4-10x^2+1$ is irreducible in $\Q$, it is the minimal polynomial of $\sqrt{2}+\sqrt{3}$.

    The minimal polynomial of $\sqrt{2}+\sqrt{3}$ over $\Q(\sqrt{2})$ is different. Note that
    \begin{align*}
        x=\sqrt{2}+\sqrt{3} &\implies x-\sqrt{2}=\sqrt{3} \\
        &\implies x^2-2\sqrt{2}x+2 = 3 \\
        &\implies x^2-2\sqrt{2}x-1=0.
    \end{align*}
    As $x^2-2\sqrt{2}x-1$ does not have a root in $\Q(\sqrt{2})$, it is irreducible and is the minimal polynomial.
\end{example}

\begin{definition}
    If $E/F$ is an extension, $E$ is a vector space over $F$. Denote by $[E:F]$ the dimension of $E$ as an $F$-vector space. If $[E:F]$ is finite, we say that an extension $E/F$ is \emph{finite}.
\end{definition}

\begin{example}
    We observe whether each extension in Example \ref{Example 3.1.2.} is finite or not.
    \begin{itemize}
        \item $\C$ is a finite extension of $\R$; $[\C:\R]=2$.
        \item $\C$ is an infinite extension of $\Q$.
        \item $\R$ is an infinite extension of $\Q$.
        \item $\Q(\sqrt{-1})$ is a finite extension of $\Q$; $[\Q(\sqrt{-1}):\Q]=2$.
        \item For an indeterminate $x$, $\Q(x)$ is an infinite extension of $\Q$. The basis
        \begin{align*}
            \{1,x,x^2,x^3,\dots\}
        \end{align*}
        is infinite.
    \end{itemize}
\end{example}

\begin{proposition}\label{Proposition 3.1.11.}
    If $E/F$ and $F/k$ are extensions, then
    \begin{align*}
        [E:k]=[E:F][F:k].
    \end{align*}
\end{proposition}
\begin{proof}
    Let $\{e_i:i\in I\}$ be a basis of $E$ as a $F$-vector space and let $\{f_j:j\in J\}$ be a basis of $F$ as a $k$-vector space. It suffices to show that
    \begin{align*}
        A = \{e_if_j:i\in I,j\in J\}
    \end{align*}
    is a basis of $E$ as a $k$-vector space. First, for arbitrary $e\in E$,
    \begin{align*}
        e=\sum_{i\in I} d_ie_i
    \end{align*}
    for some $d_i\in F$, where all but finitely many $d_i$'s are zero. Here, nonzero $d_i$'s are in $F$, so
    \begin{align*}
        d_i = \sum_{j\in J}c_{ij}f_j
    \end{align*}
    for some $c_{ij}\in k$, where all but finitely many $c_{ij}$'s are zero where $i$ is fixed. Hence
    \begin{align*}
        e = \sum_{i\in I,j\in J} c_{ij}e_if_j
    \end{align*}
    for some $c_{ij}\in k$, where all but finitely many $c_{ij}$'s are zero. Here, we obtain that $A$ spans $E/k$. Also, as $\{e_i:i\in I\}$ and $\{f_j:f\in J\}$ are bases, the choice of $d_i$ and $c_{ij}$ is unique. Therefore $A$ is a basis of $E/k$. Note that
    \begin{align*}
        [E:k] = |A| = [E:F][F:k].
    \end{align*}
\end{proof}

\begin{definition}
    Let $E/F$ be an extension. For $\alpha\in E$, denote by $F(\alpha)$ the subfield of $E$ generated by $\alpha$.
\end{definition}

\begin{proposition}
    Suppose that $\alpha$ is algebraic over $F$. Let $F[\alpha]$ be the $F$-algebra generated by $\alpha$. Then $F[\alpha]=F(\alpha)$.
\end{proposition}
\begin{proof}
    It suffices to show that $F[\alpha]$ is a field. Consider the map $e_\alpha:F[x]\rightarrow F[\alpha]$ mapping $x$ to $\alpha$. It is clear that $e_\alpha$ is surjective. By the isomorphism theorems, $e_\alpha$ induces an isomorphism
    \begin{align*}
        F[x]/(m(x))\rightarrow F[\alpha]
    \end{align*}
    where $m(x)$ is the minimal polynomial of $\alpha$. Since $F[x]$ is a PID, $(m(x))$ is maximal and $F[x]/(m(x))$ is a field.
\end{proof}

\clearpage
\begin{proposition}
    Let $\alpha\in E$ be algebraic over $F$ with minimal polynomial $m(x)$. Then
    \begin{align*}
        [F(\alpha):F]=\deg m(x).
    \end{align*}
\end{proposition}
\begin{proof}
    It follows from the fact that $F[x]/(m(x))\simeq F[\alpha]=F(\alpha)$.
\end{proof}

\begin{definition}
    An extension $E/F$ is \emph{algebraic} if every $\alpha\in E$ is algebraic over $F$.
\end{definition}

\begin{proposition}
    A finite extension is algebraic.
\end{proposition}
\begin{proof}
    Let $E/F$ be a finite extension and let $\alpha\in E$. Let $[E:F]=n$. Then
    \begin{align*}
        1,\alpha,\alpha^2,\dots,\alpha^n
    \end{align*}
    is linearly dependent over $F$. Hence there exists a polynomial of degree at most $n$ having $\alpha$ as a root. Therefore $\alpha$ is algebraic over $F$.
\end{proof}
\begin{remark}
    The converse is false; take $F_n=\Q(2^{1/2^n})$. Then
    \begin{align*}
        F_0\le F_1\le F_2\le \cdots
    \end{align*}
    and $[F_n:F_{n-1}]=2$ for all $n\ge 1$. Let $E=\bigcup_{n\ge 0}F_n$. Then $E/F$ is an infinite algebraic extension of $F_0=\Q$. 
\end{remark}

\begin{example}
    Let $\alpha=\sqrt{2}+\sqrt{3}$. As the minimal polynomial of $\alpha$ in $\Q$ is $x^4-10x^2+1$, we obtain $[\Q(\alpha):\Q]=4$. The elements of $\Q(\alpha)$ can be expressed as a linear combination of $1,\alpha,\alpha^2,\alpha^3$. For example,
    \begin{align*}
        \frac{1}{1+\sqrt{2}-\sqrt{3}}= -\frac{1}{8}-\frac{9}{8}\alpha+\frac{1}{8}\alpha^2+\frac{1}{8}\alpha^3.
    \end{align*}
\end{example}

\begin{example}
    Let $F=\Z/5\Z$ and $E=F[x]/(x^2+x+1)$. Denote by $\alpha$ the residue class of $x\in F[x]$ in $E$. Note that $x^2+x+1$ is irreducible in $F[x]$. (As $x^2+x+1$ is a quadratic polynomial, it is reducible if and only if it has a root in $F$.) The minimal polynomial of $\alpha+2$ over $F$ is
    \begin{align*}
        (x-2)^2+(x-2)+1 = x^2 - 3x + 3.
    \end{align*}
\end{example}

\begin{example}
    Let $F=\Z/2\Z$ and $E=F[x]/(x^3+x+1)$. Denote by $\alpha$ the residue class of $x\in F[x]$ in $E$. Note that $x^3+x+1$ is irreducible in $F[x]$. (As $x^3+x+1$ has degree $3$, if it is reducible it should be reduced into polynomials of degree $1$ and $2$. If it has a factor a polynomial of degree $1$, it should have a root in $F$.) We can observe that $F(\alpha)=E$. For example, we can find $a,b,c\in F$ such that
    \begin{align*}
        \frac{1}{\alpha^2+1}=a+b\alpha+c\alpha^2.
    \end{align*}
    Note that
    \begin{align*}
        1 &=(1+\alpha^2)(a+b\alpha+c\alpha^2) \\
        &= (a+b)+(b+c)\alpha+(a+c)\alpha^2.
    \end{align*}
    Hence $a=b=\frac{1}{2}$ and $c=-\frac{1}{2}$.
\end{example}

\begin{example}
    Let $E=\C(t)$ and $F=\C(t^2+1)$. As the minimal polynomial of $t$ over $F$ is $x^2+1$, $[E:F]=2$. 
\end{example}


\section{Compositum of fields}

\begin{definition}
    Suppose that $E$ and $F$ are extensions of $k$ and they are contained in a common field $L$. We depict the situation by the diagram:
    \begin{figure}[H]
        \centering
        \begin{tikzcd}
            & L\arrow[dl, dash]\arrow[dr, dash] & \\
            E\arrow[dr, dash] & & F\arrow[dl, dash] \\
            & k &
        \end{tikzcd}
    \end{figure}
    \noindent
    The \emph{compositum} $EF$ of $E$ and $F$ is the smallest subfield of $L$ containing both $E$ and $F$.
\end{definition}
\begin{remark}
    Without $L$, the compositum cannot be defined.
\end{remark}
\begin{remark}
    Compositum can be defined for arbitrary families of subfields of a given field $E$. Let $\{F_i:i\in I\}$ be the family of subfields of $E$. The compositum of $F_i$'s is defined to be the smallest subfield of $E$ containing $F_i$ for every $i\in I$.
\end{remark}

\begin{definition}
    Let $E/F$ be an extension and let $\alpha_1,\dots,\alpha_n\in E$. We denote $F(\alpha_1,\dots,\alpha_n)$ by the compositum of all $F(\alpha_i)$'s. We say $E/F$ is \emph{finitely generated} if there exists $\alpha_1,\dots,\alpha_n\in E$ such that $E=F(\alpha_1,\dots,\alpha_n)$.
\end{definition}

\begin{proposition}
    Let $E=F(\alpha_1,\dots,\alpha_n)/F$ be an extension. Then
    \begin{align*}
        E = \left\{ \frac{f(\alpha_1,\dots,\alpha_n)}{g(\alpha_1,\dots,\alpha_n)} : f,g\in F[x_1,\dots,x_n],g(\alpha_1,\dots,\alpha_n)\neq 0 \right\}.
    \end{align*}
\end{proposition}
\begin{proof}
    Let $E'$ denote the right hand side of the equation. Considering $f(x_1,\dots,x_n)=x_i$ and $g(x)=1$, it is clear that $E'$ contains $\alpha_1,\dots,\alpha_n$ and that $E\subset E'$. Conversely, note that $E$ is closed under addition, multiplication and non-zero division, and we obtain $E'\subset E$. Therefore $E=E'$.
\end{proof}

\begin{example}
    The compositum $F$ of $\Q(\sqrt{2})$ and $\Q(\sqrt{3})$ in $\R$ is equal to $\Q(\sqrt{2}+\sqrt{3})$. As $F$ contains $\sqrt{2}$ and $\sqrt{3}$, it is clear that $F$ contains $\sqrt{2}+\sqrt{3}$. So $\Q(\sqrt{2}+\sqrt{3})\subset F$. Conversely, note that
    \begin{align*}
        \frac{1}{2}\left((\sqrt{2}+\sqrt{3})^3 - 9(\sqrt{2}+\sqrt{3})\right) = \sqrt{2},\quad  -\frac{1}{2}\left((\sqrt{2}+\sqrt{3})^3 - 11(\sqrt{2}+\sqrt{3})\right) = \sqrt{3}.
    \end{align*}
    Hence $\Q(\sqrt{2}+\sqrt{3})$ contains $\Q(\sqrt{2})$ and $\Q(\sqrt{3})$, and contains $F$. Therefore $F=\Q(\sqrt{2}+\sqrt{3})$.
\end{example}

\begin{example}
    Let $\zeta=\exp \left(\frac{2\pi i}{3}\right)$. Let $E=\Q(2^{1/3})$ and $F = \Q(\zeta 2^{1/3})$. Consider $E$ and $F$ as subfields of $\C$. As the minimal polynomial of $2^{1/3}$ over $\Q$ is $x^3-2$, $[E:\Q]=3$. As the minimal polynomial of $\zeta$ over $E$ is $x^2+x+1$, $[EF:E]=2$ and $[EF:\Q]=6$. Also, as the minimal polynomial of $\zeta 2^{1/3}$ over $\Q$ is $x^3-2$, so $[F:\Q]=3$ and $[EF:F]=2$. As $[E:\Q]=3$ is not divisible by $2$, there does not exist any subfield $E'\subset E$ such that $[E':\Q]=2$.
\end{example}


\section{Distinguished class of extensions}
\begin{definition}
    Let $\mathcal{E}$ be the class of all field extensions of $F$. Let $\mathcal{C}\subset \mathcal{E}$ be a class of extensions. We say $\mathcal{C}$ is distinguished if
    \begin{enumerate}
        \item $E/F\in \mathcal{C}$ and $F/k\in\mathcal{C}$ if and only if $E/k\in\mathcal{C}$ for every tower of extensions $E/F/k$,
        \item if $E/k\in \mathcal{C}$ and $F/k$ is any extension, then $EF/F\in\mathcal{C}$,
        \item if $E/k, F/k\in\mathcal{C}$ and both $E,F$ are contained in $L$, then $EF/k\in\mathcal{C}$.
    \end{enumerate}
\end{definition}

\begin{proposition}
    The class of all finite extensions is distinguished.
\end{proposition}
\begin{proof}
    First, by Proposition \ref{Proposition 3.1.11.}, if $E/F$ and $F/k$ are finite extensions, then $E/k$ is also finite.
    
    If $E/k$ is a finite extension and $F/k$ is any extension, let $e_1,\dots,e_n$ be the basis vectors of $E$ as a vector field on $k$. Then the vector space of $F$ spanned by $e_1,\dots,e_n$ includes $E$, so it includes $EF$. Hence $EF/F$ is a finite extension.

    If $E/k$ and $F/k$ are both finite extensions, let $B$ and $B'$ be bases of $E/k$ and $F/k$, respectively. Then $BB'$ is a basis of $EF/k$, so $EF/k$ is also a finite extension.
\end{proof}

\begin{proposition}
    The class of all algebraic extensions is distinguished.
\end{proposition}
\begin{proof}
    First, if $E/F$ and $F/k$ are algebraic extensions, for all $e\in E$ there exists a polynomial
    \begin{align*}
        a_n x^n + \cdots + a_1 x + a_0 \in F[x]
    \end{align*}
    having $e$ as a root. For $i=0,1,\dots,n$, $a_i\in F$ are algebraic over $k$ and we obtain that $[k(a_i):k]$ is finite. As the class of all finite extensions is distinguished, we obtain that
    \begin{align*}
        [k(a_0,\dots,a_n):k]
    \end{align*}
    is finite. Hence $[k(a_0,\dots,a_n,e):k]=n[k(a_0,\dots,a_n:k)]$ is finite. From
    \begin{align*}
        [k(a_0,\dots,a_n,e):k] = [k(a_0,\dots,a_n,e):k(e)][k(e):k],
    \end{align*}
    $[k(e):k]$ is also finite, and $e$ is algebraic over $k$. Therefore $E/k$ is an algebraic extension.

    If $E/k$ is an algebraic extension and $F/k$ is any extension, then let
    \begin{align*}
        H=\{h\in EF:h\text{~algebraic~in~}F\}.
    \end{align*}
    It is clear that $F<H<EF$. As every element in $E$ is algebraic in $k$, it is algebraic in $F$ and we obtain $E<H$. Therefore $H=EF$ and $EF/F$ is an algebraic extension.

    If $E/k$ and $F/k$ are algebraic extensions, let
    \begin{align*}
        H=\{h\in EF:h\text{~algebraic~in~}k\}.
    \end{align*}
    Then similarly as above, $H$ contains $E$ and $F$. Therefore $EF/k$ is an algebraic extension.
\end{proof}


\clearpage
\section{Algebraically closed fields}
\begin{definition}
    A field $F$ is \emph{algebraically closed} if every algebraic extensions $E/F$ satisfies $E=F$.
\end{definition}

\begin{proposition}
    A field $F$ is algebraically closed if any irreducible polynomial in $F[x]$ has degree one.
\end{proposition}
\begin{proof}
    If there exists a polynomial $f(x)\in F[x]$ with $\deg f>1$, then there exists an algebraic extension $F(\alpha)$ where $\alpha\not\in F$ is a root of $f(x)$. As $F(\alpha)\neq F$, $F$ is not algebraically closed, which is a contradiction.
\end{proof}

\begin{theorem}
    The field $\C$ of complex numbers is algebraically closed.
\end{theorem}
\begin{proof}
    If there exists an irreducible polynomial $f(x)\in\C[x]$ with $\deg f>1$, then $1/f(x)$ is a bounded entire function. This contradicts with Lioville's theorem.
\end{proof}
\begin{remark}
    Every irreducible polynomial in $R[x]$ has degree at most $2$. Let $f(x)\in \R[x]$ be irreducible. If $\deg f>1$, then it has a root in $\C$, say $a+ib$. Then the conjugate $a-ib$ is also a root of $f(x)$, and hence $f(x)$ is divisible by $(x-a)^2+b^2\in \R[x]$. Therefore $\deg f\le 2$.
\end{remark}

\begin{example}
    Let $E\subset\C$ be the set of all numbers algebraic over $\Q$. Let $E'$ be an algebraic extension of $E$. Then as the class of algebraic extensions is distinguished, $E'/\Q$ is also an algebraic extension. As every element of $E'$ is algebraic over $\Q$, $E'=E$. Therefore $E$ is algebraically closed.
\end{example}

\begin{proposition}
    Let $E$ be a maximal element in the class of all algebraic extensions. Then $E$ is algebraically closed.
\end{proposition}
\begin{proof}
    If $E'/E$ is an algebraic extension, then $E'=E$ by maximality.
\end{proof}
\begin{remark}
    Instead of the class of all algebraic extensions, if suffices to consider the set of all fields whose underlying set is $F\times\Z_{\ge 0}$.
\end{remark}

\begin{definition}
    An algebraic extension $E/F$ is called an \emph{algebraic closure} of $F$ if $E$ is algebraically closed.
\end{definition}

\begin{proposition}
    Suppose that $E/F$ is an algebraic closure and $F'/F$ is an algebraic extension. There is an embedding $\sigma:F'\rightarrow E$ such that $\sigma|_F$ is the identity. (That is, it is an embedding over $F$.)
\end{proposition}
\begin{proof}
    Form a poset of $F$-embeddings $(K,\tau)$ where $F'/K/F$ is a subextension and $\tau:K\rightarrow E$ is an embedding over $F$. We can find a maximal element using Zorn's lemma, say $(K_0,\rho)$. We claim that $K_0=F'$. Take $\alpha\in F'$. Let $m(x)$ be the minimal polynomial of $\alpha$ over $K_0$. Let $m_\rho(x)$ be the polynomial over $\rho(K_0)$ obtained by applying $\rho$ coefficient-wise to $m(x)$. If $m(x)$ is not linear, then there exists another root $\beta$, and we can extend $(K_0,\rho)$ to $(K_0(\alpha),\rho')$ by defining $\rho'(\alpha)=\beta$. This contradicts with the maximality, and hence $m(x)$ is linear. Therefore $K_0=F'$.
\end{proof}

\begin{proposition}
    Two algebraic closures of $F$ are isomorphic over $F$.
\end{proposition}
\begin{proof}
    Let $E_1,E_2$ be two algebraic closures of $F$. Since $E_1/F$ is algebraic, there is an embedding $\sigma:E_1\rightarrow E_2$ over $F$. Since $E_1$ is algebraically closed, so is $\sigma(E_1)$. Also, $E_2/\sigma(E_1)$ is algebraic as $E_2/F$ is algebraic. Therefore $\sigma(E_1)=E_2$.
\end{proof}


\section{Splitting fields}

\begin{proposition}
    Suppose that $F$ is algebraically closed. Every polynomial $f(x)\in F[x]$ is a product of linear polynomials.
\end{proposition}
\begin{proof}
    It follows from the fact that every irreducible polynomial in $F[x]$ is linear.
\end{proof}

\begin{definition}
    A polynomial \emph{splits completely} if it is a product of linear polynomials.
\end{definition}

\begin{corollary}
    Let $f(x)\in F[x]$ be a polynomial. There exists an algebraic extension $E/F$ such that $f(x)$ splits completely in $E[x]$.
\end{corollary}
\begin{proof}
    Take $E$ to be the algebraic closure of $F$.
\end{proof}

\begin{definition}
    A \emph{splitting field} of $f(x)\in F[x]$ is an extension $E/F$ such that
    \begin{align*}
        f(x) = u\prod_{i=1}^n (x-\alpha_i)
    \end{align*}
    with $\alpha_i\in E$ and $E=F(\alpha_1,\dots,\alpha_n)$.
\end{definition}

\begin{proposition}
    Every polynomial in $F[x]$ has a splitting field. Moreover, it is a finite extension of $F$.
\end{proposition}
\begin{proof}
    Let $f(x)\in F[x]$ be a polynomial. Let $E'$ be an algebraic closure of $F$. Then
    \begin{align*}
        f(x) = u\prod_{i=1}^n (x-\alpha_i)
    \end{align*}
    for some $\alpha_i\in E'$. Take $E=F(\alpha_1,\dots,\alpha_n)$. Then $E$ is a splitting field of $f(x)$. $E/F$ is finite as $F(\alpha_1,\dots,\alpha_i)/F(\alpha_1,\dots,\alpha_{i-1})$ is finite for all $i=1,\dots,n-1$.
\end{proof}

\begin{example}
    In this example, we find splitting fields of the polynomials in $\Q[x]$.
    \begin{itemize}
        \item For $f(x)=x^2-1$, $\Q$ already splits $f(x) = x^2-1=(x-1)(x+1)$. Therefore $\Q$ itself is the splitting field of $f(x)$.
        \item For $f(x) = x^2-2$, $\Q(\sqrt{2})$ is the splitting field of $f(x)$.
        \item For $f(x) = x^2+3$, $\Q(\sqrt{3}i)$ is the splitting field of $f(x)$.
        \item For $f(x) = x^3-1$, $\Q(e^{2\pi i/3})$ is the splitting field of $f(x)$.
        \item For $f(x) = x^3-2$, $\Q(e^{2\pi i/3},\sqrt[3]{2})$ is the splitting field of $f(x)$.
        \item For $f(x) = x^3+3$, $\Q(e^{\pi i/3},\sqrt[3]{3})$ is the splitting field of $f(x)$.
    \end{itemize}
\end{example}


\clearpage
\section{Normal extensions}
Let $K/k$ be an extension. We are interested in embeddings of $K/k$ into other extensions of $k$. That is, diagrams of the form
\begin{figure}[H]
    \centering
    \begin{tikzcd}
        K \arrow[rr] & & L \\
        & k \arrow[ul] \arrow[ur] &
    \end{tikzcd}
\end{figure}
\noindent
where $L/k$ is another extension.

\begin{example}
    Let $k=\Q$ and $K=Q[x]/(x^3-2)$. The diagram can be drawn as
    \begin{figure}[H]
        \centering
        \begin{tikzcd}
            \Q[x]/(x^3-2) \arrow[rr, dashed, "f"] & & \C \\
            & \Q \arrow[ul] \arrow[ur] &
        \end{tikzcd}
    \end{figure}
    \noindent
    Here, $f$ should satisfy the equation
    \begin{align*}
        f(x)^3 - 2 = 0.
    \end{align*}
    Note that
    \begin{align*}
        \Q(2^{1/3}) \neq \Q(\zeta2^{1/3})
    \end{align*}
    where $\zeta = \exp \left( \frac{2\pi i}{3} \right)$. Hence there exists two different possible embeddings.
\end{example}

\begin{theorem}
    Let $k$ be a field and $\overline{k}/k$ be an algebraic closure. Let $K/k$ be an algebraic extension contained in $\overline{k}$. The following are equivalent.
    \begin{enumerate}
        \item Every embedding of $K/k$ into $\overline{k}$ is an automorphism of $K$. That is, the embedding is an isomorphism onto $K$.
        \item $K/k$ is the splitting field (inside $\overline{k}$) of a family of polynomials.
        \item Every irreducible polynomial of $k[x]$ which has a root in $K$ splits completely in $K[x]$.
    \end{enumerate}
\end{theorem}
\begin{proof}
    Assume 1 holds. Let $\alpha\in K$ and let $f_\alpha(x)$ be the minimal polynomial of $\alpha$ over $k$. Let $\beta$ be a zero of $f_\alpha(x)$ in $\overline{k}$. There exists an isomorphism $k(\alpha)\rightarrow k(\beta)$ over $k$, mapping $\alpha$ to $\beta$. We can extend this to an embedding $\sigma:K\hookrightarrow\overline{k}$. By hypothesis, this is an automorphism of $K$, and $\sigma(\alpha)=\beta\in K$. This implies that every zero of $f_\alpha(x)$ is in $K$. Therefore $K/k$ is the splitting field of a family of polynomials
    \begin{align*}
        \{f_\alpha:\alpha\in K\}.
    \end{align*}
    One can also note that 3 holds.

    Conversely, assume 2 holds. Let $K/k$ be the splitting field of a family of polynomials $\{f_i\}_{i\in I}$. Let $\sigma:K\hookrightarrow\overline{k}$ be an embedding over $k$. If $\alpha\in K$ is a zero of some $f_i$, then $\sigma(\alpha)$ is also a zero of $f_i$. Hence $\sigma(\alpha)\in K$. As $K$ is generated by the zeros of $\{f_i\}$, $\sigma$ maps $k$ to $k$, and therefore $\sigma$ is an automorphism. Thus 1 holds.

    Now, assume 3 holds. Let $\sigma:K\hookrightarrow\overline{k}$ be an embedding over $k$. Let $\alpha\in K$ and let $f(x)$ be its minimal polynomial over $k$. As $f(x)\in k[x]$, $\sigma(\alpha)$ is also a zero of $f(x)$, and as $f(x)$ splits completely in $K[x]$ we obtain that $\sigma(\alpha)\in K$. Similarly as above, we conclude that $\sigma$ is an automorphism of $K$ and therefore 1 holds.
\end{proof}

\begin{definition}
    A field extension $K/k$ is called \emph{normal} if every irreducible polynomial of $k[x]$ which has a root in $K$ splits completely in $K[x]$.
\end{definition}
\begin{remark}
    The normality does not depend on the choice of an embedding $K\hookrightarrow \overline{k}$. Let $K/k$ be an algebraic extension and $\sigma,\tau\hookrightarrow\overline{k}$ be embeddings. Then $\sigma(K)/k$ is normal if and only if $\tau(K)/k$ is normal.
\end{remark}

\begin{example}
    As $\C/\R$ is a splitting field of $x^2+1$, it is a normal extension.
\end{example}

\begin{example}
    $\Q(\sqrt{3})/\Q$ is a splitting field of $x^2-3$, and it is a normal extension. However, $\Q(\sqrt[3]{3})/\Q$ is a splitting field of $x^3-3$, and as $x^3-3$ does not split completely in $\Q(\sqrt[3]{3})$ (note that $\sqrt[3]{3}\exp \left( \frac{2\pi i}{3} \right)\notin \Q(\sqrt[3]{3})$), it is not a normal extension. Let $\zeta_3=\exp \left( \frac{2\pi i}{3} \right)$. Then $\Q(\sqrt[3]{3}\zeta_3)/\Q(\zeta_3)$ is a splitting field of $x^3-3$, and it is a normal extension.
\end{example}

\begin{example}
    Let $\zeta_n = \exp \left( \frac{2\pi i}{n} \right)$. $\Q(\zeta_n)/\Q$ is a splitting field of $x^n-1$, and the roots of $x^n-1$ are $\zeta_n^k$, $k=0,1,\dots,n-1$. Hence it is a normal extension.
\end{example}

\begin{example}
    $\Q(\zeta_n,\sqrt[n]{2})/\Q$ is a splitting field of $x^n-2$, and the roots of $x^n-2$ are $\sqrt[n]{2}\zeta_n^k$, $k=0,1,\dots,n-1$. Hence it is a normal extension.
\end{example}

\begin{theorem}
    Normal extensions remain normal under lifting. For a tower of extensions $K/E/k$, if $K/k$ is normal, then $K/E$ is normal. If $K_1/k, K_2/k$ are normal and are contained in some field $L$, then $K_1K_2/k$ is normal, and also $K_1\cap K_2$ is normal.
\end{theorem}
\begin{proof}
    Omitted.
\end{proof}


\clearpage
\section{Separable extensions}
\begin{definition}
    Let $k$ be a field. A polynomial $f(x)\in k[x]$ is called \emph{separable} if it has no multiple roots in every algebraic extension on $k$.
\end{definition}

\begin{proposition}
    Let $k$ be a field. Suppose that $f(x)\in k[x]$ has no repeated roots in $\overline{k}$. Then it is separable.
\end{proposition}
\begin{proof}
    Let $K/k$ be an algebraic extension. Choose an embedding $\sigma:K\rightarrow \overline{k}$ over $k$. If $\alpha\in K$ is a root of $f(x)$, then so is $\sigma(\alpha)$ because $\sigma$ is over $k$. Therefore $f(x)$ has no repeated roots in $K$, and is separable.
\end{proof}

\begin{proposition}
    A polynomial $f(x)\in k[x]$ is separable if and only if $f(x)$ and $f'(x)$ are relatively prime.
\end{proposition}
\begin{proof}
    If $f(x)$ has multiple roots, say $\alpha$, then $f(x)$ and $f'(x)$ are both divisible by $x-\alpha$. Conversely, if $f(x)$ and $f'(x)$ are not relatively prime, there exists $\alpha\in\overline{k}$ such that $x-\alpha$ divides $f(x)$ and $f'(x)$. Then $f(x)$ has multiple root $\alpha$.
\end{proof}

\begin{example}
    The polynomial $f(x)=x^n-1\in\Q[x]$ is separable, as $f'(x)=nx^{n-1}$ and it is relatively prime with $f(x)$.
\end{example}

\begin{example}
    Let $p$ be a prime and let $k=\Z/p\Z$. The polynomial $f(x)=x^p-t\in k[x]$ is not separable. Indeed $f'(x)=0$, so $(f(x),f'(x))=(f(x))$ is a proper ideal of $k[x]$. In fact, as
    \begin{align*}
        x^p-t - (x-t^{1/p})^p,
    \end{align*}
    $f(x)$ has multiple roots and is not separable.
\end{example}

\begin{definition}
    An element $\alpha\in K$ is called \emph{separable over a subfield $k\subset K$} if the minimal polynomial of $\alpha$ over $k$ is separable.
\end{definition}

\begin{definition}
    An extension $K/k$ is called \emph{simple} if $K=k(\alpha)$ for some $\alpha\in K$.
\end{definition}

\begin{proposition}
    For a simple extension $K=k(\alpha)$ over $k$ of degree $n$, $\alpha$ is separable over $k$ if and only if there exists $n$ distinct embeddings $K\hookrightarrow\overline{k}$.
\end{proposition}
\begin{proof}
    Let $f(x)$ be the minimal polynomial of $\alpha$ over $k$. Then $K\simeq k[x]/(f(x))$. Each embedding $K\hookrightarrow\overline{k}$ corresponds to the roots of $f(x)$ in $\overline{k}$. The number of distinct embeddings is the number of distinct roots of $f(x)$ in $\overline{k}$.
\end{proof}

Note that the separability of a polynomial is preserved under field extension. If $f(x)\in k[x]$ and $K/k$ is an extension, then $f(x)$ is separable in $k[x]$ if and only if it is separable in $K[x]$. In fact, the separability is determined by the factorization of $f(x)$ in $\overline{k}[x]$. However, separability of an element depends on the field under consideration. For example, take $k=\Z/p\Z(t)$. Then $t^{1/p}$ is not separable over $k$, but it is separable over $k(t^{1/p})$.\\

Next, we consider the separability for non-simple extensions.

\begin{definition}
    Suppose that $K/k$ is a finite extension and $\overline{k}/k$ is an algebraic closure. We say $K/k$ is \emph{separable} if there are $[K:k]$ distinct embeddings $K\rightarrow \overline{k}$.
\end{definition}
\begin{remark}
    In fact, the primitive element theorem (Theorem \ref{Primitive Element Theorem}) states that every finite separable extension is simple. Hence this notion reduces to the original definition of separable extensions.
\end{remark}

\begin{definition}
    Let $K/k$ be a field extension. Define $[K:k]_s$ to be the number of distinct embeddings of $K$ into an algebraic closure of $k$. We call it the \emph{separable degree} of $K/k$.
\end{definition}

\begin{proposition}
    Let $F/K/k$ be a tower of extensions. Then
    \begin{align*}
        [F:K]_s[K:k]_s=[F:k]_s.
    \end{align*}
\end{proposition}
\begin{proof}
    Let $i:k\rightarrow\overline{k}$ be a trivial embedding. Let $\{\sigma_i:i\in I\}$ be the family of distinct extensions of $i$ to $K$, and for each $i$, let $\{\tau_{ij}\}$ be the family of distinct extensions of $\sigma_i$ to $F$. Note that
    \begin{align*}
        |\{\tau_{ij}\}|=[F:K]_s[K:k]_s.
    \end{align*}
    As the embedding $F\rightarrow\overline{k}$ fixing $k$ should be one of $\{\tau_{ij}\}$, we obtain the result.
\end{proof}

\begin{proposition}
    Let $K/k$ be a finite extension. Then $[K:k]_s\le [K:k]$.
\end{proposition}
\begin{proof}
    Consider the tower of simple extensions
    \begin{align*}
        K=k(\alpha_1,\dots,\alpha_n)/\cdots/ k(\alpha_1,\alpha_2)/k(\alpha_1)/k.
    \end{align*}
    It suffices to show for simple extensions. Let $f(x)$ be the minimal polynomial of $\alpha$ over $k$. Then $[K:k]$ equals the degree of $f(x)$, and $[K:k]_s$ equals the number of distinct roots of $f(x)$ in $\overline{k}$. Therefore $[K:k]_s\le [K:k]$.
\end{proof}

\begin{proposition}
    Let $F/K/k$ be a tower of extensions. The extension $F/k$ is separable if and only if both $F/K$ and $K/k$ are separable.
\end{proposition}
\begin{proof}
    If $F/K$ and $K/k$ are separable, then
    \begin{align*}
        [F:k]=[F:K][F:k]=[F:K]_s[K:k]_s=[F:k]_s,
    \end{align*}
    and $F/k$ is separable. Conversely, if $F/k$ is separable, then
    \begin{align*}
        [F:k]_s=[F:K]_s[F:k]_s\le [F:K][K:k]=[F:k]=[F:k]_s,
    \end{align*}
    and as the equality should hold, we obtain $[F:K]=[F:K]_s$ and $[K:k]=[K:k]_s$. Therefore $F/K$ and $K/k$ are separable.
\end{proof}


\clearpage
\section{Primitive element theorem}
As mentioned before, in this section, we prove the primitive element theorem, which states that every finite separable extension is simple. We first handle the case of finite fields.

\begin{proposition}
    Let $K$ be a field and $A\subset K^\times$ be a finite subgroup. Then $A$ is cyclic. In particular, $K^\times$ is cyclic if $K$ is finite.
\end{proposition}
\begin{proof}
    Let $n$ be the least common multiple of the orders of elements in $A$. Consider
    \begin{align*}
        \phi(x)=x^n-1=0.
    \end{align*}
    Note that $\phi(x)$ has at most $n$ distinct roots in any given field. Hence $A$ has at most $n$ distinct elements. On the other hand, $A$ must have an element of order exactly $n$ (think of a Chinese remainder theorem), say $g$. Thus, $A$ is generated by $g$ and has exactly $n$ elements.
\end{proof}

\begin{corollary}\label{Corollary 3.8.2.}
    If $k$ is a finite field, then any finite extension $K/k$ is simple.
\end{corollary}
\begin{proof}
    If $k$ is finite and $[K:k]=n$, then $|K|=|k|^n$ is also finite. Then $K^\times$ is generated by some $\alpha\in K^\times$. Then $K=k(\alpha)$ and $K/k$ is simple.
\end{proof}

Corollary \ref{Corollary 3.8.2.} gives the primitive element theorem for finite fields. Now we have to prove the primitive element theorem for infinite $k$. First we assume that a finite extension $K/k$ has only finitely many intermediate fields.

\begin{lemma}
    Let $K/k$ be a finite extension with only finite intermediate fields. Let $k'/k$ be a subextension of $K/k$. For $\alpha,\beta\in K$ with $K=k'(\alpha,\beta)$, then $K$ is simple over $k'$.
\end{lemma}
\begin{proof}
    Consider $k'(\alpha+c\beta)$ with $c\in k$. Since $k$ is infinite and there are only finitely many intermediate fields, there exists distinct $c_1,c_2\in k$ such that
    \begin{align*}
        k'(\alpha+c_1\beta) = k'(\alpha+c_2\beta).
    \end{align*}
    Then
    \begin{align*}
        (c_1-c_2)\beta= (\alpha+c_1\beta)-(\alpha+c_2\beta) \in  k'(\alpha+c_1\beta)
    \end{align*}
    so it contains $\beta$. Also,
    \begin{align*}
        \alpha = (\alpha+c_1\beta) - c_1\beta\in k'(\alpha+c_1\beta).
    \end{align*}
    Therefore $k'(\alpha+c_1\beta) = k'(\alpha,\beta) = K$ and $K$ is simple over $k'$.
\end{proof}

\begin{corollary}
    Let $K/k$ be a finite extension with only finite intermediate fields. Then $K/k$ is simple.
\end{corollary}
\begin{proof}
    Use induction for $K/k = k(\alpha_1,\dots,\alpha_n)/k$.
\end{proof}

Conversely, assume that $K/k$ is simple. We here show that a finite extension $K/k$ has only finitely many intermediate fields.

\begin{lemma}
    Let $K=k(\alpha)$ has degree $n$ over $k$. Then there exists at most $2^n$ intermediate extensions.
\end{lemma}
\begin{proof}
    Let $f(x)$ be the minimal polynomial of $\alpha$ over $k$. For each intermediate extension $E/k$, the minimal polynomial $g_E(x)$ of $\alpha$ over $E$ is a factor of $f(x)$. Now, let $k'$ be the intermediate extension of $K/k$ generated by the coefficients of $g_E(x)$. As $g_E(x)$ is irreducible over $E$ and $g_E(x)\in k'[x]$, $g_E(x)$ is irreducible over $k'$. Now consider the tower of extensions $K/E/k'/k$. As
    \begin{align*}
        [K:E] = \deg (g_E(x)) = [K:k'],
    \end{align*}
    we obtain $E=k'$. Hence we obtain that every intermediate extension maps injectively to a factor of $f(x)$. As $\deg f(x)=n$, it has at most $2^n$ factors. Therefore $K/k$ has at most $2^n$ intermediate extensions.
\end{proof}

Now we prove our main theorem.

\begin{theorem}\label{Primitive Element Theorem}
    Every finite separable extension is simple
\end{theorem}
\begin{proof}
    With induction, it suffices to show for the finite separable extension $K/k$ where $K=k(\alpha,\beta)$. If $k$ is finite, the theorem holds by Corollary \ref{Corollary 3.8.2.}. Assume $k$ to be infinite. Let $n=[K:k]$ and let $\sigma_1,\dots,\sigma_n$ be distinct embeddings of $K$ into $\overline{k}$ over $k$. Then the polynomial
    \begin{align*}
        P(x) = \prod_{i\neq j} (\sigma_i a+x\sigma_i\beta - \sigma_j \alpha - x\sigma_j\beta)
    \end{align*}
    is a nonzero polynomial. Since $k$ is infinite, $P(c)\neq 0$ for some $c\in k$. Consider $\gamma=\alpha+c\beta$. Then $\sigma_i(\gamma)$ are distinct, and $[k(\gamma):k]\ge n$. Therefore $k(\gamma)=k(\alpha,\beta)$ and the extension is simple.
\end{proof}

\begin{example}
    Let $k$ be a field of characteristic $p$. Consider the extension
    \begin{align*}
        k(x,y)/k(x^p,y^p)
    \end{align*}
    which is finite of degree $p^2$. For each $a\in k$, $k(x+ay)/k$ is an extension of degree $p$, since
    \begin{align*}
        (x+ay)^p = x^p+a^py^p \in k(x^p,y^p)
    \end{align*}
    but $x+ay\notin k(x^p,y^p)$. In particular, there are infinitely many intermediate field extension when $k$ is infinite.
\end{example}


\clearpage
\section{Finite fields}

For a prime $\ell$, put $\mathbb{F}_\ell:=\Z/\ell\Z$, which is a field with $\ell$ elements. Let $k$ be a finite field and let $p$ be its characteristic. As a field of characteristic zero is infinite, $p$ is nonzero. The image of the unique homomorphism
\begin{align*}
    \Z\rightarrow k
\end{align*}
is isomorphic to $\mathbb{F}_p$. Hence, $k$ is an extension of $\mathbb{F}_p$. This shows that to study finite fields, it suffices to study finite extensions of $\mathbb{F}_p$.

\begin{lemma}\label{Lemma 3.9.1.}
    Let $k$ be a finite field with $|k|=q=p^n$ for some prime $p$. If $\alpha\in k$, then $\alpha^q=\alpha$.
\end{lemma}
\begin{proof}
    By Legendre's theorem, for all $\alpha\in k^\times$ the order of $\alpha$ divides $|k^\times|=q-1$.
\end{proof}

\begin{proposition}
    Let $k$ be a finite field with $|k|=q=p^n$ for some prime $p$. Then $k$ is a splitting field of $x^q-x$ over $\mathbb{F}_p$.
\end{proposition}
\begin{proof}
    By Lemma \ref{Lemma 3.9.1.}, every element of $k$ is a zero of $x^q-x$. Also, we know that $|k|=q$. Therefore
    \begin{align*}
        x^q-x = \prod_{\alpha\in k}(x-\alpha)
    \end{align*}
    and we obtain the result.
\end{proof}

The proposition above shows that a field with $q=p^n$ elements is a splitting field of $x^q-x$ over $\mathbb{F}_p$. Now we consider its converse, and try to find a field extension $\mathbb{F}_q/\mathbb{F}_p$ of degree $n$.

\begin{proposition}
    Let $k_q$ be a splitting field of $x^q-x$ over $\mathbb{F}_p$. Then $k_q$ has $q$ elements.
\end{proposition}
\begin{proof}
    Let $k'\subset k_q$ be the set of zeros of $x^q-x$. For $a\in k'$,
    \begin{align*}
        (a^{-1})^q = a^{-q} = (a^q)^{-1} = a^{-1}
    \end{align*}
    and $a^{-1}\in k'$.
    Also, if $a,b\in k'$, then
    \begin{align*}
        (a+b)^q = (a+b)^{p^n} = a^{p^n}+b^{p^n} = a^q + b^q = a+b
    \end{align*}
    and
    \begin{align*}
        (ab)^q = a^q b^q = ab,
    \end{align*}
    so $a+b,ab\in k'$. Therefore $k'$ is a subfield of $k_q$. By the definition of a splitting field, $k'=k_q$. The separability of $x^q-x$ implies that $k'$ has exactly $q$ elements.
\end{proof}

Summarizing the results above, we obtain the following theorem.

\begin{theorem}
    Let $q=p^n$ be a power of a prime $p$. There exists a field $\mathbb{F}_q$ with $q$ element, and it is unique up to isomorphisms.
\end{theorem}

We can also take the base field having order a prime power $q=p^n$.

\begin{theorem}
    Let $p$ be a prime and let $q=p^n$. Let $\mathbb{F}_q$ be a finite field with $q$ elements, For every $m\ge 1$, there exists an extension of $\mathbb{F}_q$ with $q^m$ elements. Moreover, two such extensions are isomorphic over $\mathbb{F}_q$. 
\end{theorem}
\begin{proof}
    Note that $\mathbb{F}_q$ is a finite extension of $\mathbb{F}_p$, and there exists a finite extension $k/\mathbb{F}_p$ such that $k = q^m = p^{nm}$. $k$ is a splitting field of $x^{q^m}-x$, which has a factor $x^q-x$. Therefore $k$ has a subfield isomorphic to $\mathbb{F}_q$.
\end{proof}

\begin{lemma}
    Let $\sigma_q(a) = a^q$ for $a\in\overline{\mathbb{F}}_q$. Then $\sigma_q$ is a field automorphism of $\overline{\mathbb{F}}_q$ over $\mathbb{F}_q$.
\end{lemma}
\begin{proof}
    Note that $\sigma_q(0)=0$, $\sigma_q(1) = 1$, and for $a,b\in\overline{\mathbb{F}}_q$,
    \begin{align*}
        \sigma_q(a+b) = (a+b)^q = a^q+b^q = \sigma_q(a) + \sigma_q(b),\quad \sigma_q(ab) = (ab)^q = a^qb^q = \sigma_q(a)\sigma_q(b).
    \end{align*}
    As $\mathbb{F}_q$ is a splitting field of $x^q-x$, $\sigma_q$ fixes $\mathbb{F}_q$. Therefore $\sigma_q$ is a field automorphism of $\overline{\mathbb{F}}_q$ over $\mathbb{F}_q$.
\end{proof}

\begin{theorem}
    Consider the extension $\mathbb{F}_{q^m}/\mathbb{F}_q$. The $m$ automorphisms $1,\sigma_q,\dots,\sigma_{q}^{m-1}$ are the only automorphisms of $\mathbb{F}_{q^m}/\mathbb{F}_q$, and they are distinct.
\end{theorem}
\begin{proof}
    As $\mathbb{F}_{q^m}/\mathbb{F}_q$ is a simple extension of degree $m$, it has at most $m$ automorphisms. Hence it suffices to show that the automorphisms are distinct. Take $\alpha\in\mathbb{F}_{q^m}$ to be a multiplicative generator. Then $\sigma_q^i(\alpha)=\alpha$ if and only if $q^{i}-1$ is a multiple of $q^m-1$. Therefore the automorphisms are all distinct.
\end{proof}

We conclude that over a finite field, every finite extension is normal and separable. For every $n$, there exists a unique such extension whose automorphism group is a cyclic group of order $n$.


\clearpage
\section{Purely inseparable fields and perfect fields}
\begin{lemma}
    Let $K=k(\alpha)$ be a simple extension of $k$. Let $f(x)$ be the minimal polynomial of $\alpha$ over $k$. If the characteristic of $k$ is zero, then $f(x)$ is separable.
\end{lemma}
\begin{proof}
    If $[K:k]=1$, the case is trivial. Otherwise, $\deg f(x)\ge 1$ and $f'(x)\not\equiv 0$. As $f(x)$ is irreducible, $f(x)$ and $f'(x)$ are relatively prime, which implies that $f(x)$ is separable.
\end{proof}
\begin{remark}
    If $k$ has an inseparable extension, it should have a nonzero characteristic.
\end{remark}

Let $k$ be a field of characteristic $p$. If $K=k(\alpha)/k$ is inseparable, then the minimal polynomial $f(x)$ of $\alpha$ over $k$ should satisfy $f'(x)=0$.

\begin{lemma}
    Let $k$ be a field of characteristic $p$. Let $K=k(\alpha)/k$ be a simple extension, and let $f(x)$ be the minimal polynomial of $\alpha$ over $k$. If $f'(x) = 0$, then $f(x) = g(x^p)$ for some $g(x)\in k[x]$.
\end{lemma}
\begin{proof}
    Note that $(x^m)'=mx^{m-1}$. The polynomial $f(x)$ can only have terms of degree multiples of $p$.
\end{proof}

Inductively, define
\begin{align*}
    f(x) = g_1(x^p) = g_2(x^{p^2}) = \cdots=g_n(x^{p^n}) 
\end{align*}
which terminates when $g_n'(x) = 0$. Then $K/k$ splits as
\begin{align*}
    K(\alpha) / k(\alpha^{p^n})/k
\end{align*}
and
$k(\alpha^{p^n})/k$ is separable. Also, we obtain
\begin{align*}
    [K:k]_s = \deg g_n(x).
\end{align*}

\begin{proposition}
    If $K/k$ is a finite extension, then $[K:k]_s$ divides $[K:k]$.
\end{proposition}
\begin{proof}
    For simple extensions, define $g_1,\dots,g_n$ as above. As
    \begin{align*}
        [K:k]=\deg f(x) = p^n \cdot \deg g_n(x) = p^n \cdot [K:k]_s
    \end{align*}
    we obtain the result. For general cases, we can decompose $K/k$ into a finite number of simple extensions.
\end{proof}

\begin{definition}
    The \emph{inseparable degree} of an extension $K/k$ is defined to be
    \begin{align*}
        [K:k]_i = [K:k]/[K:k]_s.
    \end{align*}
    $K/k$ is called \emph{purely inseparable} if $[K:k]_s=1$.
\end{definition}

\begin{proposition}\label{Proposition 3.10.5.}
    Let $K/k$ be a purely inseparable extension. For every $\alpha\in K$, there exists some $n$ such that $\alpha^{p^n}\in k$.
\end{proposition}
\begin{proof}
    $k(\alpha)/k$ can be decomposed as
    \begin{align*}
        k(\alpha)/k(\alpha^{p^n})/k.
    \end{align*}
    As $k$ is purely inseparable, $k(\alpha^{p^n})=k$, and therefore $\alpha^{p^n}\in k$.
\end{proof}

\begin{definition}
    A field $k$ of characteristic $p$ is called \emph{perfect} if $k^p=k$.
\end{definition}

\begin{proposition}
    If $k$ is a perfect field, then every finite extension is separable.
\end{proposition}
\begin{proof}
    Suppose that there exists an extension $K/k$ that is not separable. Then there exists an intermediate field $S$ where $S/k$ is purely inseparable. By Proposition \ref{Proposition 3.10.5.}, for every element $\alpha\in S$ there exists some $n$ such that $\alpha^{p^n}\in k = k^{p^n}$. Hence there exists $\beta\in k$ such that $\alpha^{p^n}=\beta^{p^n}$. Then
    \begin{align*}
        \alpha^{p^n}-\beta^{p^n} = (\alpha-\beta)^{p^n} = 0
    \end{align*}
    and $\alpha=\beta$. Therefore $S=k$, which is a contradiction.
\end{proof}
\begin{remark}
    Sometimes a perfect field is defined as a field such that every finite extension is separable.
\end{remark}


\clearpage
\section{Galois extensions: definition and examples}
\begin{definition}
    A finite extension $K/k$ is called \emph{Galois} if it is normal and separable.
\end{definition}

\begin{example}
    Suppose that $k$ is a field in which $2$ is invertible (i.e. $k$ is not of characteristic $2$.) Then, every quadratic extensions of $k$ is Galois. Note that $K=k(\alpha)$ for $\alpha\in K\setminus k$. Let
    \begin{align*}
        f(x) = x^2+ax+b=(x-p)^2-q\in k[x]
    \end{align*}
    be a minimal polynomial of $\alpha$ over $k$. Then $\alpha-p\in K$, so
    \begin{align*}
        K\simeq k[x]/(x^2-q).
    \end{align*}
    As $K$ is not of characteristic $2$, it is normal and separable, and therefore invertible. In particular, $\Q(\sqrt{d})/\Q$ is a quadratic Galois extension for any square-free integer $d\neq 1$.
\end{example}

\begin{example}
    For an integer $n\ge 1$, let $\zeta_n=\exp \left( \frac{2\pi i}{n} \right)$. Note that the roots of $x^n-1$ in $\C$ are $\zeta_n^k$ for $k=0,1,\dots,n-1$. Hence $\Q(\zeta_n)$ is a splitting field of $x^n-1$ over $\Q$, and therefore $\Q(\zeta_n)/\Q$ is a Galois extension.
\end{example}

\begin{proposition}
    Suppose that $K_1/k$ and $K_2/k$ are Galois. Then the compositum $K_1K_2$ is Galois over $k$.
\end{proposition}
\begin{remark}
    In general, compositum of $K_1$ and $K_2$ are defined when there exists a field $L$ containing $K_1$ and $K_2$. If $K_1/k$ and $K_2/k$ are Galois, then the embedding $K_i\hookrightarrow\overline{k}$ has a well-defined image. Hence the composition is well-defined. More generally, a Galois extension $K/k$ corresponds to a unique subfield of $\overline{k}$ isomorphic to $K$.
\end{remark}
\begin{proof}
    As a compositum of two normal extensions is normal, $K_1K_2/k$ is a normal extension. Also, as $K_1K_2$ is generated by elements of $K_1$ and $K_2$, which are all separable over $k$, we observe that $K_1K_2/k$ is separable. Therefore $K_1K_2$ is Galois over $k$.
\end{proof}

\begin{definition}
    Let $K/k$ be a finite extension. An \emph{automorphism of $K$ over $k$} is a field automorphism
    \begin{align*}
        \sigma:K\rightarrow K
    \end{align*}
    such that $\sigma(a)=a$ for all $a\in k$.
\end{definition}

\begin{example}
    If $K=k(\alpha)$, then an automorphism of $K$ over $k$ is determined by $\sigma(\alpha)$. Moreover, $\sigma(\alpha)$ and $\alpha$ have the same minimal polynomial over $k$; they are conjugate over $k$.
\end{example}

\begin{example}
    For $\Q(\sqrt{d})/\Q$, $\sigma:\sqrt{d}\mapsto \sqrt{d}$ or $\sigma:\sqrt{d}\mapsto -\sqrt{d}$ determines an automorphism of $\Q(\sqrt{d})$ over $\Q$.
\end{example}

\begin{proposition}
    For every integer $n\ge 1$,
    \begin{align*}
        [\Q(\zeta_n):\Q]=\varphi(n)
    \end{align*}
    where $\varphi(n)$ is the Euler's totient function.
\end{proposition}
\begin{proof}
    Let $f(x)$ be the minimal polynomial of $\zeta_n$ over $\Q$. $\zeta_n$ is a zero of $x^n-1$, so $f(x)$ is a factor of $x^n-1$, and the candidates for zeros of $f(x)$ are
    \begin{align*}
        1,\zeta_n,\zeta_n^2,\dots,\zeta_n^{n-1}.
    \end{align*}
    If $\gcd (n,k)=g\neq 1$, then $\zeta_n^k = \zeta_{n/g}^{n/k}$. As $\Q(\zeta_{n/g})$ does not contain $\zeta_n$, it is an intermediate field strictly smaller than $\Q(\zeta_n)$. As $\zeta_n^k\in \Q(\zeta_{n/g})$, $\zeta_n^k$ is not a zero of $f(x)$. Otherwise, if $\gcd (n,k)=1$, consider the automorphism
    \begin{align*}
        \Q(\zeta_n) & \rightarrow \Q(\zeta_n) \\
        \zeta_n &\mapsto \zeta_n^k.
    \end{align*}
    This is an automorphism, as the order of $\zeta_n^k$ is $n$. Hence $\zeta_n^k$ is a zero of $f(x)$. Therefore, $[\Q(\zeta_n):\Q]=\deg f(x)$ equals the number of positive integers smaller than $n$ and is relatively prime to $n$, which is $\varphi(n)$.
\end{proof}

\begin{corollary}
    The minimal polynomial $\Phi_n$ of $\zeta_n$ is
    \begin{align*}
        \Phi_n(x) = \prod_{t\in (\Z/n\Z)^\times} (x-\zeta_n^t).
    \end{align*}
    Moreover, for $\Q(\zeta_n)/\Q$, $\sigma_t:\zeta_n\mapsto\zeta_n^t$ determines an automorphism of $\Q(\zeta_n)$ over $\Q$ if and only if $t$ is relatively prime to $n$.
\end{corollary}

\begin{example}
    Consider $F=\Q(\sqrt{2},\sqrt{3})$, which is Galois over $\Q$. The automorphism of $F$ over $\Q$ maps $\sqrt{2}$ to $\pm\sqrt{2}$ and maps $\sqrt{3}$ to $\pm\sqrt{3}$. The group of automorphisms is of order $4$ and is isomorphic to $\Z/2\Z\times\Z/2\Z$.
\end{example}

\begin{example}
    Let $p$ be a prime. Suppose that $k$ is a field with a primitive $p$-th root of unity $\zeta\in k$. Choose
    \begin{align*}
        c\in k^\times-(k^\times)^p
    \end{align*}
    and consider
    \begin{align*}
        K=k(c^{1/p})
    \end{align*}
    which is Galois over $k$. If $\sigma$ is an automorphism of $K$ over $k$, then
    \begin{align*}
        \sigma(c^{1/p}) = \zeta^a c^{1/p}
    \end{align*}
    for some $a\in \Z/p\Z$. The group of automorphisms of $K$ over $k$ is cyclic of order $p$.
\end{example}


\clearpage
\section{Galois extensions: main theorem and proof}

\begin{theorem}[Main theorem of Galois theory]
    Let $K/k$ be a Galois extension. Let $G=\mathrm{Aut}(K/k)$ be the group of automorphisms of $K$ over $k$. The poset of subgroups of $G$ is anti-equivalent to the poset of sub-extensions of $K/k$ via
    \begin{align*}
        H\mapsto K^H
    \end{align*}
    and its inverse
    \begin{align*}
        E\mapsto \mathrm{Aut}(K/E).
    \end{align*}
\end{theorem}
\begin{remark}
    Given the anti-equivalence, it must satisfy
    \begin{align*}
        [K:K^H]=|H|,\quad [K:E]=|\mathrm{Aut}(K/E)|.
    \end{align*}
\end{remark}

Before proof, we first verify the anti-equivalence for some examples of Galois extensions.

\begin{example}
    Consider the Galois extension $\mathbb{F}_8/\mathbb{F}_2$. As $[\mathbb{F}_8:\mathbb{F}_2]=3$, there are no intermediate fields. Observe the group of automorphisms:
    \begin{align*}
        \mathrm{Aut(\mathbb{F}_8/\mathbb{F}_2)} & \simeq \Z/3\Z \\
        \sigma_2 &\mapsto 1
    \end{align*}
    where $\sigma_2(a) = a^2$.
\end{example}

\begin{example}
    Consider the Galois extension $\Q(\zeta_5)/\Q$. We can find the anti-equivalence:
    \begin{figure}[H]
        \centering
        \begin{tikzcd}
            \Q(\zeta_5)\arrow[d, dash] & & \{1\}\arrow[d, dash] \\
            \Q(\zeta_5+\zeta_5^{-1})\simeq \Q(\sqrt{5})\arrow[d, dash] & & \{\pm 1\}\arrow[d, dash] \\
            \Q & & \mathrm{Aut}(\Q(\zeta_5)/\Q)\simeq (\Z/5\Z)^\times
        \end{tikzcd}
    \end{figure}
\end{example}

\begin{example}
    Consider the Galois extension $\Q(\sqrt{2},\sqrt{3})/\Q$. We can find the anti-equivalence:
    \begin{figure}[H]
        \centering
        \begin{tikzcd}[column sep=small]
            & \Q(\sqrt{2},\sqrt{3})\arrow[dl, dash]\arrow[d, dash]\arrow[dr, dash] & & & \{1,1\}\arrow[dl, dash]\arrow[d,dash]\arrow[dr, dash] & \\
            \Q(\sqrt{2})\arrow[dr, dash] & \Q(\sqrt{6})\arrow[d,dash] & \Q(\sqrt{3})\arrow[dl, dash] & \langle (-1,1)\rangle \arrow[dr, dash] & \langle (-1,-1)\rangle\arrow[d, dash] & \langle (1,-1)\rangle \arrow[dl, dash] \\
            & \Q & & & \Z/2\Z\times \Z/2\Z &
        \end{tikzcd}
    \end{figure}
\end{example}

\begin{example}
    Consider the Galois extension $\mathbb{F}_{16}/\mathbb{F}_2$. We can find the anti-equivalence:
    \begin{figure}[H]
        \centering
        \begin{tikzcd}
            \mathbb{F}_{16}\arrow[d, dash] & & \{1\}\arrow[d, dash] \\
            \mathbb{F}_4\arrow[d, dash] & & \Z/2\Z\arrow[d, dash] \\
            \mathbb{F}_2 & & \mathrm{Aut}(\mathbb{F}_{16}/\mathbb{F}_2)\simeq \Z/4\Z
        \end{tikzcd}
    \end{figure}
\end{example}

\begin{example}
    Consider the Galois extension $\Q(\zeta_4,3^{1/4})/\Q(\zeta_4)$. We can find the anti-equivalence:
    \begin{figure}[H]
        \centering
        \begin{tikzcd}
            \Q(\zeta_4, 3^{1/4})\arrow[d, dash] & & \{1\}\arrow[d, dash] \\
            \Q(\zeta_4, 3^{1/2})\arrow[d, dash] & & \Z/2\Z\arrow[d, dash] \\
            \Q(\zeta_4) & & \mathrm{Aut}(\Q(\zeta_4,3^{1/4})/\Q(\zeta_4)\simeq \Z/4\Z
        \end{tikzcd}
    \end{figure}
\end{example}

Now we go to the proof of the main theorem. We first prove for the injectivity.

\begin{theorem}
    Let $K$ be a Galois extension of $k$ and let $G=\mathrm{Aut}(K/k)$. Then $k=K^G$.
\end{theorem}
\begin{proof}
    It suffices to show that $k(\alpha)=k$ for all $\alpha\in K^G$. Let $\sigma$ be an embedding
    \begin{align*}
        \sigma:k(\alpha)\rightarrow \overline{k}
    \end{align*}
    over $k$. Then, the image of $\sigma$ is contained in $K$. Since $\alpha\in K^G$, $\sigma(\alpha)=\alpha$. Hence, $\sigma$ is the identity on $k(\alpha)$, and $[k(\alpha):k]_s=1$. As $K/k$ is separable,
    \begin{align*}
        [k(\alpha):k]=[k(\alpha):k]_s=1
    \end{align*}
    and $k=k(\alpha)$.
\end{proof}
\begin{remark}
    The theorem gives the injectivity; for two intermediate fields $E_1$, $E_2$, if
    \begin{align*}
        \mathrm{Aut}(K/E_1)=\mathrm{Aut}(K/E_2),
    \end{align*}
    then
    \begin{align*}
        E_1=K^{\mathrm{Aut}(K/E_1)}=K^{\mathrm{Aut}(K/E_2)}=E_2.
    \end{align*}
\end{remark}

Next, we prove the Artin's theorem to show the one-to-one correspondence.

\begin{lemma}\label{Lemma 3.12.8.}
    Let $E/k$ be a finite algebraic separable extension. If every $\alpha\in E$ has degree at most $n$ over $k$, then $[E:k]\le n$.
\end{lemma}
\begin{proof}
    By primitive element theorem, $E=k(\alpha)$ for some $\alpha\in E$. As $\alpha$ has degree at most $n$ and $[E:k]$ equals the degree of $\alpha$, we obtain the result. 
\end{proof}

\begin{theorem}[Artin]
    Let $K$ be a field and $G$ be a finite group of automorphisms of $K$. Put $k=K^G$. Then, $\mathrm{Aut}(K/k)=G$ and $[K:k]=|G|$.
\end{theorem}
\begin{proof}
    Let $\alpha\in K$. Let
    \begin{align*}
        G= \{\sigma_1,\dots,\sigma_n\}
    \end{align*}
    and consider the polynomial
    \begin{align*}
        f(x) = \prod_{i=1}^n (x-\sigma_i(\alpha)).
    \end{align*}
    As $G$ contains an identity automorphism, $f(x)$ has $\alpha$ as zero. For arbitrary $\tau\in G$, $G=\tau G$, so we obtain $f(x) = \tau(f)(x)$. Hence $f(x)\in k[x]$. This implies that the minimal polynomial of $\alpha$ divides $\deg f(x) = n$. Therefore the degree of $\alpha$ is at most $n$, and Lemma \ref{Lemma 3.12.8.} implies that $[K:k]\le n=|G|$.

    On the other hand, as $K/k$ is simple, $\mathrm{Aut}(K/k)=G$ and
    \begin{align*}
        |G| = |\mathrm{Aut}(K/k)| \le [K:k].
    \end{align*}
    Therefore $[K:k]=|G|$.
\end{proof}

\begin{definition}
    If $K/k$ is Galois, we write
    \begin{align*}
        \mathrm{Gal}(K/k)=\mathrm{Aut}(K/k)
    \end{align*}
    and call it the \emph{Galois group of $K$ over $k$}.
\end{definition}

Now we consider the Galois group of a polynomial.

\begin{definition}
    Let $f(x)\in k[x]$ be irreducible and separable. The \emph{Galois group} $G_f$ of $f(x)$ refers to $\mathrm{Gal}(K/k)$ where $K/k$ is the splitting field of $f(x)$.
\end{definition}
\begin{remark}
    Let
    \begin{align*}
        f(x)=\prod_{i=1}^n (x-\alpha_i),\quad \alpha_i\in \overline{k}.
    \end{align*}
    Then $K=k(\alpha_1,\dots,\alpha_n)$. For all $\tau\in G_f$, $\tau(f(x))=f(x)$ and $\tau(x-\alpha_i)|f(x)$. Hence we obtain that $G_f$ acts on the set $\{\alpha_1,\dots,\alpha_n\}$. Consider the mapping $G\rightarrow S_n$. If $\tau$ is in the kernel of the mapping, then $\tau(\alpha_i)=\alpha_i$ for all $i$, and $\tau$ should be an identity. Hence the kernel is trivial, and the mapping is injective. Also, note that the mapping $G\rightarrow S_n$ is transitive.
\end{remark}

\begin{example}
    In this example, we find the Galois group for polynomials.
    \begin{enumerate}
        \item Consider $f(x) = x^3-x-1$. The Galois group $G$ of $f(x)$ over $\Q$ is isomorphic to a transitive subset of $S_3$. As the discriminant $-4(-1)^3 - 27(-1)^2 = -23$ is not a square, $G\simeq S_3$.
        \item Consider $f(x) = x^3-10$. The zeros of $f(x)$ in $\C=\overline{\Q}$ are $10^{1/3}, 10^{1/3}\zeta_3,10^{1/3}\zeta_3^2$, so the splitting field of $f(x)$ over $\Q$ is $\Q(10^{1/3},\zeta_3)$. As $[\Q(10^{1/3},\zeta_3):\Q]=6$, the order of the Galois group of $f(x)$ over $\Q$ is $6$. The possible choices for automorphism $\sigma$ of $\Q(10^{1/3},\zeta_3)$ over $\Q$ are
        \begin{align*}
            \sigma(10^{1/3}) &\in \{ 10^{1/3},10^{1/3}\zeta_3,10^{1/3}\zeta_3^2 \}, \\
            \sigma(\zeta_3) &\in \{ \zeta_3, \zeta_3^2 \}.
        \end{align*}
        Hence the Galois group is isomorphic to $S_3$.
    \end{enumerate}
\end{example}

\begin{example}
    One possible question is that given a transitive subgroup $G$ of $S_n$, can we classify all $f$ with $G_f=G$. The degree $2$ case is trivial, as the only transitive subgroup of $S_2$ is $S_2$. For $n=3$, there exists a proper transitive subgroup of $S_3$, as $A_3$. Consider an irreducible cubic polynomial
    \begin{align*}
        f(x)=x^3+ax+b\in k[x].
    \end{align*}
    (Here we may cancel out the quadratic term, as we may simply consider $f(x-k)$.) Let
    \begin{align*}
        x^3+ax+b=(x-\alpha)(x-\beta)(x-\gamma)
    \end{align*}
    be the factorization in $\overline{k}$ and let
    \begin{align*}
        d = (\alpha-\beta)(\beta-\gamma)(\gamma-\alpha),\quad \Delta = d^2.
    \end{align*}
    ($\Delta$ is called the \emph{discriminant} of $f(x)$.) With some calculations, we obtain
    \begin{align*}
        \Delta=-4a^3-27b^2.
    \end{align*}
    If $d\in k$, then the Galois automorphisms $\sigma$ should preserve $d$ and
    \begin{align*}
        (\alpha-\beta)(\beta-\gamma)(\gamma-\alpha) = (\sigma(\alpha)-\sigma(\beta))(\sigma(\beta)-\sigma(\gamma))(\sigma(\gamma)-\sigma(\alpha)).
    \end{align*}
    Hence the Galois group of $f(x)$ is isomorphic to $A_3$. Otherwise, the Galois group of $f(x)$ is isomorphic to $S_3$.
\end{example}




\clearpage
\section{Norm and trace}

\begin{definition}
    Let $K/k$ be a finite extension. For $\alpha\in K$, consider
    \begin{align*}
        m_\alpha:K&\rightarrow K \\
        x&\mapsto \alpha x.
    \end{align*}
    It is $k$-linear, so we may put
    \begin{align*}
        N'_{K/k}(\alpha):=\det (m_\alpha).
    \end{align*}
    On the other hand, put
    \begin{align*}
        N_{K/k}(\alpha) := \left( \prod_{\sigma:K\rightarrow\overline{k}}\sigma(\alpha) \right)^{[K:k]_i}.
    \end{align*}
\end{definition}

\begin{proposition}
    For every $\alpha\in K$, $N_{K/k}(\alpha)=N_{K/k}'(\alpha)$.
\end{proposition}
\begin{proof}
    First, we consider the special case $K=k(\alpha)$. The minimal polynomial $f(x)$ for $\alpha$ over $k$ is factored as
    \begin{align*}
        f(x) = \prod_{\sigma:k(\alpha)\rightarrow \overline{k}}(x-\sigma(\alpha))^{[k(\alpha):k]_i}.
    \end{align*}
    Hence the 
\end{proof}

\begin{definition}
    We call $N_{K/k}(\alpha)$ the \emph{norm of $\alpha$ along $K/k$}.
\end{definition}

The trace is the additive counterpart of the norm.

\begin{definition}
    Define
    \begin{align*}
        \mathrm{Tr}'_{K/k}(\alpha) = \mathrm{Tr}(m_\alpha)
    \end{align*}
    and
    \begin{align*}
        \mathrm{Tr}_{K/k}(\alpha) = [K:k]_i\left( \sum_{\sigma:K\rightarrow\overline{k}} \sigma(\alpha) \right).
    \end{align*}
\end{definition}
\begin{remark}
    Note that $[K:k]_i>1$ implies that $\mathrm{Tr}_{K:k}(\alpha)=0$ for all $\alpha\in K$. Hence this characterizes separability.
\end{remark}

\begin{proposition}
    $\mathrm{Tr}_{K/k}'(\alpha) = \mathrm{Tr}_{K/k}(\alpha)$ for all $\alpha\in K$.
\end{proposition}

\begin{proposition}
    Let $k(\alpha)/k$ be a separable extension. Then, 
    \begin{align*}
        \mathrm{Tr}_{k(\alpha)/k}:k(\alpha)\rightarrow k
    \end{align*}
    is a nonzero $k$-linear map.
\end{proposition}
\begin{proof}
    Let $R=k[x]/f(x)$ where $f(x)$ is the minimal polynomial of $\alpha$. Let $\overline{R}=\overline{k}[x]/f(x)$. 
\end{proof}

\begin{proposition}
    Let $K/k$ be a finite separable extension. Then the $k$-bilinear map
    \begin{align*}
        K\times K &\rightarrow k \\
        (x,y) &\mapsto \mathrm{Tr}_{K/k}(xy)
    \end{align*}
    is perfect.
\end{proposition}
\begin{proof}
    The map
    \begin{align*}
        x\mapsto \left( y\mapsto \mathrm{Tr}_{K/k}(xy) \right)
    \end{align*}
    is injective since $\mathrm{Tr}_{K/k}$ is non-zero and $xK=K$ for all $x\in k^\times$. Injective linear maps from $K$ to $K^*$ are isomorphisms.
\end{proof}

Below are some examples of the computations of norm and trace.

\begin{example}
    Consider a finite extension $\Q(\sqrt{3})/\Q$ and let $\alpha=1+\sqrt{3}$. The norm can be calculated as
    \begin{align*}
        N_{\Q(\sqrt{3})/\Q}(1+\sqrt{3}) = (1+\sqrt{3})(1-\sqrt{3})=-2
    \end{align*}
    and the trace can be calculated as
    \begin{align*}
        \mathrm{Tr}_{\Q(\sqrt{3})/\Q}(1+\sqrt{3})=(1+\sqrt{3})+(1-\sqrt{3})=2.
    \end{align*}
\end{example}

\begin{example}
    Consider a finite extension $\Q(2^{1/3})/\Q$ and let $\alpha=2^{1/3}$. The norm can be calculated as
    \begin{align*}
        N_{\Q(2^{1/3})/\Q}(2^{1/3}) = \det \begin{pmatrix}
            0 & 0 & 2 \\
            1 & 0 & 0 \\
            0 & 1 & 0
        \end{pmatrix} = 2
    \end{align*}
    and the trace can be calculated as
    \begin{align*}
        \mathrm{Tr}_{\Q(2^{1/3})/\Q}(2^{1/3}) = \mathrm{Tr} \begin{pmatrix}
            0 & 0 & 2 \\
            1 & 0 & 0 \\
            0 & 1 & 0
        \end{pmatrix} = 0.
    \end{align*}
\end{example}

\begin{example}
    Consider a finite extension $\Q(2^{1/3},\zeta_3)/\Q$ and let $\alpha=2^{1/3}$. The norm can be calculated as
    \begin{align*}
        N_{\Q(2^{1/3},\zeta_3)/\Q}(2^{1/3}) = (2^{1/3})^2(2^{1/3}\zeta_3)^2(2^{1/3}\zeta_3^2)^2=4
    \end{align*}
    and the trace can be calculated as
    \begin{align*}
        \mathrm{Tr}_{\Q(2^{1/3},\zeta_3)/\Q}(2^{1/3}) = 2(2^{1/3})+2(2^{1/3}\zeta_3)+2(2^{1/3}\zeta_3^2)=0.
    \end{align*}
\end{example}

\begin{example}
    Consider a finite extension $\Q(2^{1/3},\zeta_3)/\Q(\zeta_3)$ and let $\alpha=2^{1/3}$. The norm can be calculated as
    \begin{align*}
        N_{\Q(2^{1/3},\zeta_3)/\Q(\zeta_3)}(2^{1/3}) = (2^{1/3})(2^{1/3}\zeta_3)(2^{1/3}\zeta_3^2)=2
    \end{align*}
    and the trace can be calculated as
    \begin{align*}
        \mathrm{Tr}_{\Q(2^{1/3},\zeta_3)/\Q}(2^{1/3}) = (2^{1/3})+(2^{1/3}\zeta_3)+(2^{1/3}\zeta_3^2)=0.
    \end{align*}
\end{example}

\begin{example}
    Consider a finite extension $\Q(\zeta_5)/\Q$ and let $\alpha=\zeta_5$. The norm can be calculated as
    \begin{align*}
        N_{\Q(\zeta_5)/\Q}(\zeta_5) =  \zeta_5\cdot\zeta_5^2\cdot\zeta_5^3\cdot\zeta_5^4 = 1
    \end{align*}
    and the trace can be calculated as
    \begin{align*}
        \mathrm{Tr}_{\Q(\zeta_5)/\Q}(\zeta_5) = \zeta_5+\zeta_5^2+\zeta_5^3+\zeta_5^4 = -1.
    \end{align*}
\end{example}

\begin{example}
    Consider a finite extension $\mathbb{F}_2(\alpha)/\mathbb{F}_2$ where $\alpha^2+\alpha=1$. As the minimal polynomial of $\alpha$ is $x^2+x-1$, the norm of $\alpha$ is $-1$ and the trace of $\alpha$ is $-1$.
\end{example}

\begin{example}
    Consider a finite extension $\mathbb{F}_5(\alpha)/\mathbb{F}_5$ where $\alpha^3+\alpha+1=0$. Consider $\alpha^2$. The minimal polynomial of $\alpha$ is $x^3+x+1$, and let $a,b,c$ be the roots. Then
    \begin{align*}
        N_{\mathbb{F}_5(\alpha)/\mathbb{F}_5}(\alpha^2) = a^2b^2c^2 = (abc)^2 = 1
    \end{align*}
    and
    \begin{align*}
        \mathrm{Tr}_{\mathbb{F}_5(\alpha)/\mathbb{F}_5}(\alpha^2) = a^2+b^2+c^2 = (a+b+c)^2 - 2(ab+bc+ca) = 0-2\cdot 1 = -2 = 3.
    \end{align*}
\end{example}

\begin{example}
    Consider a finite extension $\mathbb{F}_5(\alpha)/\mathbb{F}_5$ where $\alpha^3+\alpha+1=0$. Consider $\alpha+3$. The minimal polynomial of $\alpha$ is $x^3+x+1$, and let $a,b,c$ be the roots. Then
    \begin{align*}
        N_{\mathbb{F}_5(\alpha)/\mathbb{F}_5}(\alpha+3) &= (a+3)(b+3)(c+3)\\ &= (abc)+3(ab+bc+ca)+9(a+b+c)+27\\ &= 0+3\cdot 1 + 9\cdot (-1) + 27 = 21 = 4
    \end{align*}
    and
    \begin{align*}
        \mathrm{Tr}_{\mathbb{F}_5(\alpha)/\mathbb{F}_5}(\alpha^2) = (a+3)+(b+3)+(c+3) = 9 = 4.
    \end{align*}
\end{example}


\clearpage
\section{Applications of Galois theory: solvability}

\begin{definition}
    Let $K$ be a field and $L/K$ be a finite Galois extension. We say $L/K$ is \emph{generated by radicals} if there is a tower of extensions
    \begin{align*}
        K=K_0\subset K_1\subset \cdots \subset K_n=L
    \end{align*}
    such that each $K_i$ is of the form $K_i=K_{i-1}(\alpha_i^{1/d_i})$ for some $\alpha_i\in K_{i-1}$ and $d_i\in\Z$. We always assume that $d_i$ is invertible in $K$.
\end{definition}
\begin{remark}
    We do not lose generality by assuming that all $d_i$'s are equal.
\end{remark}

\begin{proposition}\label{Proposiiton 3.14.2.}
    Let $d\in\Z$ and assume that $K$ contains a primitive $d$-th root of unity. Suppose that $L/K$ is of the form
    \begin{align*}
        K=K_0\subset K_1\subset \cdots\subset K_n=L
    \end{align*}
    with $K_i=K_{i-1}(\alpha_i^{1/d})$ for $\alpha_i\in K_{i-1}$. Then, the Galois closure of $L/K$ has a solvable Galois group.
\end{proposition}
\begin{proof}
    We use induction on $n$. If $n=1$, there exists a mapping
    \begin{align*}
        \mathrm{Gal}(k(\alpha^{1/d})/k) &\rightarrow \Z/d\Z \\
        \sigma &\mapsto m_\sigma
    \end{align*}
    where $\sigma(\alpha^{1/d})/\alpha^{1/d}=\zeta^{m_\sigma}$. Hence the Galois group of $L/K$ is abelian and the assertion holds true. Now, suppose that the assertion holds true for the extension $K_{n-1}/K$. Let $F'$ be the Galois closure of $K_{n-1}$ over $K$. Then, by the induction hypothesis, $F'/K$ is solvable. Let $F$ be the Galois closure of $L$. Then, we have a tower
    \begin{align*}
        K\subset F'\subset F
    \end{align*}
    so it suffices to show that $F/F'$ is abelian. Indeed, $F$ is contained in the Galois closure of $F'(\alpha_n^{1/d})$ over $F'$, which must be abelian by the argument for the case $n=1$.
\end{proof}

\begin{proposition}
    Let $d\in\Z$. Suppose that $L/K$ is of the form
    \begin{align*}
        K=K_0\subset K_1\subset \cdots\subset K_n=L
    \end{align*}
    with $K_i=K_{i-1}(\alpha_i^{1/d})$ for $\alpha_i\in K_{i-1}$. Then, the Galois closure of $L/K$ has a solvable Galois group.
\end{proposition}
\begin{proof}
    Let $\zeta_d$ be a primitive $d$-th root of unity in an algebraic closure $\overline{K}$. Let $F$ be the Galois closure of $L/K$ in $\overline{K}$. Consider the tower of extensions
    \begin{align*}
        K \subset K(\zeta_d) \subset F \subset F(\zeta_d).
    \end{align*}
    $F(\zeta_d)/K(\zeta_d)$ is solvable by \ref{Proposiiton 3.14.2.}, so $F/K(\zeta_d)$ is also solvable. Also, $K(\zeta_d)/K$ is abelian, so it is solvable. Therefore we conclude that $F/K$ is solvable.
\end{proof}

\begin{proposition}
    Suppose that the splitting field of $f(x)\in K[x]$ has a non-solvable Galois group. Then, the roots of $f(x)$ is not contained in any extension generated by radicals.
\end{proposition}
\begin{proof}
    Suppose that the roots are in $L$ where $L/K$ is an extension generated by radicals. Then, the Galois closure $F/K$ has a solvable Galois group, say $G$. Since $F$ contains the splitting field of $f(x)$, it must have a non-solvable quotient, contradicting the solvability of $G$.
\end{proof}

Now we prove the insolvability of the quintic polynomials.

\begin{lemma}
    For a prime $p$, a transitive subgroup of $S_p$ has an element of order $p$.
\end{lemma}
\begin{proof}
    If $G\subset S_p$ is transitive, then
    \begin{align*}
        |G|=p\times |G_1|
    \end{align*}
    where $G_1$ is the stabilizer of $1$. Thus $p$ divides $|G|$. Since $p$ is a prime, it implies that $G$ has an element of order $p$.
\end{proof}

\begin{proposition}
    Suppose that a polynomial $f(x)\in\Q[x]$ is irreducible of prime degree $p$ and has exactly two non-real roots in $\C$. Then, $f(x)$ is not solvable by radicals whenever $p\ge 5$.
\end{proposition}
\begin{proof}
    A subgroup of $S_p$ containing a transposition and an element of order $p$ must be $S_p$. Note that the complex conjugation represents a transposition. On the other hand, an element of order $p$ exists by the previous lemma. So the Galois group of $f(x)$ is $S_p$. As $S_n$ is not solvable for $n\ge 5$, we obtain the result.
\end{proof}

\begin{example}
    The polynomial $2x^5-5x^4+5$ is an example of a polynomial not solvable by radicals.
\end{example}


\clearpage
\section{Applications of Galois theory: constructibility}

\begin{definition}
    Let $C\subset\R^2$ be a subset. A point $x\in \R^2$ is called \emph{constructible from $C$} if $x$ satisfies
    \begin{align*}
        c\in L_1\cap L_2
    \end{align*}
    where $L_1, L_2$ are either a line or a circle containing at least two points in $C$.
\end{definition}

\begin{definition}
    For $n\ge 0$, define $\mathcal{C}_n$ recursively by putting $\mathcal{C}_0=\Q^2$ and taking $\mathcal{C}_n$ to be the set of all points constructible from $\mathcal{C}_{n-1}$. We say a real number $\alpha$ is constructible if $(0,\alpha)\in \mathcal{C}_n$ for some $n$.
\end{definition}

\begin{example}
    Let $\alpha,\beta$ be constructible. We can construct $\alpha+\beta$ and $\alpha\beta$ as below.
    \begin{figure}[H]
        \centering
        \begin{tikzpicture}[>=latex]
            \def\alphaY{3} % Position of alpha on y-axis
            \def\betaY{5}  % Position of beta on y-axis
            \def\alphaPrimeX{3} % Position of alpha' on x-axis
            \def\betaX{5}       % Position of beta on x-axis
            \def\alphaPlusBetaX{8} % Position of alpha+beta on x-axis
        
            % Draw the x and y axes
            \draw [->] (0,0) -- (\alphaPlusBetaX + 1, 0) node [right] {}; % x-axis
            \draw [->] (0,0) -- (0, \betaY + 1) node [above] {};
        
            \draw (0, \alphaY) -- (-0.1, \alphaY) node [left] {$\alpha$};
            \draw (0, \betaY) -- (-0.1, \betaY) node [left] {$\beta$};
        
            % Mark labels on x-axis
            \draw (\alphaPrimeX, 0) -- (\alphaPrimeX, -0.1) node [below] {$\alpha$};
            \draw (\betaX, 0) -- (\betaX, -0.1) node [below] {$\beta$};
            \draw (\alphaPlusBetaX, 0) -- (\alphaPlusBetaX, -0.1) node [below] {$\alpha+\beta$};
            
            \draw (0,3) arc (90:0:\alphaY);
            \draw (0,5) arc (90:0:\betaY);
            \draw [dashed] (3,0) -- (3,5);
            \draw [dashed] (0,5) -- (3,5);
            \draw [dashed] (5,0) -- (5,3);
            \draw [dashed] (0,3) -- (5,3);
            \draw (2,6) -- (\alphaPlusBetaX, 0);
        \end{tikzpicture}
    \end{figure}

    \begin{figure}[H]
        \centering
        \begin{tikzpicture}[scale=0.5, >=latex]
            % Draw the x and y axes
            \draw [->] (-5,0) -- (6, 0) node [right] {};
            \draw [->] (0,-2) -- (0, 13) node [above] {};
        
            \draw (0, -1) -- (-0.1, -1) node [left] {$-1$};
        
            \draw (-3, 0) -- (-3, -0.1) node [below] {$-\alpha$};
            \draw (4, 0) -- (4, -0.1) node [below] {$\beta$};

            \draw (0, 12) -- (-0.1, 12) node [left] {$\alpha\beta$};

            \def\circleCenterX{0.5}
            \def\circleCenterY{5.5}
            \def\circleRadius{6.519202}
            \draw (\circleCenterX, \circleCenterY) circle (\circleRadius);
            
        \end{tikzpicture}
    \end{figure}
    \noindent
    Also, if $\alpha\neq 0$, then $1/\alpha$ is constructible.
    \begin{figure}[H]
        \centering
        \begin{tikzpicture}[scale = 1.6, >=latex]
            % Draw the x and y axes
            \draw [->] (-3,0) -- (1, 0) node [right] {};
            \draw [->] (0,-2) -- (0, 2) node [above] {};
        
            \draw (0, -1) -- (0.1, -1) node [right] {$-1$};
            \draw (0, 1) -- (0.1, 1) node [right] {$1$};
        
            \draw (-2, 0) -- (-2, -0.1) node [below] {$-\alpha$};
            \draw (0.5, 0) -- (0.5, -0.1) node [below] {$1/\alpha$};

            \draw (-0.75, 0) circle (1.25);            
        \end{tikzpicture}
    \end{figure}
    \noindent
    Therefore we observe that the set of constructible numbers form a field.
\end{example}

\begin{proposition}\label{Proposition 3.15.4.}
    If $\alpha$ is a constructible positive real number, so is $\alpha^{1/2}$. In particular, for every positive integers $d$ and $k$, $d^{1/2^k}$ is constructible.
\end{proposition}
\begin{proof}
    If $\alpha$ is constructible, we can construct $\alpha^{1/2}$ as below:
    \begin{figure}[H]
        \centering
        \begin{tikzpicture}[scale = 1.6, >=latex]
            % Draw the x and y axes
            \draw [->] (-2,0) -- (2, 0) node [right] {};
            \draw [->] (0,-2) -- (0, 3) node [above] {};
        
            \draw (0, -1) -- (0.1, -1) node [right] {$-1$};
            \draw (0, 2) -- (0.1, 2) node [right] {$\alpha$};
        
            \draw (-1.414, 0) -- (-1.414, -0.1) node [below] {$-\alpha^{1/2}$};
            \draw (1.414, 0) -- (1.414, -0.1) node [below] {$\alpha^{1/2}$};

            \draw (0, 0.5) circle (1.5);            
        \end{tikzpicture}
    \end{figure}
    \noindent
    $d^{1/2^k}$ can be constructed inductively.
\end{proof}

\begin{proposition}
    An algebraic number $\alpha\in\R$ is constructible if and only if $F=\Q(\alpha)$ admits a tower of subfields
    \begin{align*}
        K_0=\Q\subset K_1\subset K_2\subset\cdots\subset K_n=F
    \end{align*}
    such that $[K_i:K_{i-1}]=2$ for all $i=1,\dots,n$.
\end{proposition}
\begin{proof}
    By Proposition \ref{Proposition 3.15.4.}, we know that $\alpha\in \R$ is constructible if $F=\Q(\alpha)$ admits such a tower of subfields. To prove the converse, suppose that $x=(x_1,x_2)\in\R^2$ is constructible from $\mathcal{C}_t\subset \R^2$. Then, $x_i$ satisfies a quadratic equation $x_i^2+a_ix_i+b_i$ with $(0,a_i),(0,b_i)\in\mathcal{C}_{i-1}$.
\end{proof}

\begin{proposition}
    Suppose that $p=2^n+1$ is a prime. Then, a regular $p$-gon can be constructed.
\end{proposition}
\begin{proof}
    Let $\alpha_p=\cos\left(\frac{2\pi i}{p}\right)$. Then $\Q(\alpha_p)/\Q$ is a cyclic extension of degree $\frac{p-1}{2} = 2^{n-1}$. Let $H=\mathrm{Gal}(\Q(\alpha_p)/\Q)$. The Galois theorem gives us a tower of subfields corresponding to
    \begin{align*}
        H\supset 2H\supset 4H\supset 8H\supset\cdots.
    \end{align*}
\end{proof}

\begin{example}
    By the proposition above, we observe that a regular $17$-gon is constructible. However, a regular $13$-gon is not constructible, as the irreducible polynomial of $\alpha_{13} = \frac{2\pi i}{13}$ is of degree $6$, which is not of power of $2$. We can also observe that the proposition above only holds when $p$ is prime. For example, a regular $9$-gon is not constructible, as the irreducible polynomial of $\alpha_9 = \cos\left(\frac{2\pi i}{9}\right)$ is of degree $3$.
\end{example}

The problem of doubling a cube refers to showing the constructibility of $2^{1/3}$.

\begin{proposition}
    $2^{1/3}$ is not constructible.
\end{proposition}
\begin{proof}
    The degree $[\Q(2^{1/3}):\Q]=3$ is odd.
\end{proof}


\end{document}